\chapter{Introduction without a supplied witness}
\label{prov:ch:wit}
\label{wit:ch}

The rule \rulename{DI} of the native scalar core takes a term: its witness.
This chapter restricts that term.  The witness is first required to be a
variable or the literal $0$, then to be the newest numerical parameter of the
current scope, with $0$ in an empty scope, so that the rule takes no term at
all.  Every restricted calculus of the chapter derives the same sequents as
its unrestricted counterpart, on all open sequents and with the identity
translation of terms and formulas.  In $\Nat$ the restriction of the
witness cannot be combined with a restriction of the targets of
\rulename{QE} (\cref{wit:thm:finite-range}); with a discharge that
substitutes a term read off its major premise, both restrictions hold
together (\cref{wit:sec:spec}).

The chapter is organized by idea.  \Cref{wit:sec:restricted} states the
restricted calculi and what they lose: a calculus with variable witnesses is
closed under renaming and not under substitution, and a calculus with a
scope-determined witness is not closed under the extension of its scope.
\Cref{wit:sec:obstructions} contains the negative results.
\Cref{wit:sec:names} reduces the completeness of a restricted calculus to the
derivability of a name for every term, and \cref{wit:sec:packet} derives the
names from a formula proved by numerical cases.  \Cref{wit:sec:newest} treats
the newest parameter, where derivations must be regenerated instead of
weakened; in the Alias/Back calculus (\cref{wit:sec:aliasback}) a fresh
alias takes the place of the names, and when \rulename{Alias} as well as
\rulename{Back} exports equations only, one naming scheme is needed and
suffices (\cref{wit:sec:ns}).  \Cref{wit:sec:guarded} treats the raw guarded clause calculus, with
and without a naming exception.  \Cref{wit:sec:recursor,wit:sec:copy} treat
two calculi on a recursor whose predecessor computes; the second combines the
restricted introduction with an inferred discharge by means of conditional
copy.  \Cref{wit:sec:spec} restricts the use of a clause to the newest
parameter as well: first with an arbitrary witness of the introduction, and
then, by a vacuous case of the selector, together with the introduction at
the newest parameter, so that no rule takes a term.
\Cref{wit:sec:recon} keeps the witness and removes the matrix instead: the
matrix of the introduction is reconstructed from the proved conclusion of its
premise.

Throughout the chapter $1$ abbreviates $\Suc 0$ and $2$ abbreviates
$\Suc\Suc 0$.

% ---------------------------------------------------------------------
\section{Restricted witnesses}
\label{wit:sec:restricted}

This section draws on \pkg{no\_witness}, \pkg{newest\_parameter},
\pkg{parameter\_witness} and \pkg{deterministic\_introduction}.

\subsection{Scopes and the restricted rules}

In $\Nat$ as stated in \cref{base:sec:kernel} the variables available to a
derivation are left implicit.  A restriction that mentions the variables in
scope needs them explicit.

\begin{definition}[Scopes]
\label{wit:def:scope}
A \emph{numerical scope} is a finite list $\Pi=(a_1,\dots,a_m)$ of distinct
variables, called \emph{parameters}.  A term or formula is \emph{over $\Pi$}
if its free variables are in $\Pi$.  The \emph{newest parameter} is
\[
 \nu(\Pi)\defeq
 \begin{cases}
  a_m, & m>0,\\
  0, & m=0.
 \end{cases}
\]
A judgment is $\Pi;\Gamma\proves A$, where $\Gamma$ and $A$ are over $\Pi$.
The scope and the context are separate lists: an assumption added to $\Gamma$
does not change $\Pi$.  The order of $\Pi$ is proof syntax and has no
numerical meaning.
\end{definition}

With scopes the rules of $\Nat$ read as follows.  \rulename{Hyp},
\rulename{I0}, \rulename{IS}, \rulename{Tr} and \rulename{QE} keep the scope,
and the terms they mention are over $\Pi$.  The three rules with an
eigenvariable append it to the scope of their premise:
\begin{gather*}
 \infrule{\Pi;\Gamma\proves\D{x}{y}B\qquad
          \Pi,b;\ \Gamma,\Q{y}B(b,y)\proves A}
         {\Pi;\Gamma\proves A}\ \rulename{DE}
 \\[1.2ex]
 \infrule{\Pi,z;\ \Gamma,B(t,0)\proves B(t,\Suc z)}
         {\Pi;\Gamma\proves\D{x}{y}B}\ \rulename{DI}
 \qquad
 \infrule{\Pi,z;\ \Gamma,E(z)\proves E(\Suc z)}
         {\Pi;\Gamma\proves\Q{y}E(y)}\ \rulename{QI}
\end{gather*}
The conclusions are over $\Pi$, the witness $t$ of \rulename{DI} is over
$\Pi$, and $E$ is a literal equation.  A derivation of $\Nat$ in the sense of
\cref{base:sec:kernel} is a derivation in this sense over any scope that
contains its free variables, and conversely.

\begin{definition}[Restrictions of the witness]
\label{wit:def:calculi}
Each of the following calculi has the terms, formulas and rules of $\Nat$,
except that the witness $t$ of \rulename{DI} is restricted.
\begin{enumerate}[label=\textup{(\roman*)}]
\item $\Nat_{0}$: $t$ is the literal $0$.
\item $\Nat_{\mathrm v}$: $t$ is a parameter of $\Pi$.
\item $\Nat_{0\mathrm v}$: $t$ is the literal $0$ or a parameter of $\Pi$.
\item $\Nat_{\nu}$: $t$ is $\nu(\Pi)$.  The rule takes no term:
\[
 \infrule{\Pi,z;\ \Gamma,B(\nu(\Pi),0)\proves B(\nu(\Pi),\Suc z)}
         {\Pi;\Gamma\proves\D{x}{y}B}\ \rulename{DI}_{\nu}
\]
\end{enumerate}
In \rulename{DI}$_\nu$ the witness is computed in the scope $\Pi$ of the
conclusion, before the eigenvariable $z$ is appended.  The eigenvariable is
never the witness.
\end{definition}

Every derivation of $\Nat_{\nu}$ is a derivation of $\Nat_{0\mathrm v}$, and
every derivation of $\Nat_{0}$, $\Nat_{\mathrm v}$ or $\Nat_{0\mathrm v}$ is
a derivation of $\Nat$.  Hence on derivable judgments
\begin{equation}
 \Nat_{\nu}\ \subseteq\ \Nat_{0\mathrm v}\ \subseteq\ \Nat .
 \label{wit:eq:sandwich}
\end{equation}
The distinction between a variable and the constant $0$ is syntactic: a
term that is provably equal to $0$ is not the literal $0$.

The derivations of reflexivity, symmetry, transitivity and congruence of
\cref{base:sec:structural} use \rulename{I0} and \rulename{Tr} only, and
equation induction is \rulename{QI} followed by \rulename{QE}.  These rules
are therefore available, with the same derivations, in every calculus of
\cref{wit:def:calculi}; equations are transformed by them without further
comment.  In particular the selector equations $\sel(0,a,b)=a$ and
$\sel(\Suc c,a,b)=b$ are available.

\subsection{What a restricted witness loses}

\begin{lemma}[Structural rules for variable witnesses]
\label{wit:lem:structural-0v}
In $\Nat_{0\mathrm v}$ and in $\Nat_{\mathrm v}$ the following are
admissible.
\begin{enumerate}[label=\textup{(\roman*)}]
\item Renaming and weakening: if $\Pi;\Gamma\proves A$, $\rho$ is an
 injective map from $\Pi$ into a scope $\Pi'$, and $\Delta$ is a context
 over $\Pi'$ that contains $\rho\Gamma$, then $\Pi';\Delta\proves\rho A$.
\item Cut: if $\Pi;\Gamma\proves A$ and $\Pi;\Gamma,A\proves B$, then
 $\Pi;\Gamma\proves B$.
\end{enumerate}
\end{lemma}

\begin{proof}
(i) By induction on the derivation, for all $\rho$ and $\Delta$.  A renaming
sends a variable witness to a variable witness and fixes the literal $0$.
At a rule with an eigenvariable, choose a fresh parameter of the receiving
scope, extend $\rho$ by it and extend $\Delta$ by the renamed assumption.
(ii) By induction on the derivation of $\Pi;\Gamma,A\proves B$.  A use of
\rulename{Hyp} for $A$ is replaced by the derivation of $A$, renamed along
the inclusion of $\Pi$ into the scope of that leaf and weakened by the
assumptions added on the way, by (i).  All other rule data are unchanged.
\end{proof}

The proof of (i) does not extend to substitution.  A derivation of
$\Nat_{0\mathrm v}$ that introduces $\D{x}{y}B$ at a parameter $a$ is not
turned into a derivation by substituting a compound term for $a$, because
the image of the witness is not a variable.  Closure under substitution
holds for $\Nat_{0\mathrm v}$ (\cref{wit:thm:zero-or-variable}), but only as
a consequence of its completeness.  The arguments of the chapter therefore
never substitute a term into a derivation of a restricted calculus; where a
formula $B(t)$ is needed for a compound $t$, it is obtained by \rulename{QE}
from a closure, or the derivation is constructed anew for that instance.

In $\Nat_{\nu}$ less is available.

\begin{lemma}[Weakening of the context at a fixed scope]
\label{wit:lem:ctxweak}
If $\Pi;\Gamma\proves A$ in $\Nat_{\nu}$ and $\Gamma\subseteq\Delta$ with
$\Delta$ over $\Pi$, then $\Pi;\Delta\proves A$ in $\Nat_{\nu}$.
\end{lemma}

\begin{proof}
By induction on the derivation.  The scope of every node is unchanged, so
every \rulename{DI}$_\nu$ selects the same witness; the eigenvariables are
not in $\Pi$ and hence not free in $\Delta$.
\end{proof}

\begin{example}[Extending the scope changes the witness]
\label{wit:ex:weakening}
In the empty scope $\Nat_{\nu}$ derives $\D{x}{y}(\sel(y,0,x)=0)$:
the selected witness is $0$, and the premise asks for
$\sel(\Suc z,0,0)=0$, an instance of \rulename{IS}.  The same tree is not a
derivation over the scope $(a)$.  There the rule selects $a$ and its premise
asks for $\sel(\Suc z,0,a)=0$ from $\sel(0,0,a)=0$, that is, for $a=0$,
which is not derivable because it is false at $a=1$ in the standard reading.
The judgment itself remains derivable over $(a)$, by
\cref{wit:thm:newest}; its derivation there is a different tree.
\end{example}

Thus a derivation of $\Nat_{\nu}$ cannot be weakened to a larger scope by
replaying its tree, and the usual proof of cut fails as well: when the
receiving derivation enters a branch with a fresh parameter, the derivation
to be inserted would have to be weakened to the larger scope.
\Cref{wit:sec:newest} replaces the weakening of derivations by families of
derivations that are constructed anew in every scope; in the raw guarded
calculus with a naming exception, weakening is recovered by naming the old
selection anew (\cref{wit:thm:weakening}).

% ---------------------------------------------------------------------
\section{Obstructions}
\label{wit:sec:obstructions}

This section draws on \pkg{newest\_parameter} for the two shapes and on
\pkg{constructed\_computation} for the simultaneous restriction.

\subsection{Each shape alone is incomplete}

\begin{proposition}[One shape is not enough]
\label{wit:prop:shapes}
\begin{enumerate}[label=\textup{(\roman*)}]
\item $\Nat_{0}$ does not derive $\D{x}{y}(\sel(y,1,x)=1)$ in the empty
 scope and context.  $\Nat$ derives it.
\item $\Nat_{\mathrm v}$ derives no formula $\D{x}{y}B$ in the empty scope
 and the empty context.  $\Nat$ derives $\D{x}{y}(0=0)$ there.
\end{enumerate}
\end{proposition}

\begin{proof}
(i) Read the terms in the natural numbers, read $\Q{y}B$ as
$B(0)\to\forall y\,B(y)$, and read $\D{x}{y}B$ at $x=0$, as
$B(0,0)\to\forall y\,B(0,y)$.  A judgment is valid if its conclusion is true
under every assignment of its parameters that makes its context true.  All
rules of $\Nat_{0}$ preserve validity.  For \rulename{Hyp}, the computation
equations, \rulename{Tr} and \rulename{QE} this is the usual verification,
and \rulename{QI} is induction in the natural numbers.  For \rulename{DE},
the premise is valid at $b=0$, where its assumption $\Q{y}B(0,y)$ is the
reading of the major premise.  For \rulename{DI} at the witness $0$, the
premise gives $B(0,\Suc z)$ for every $z$ from $B(0,0)$, and $B(0,0)$ gives
itself.  The formula $\D{x}{y}(\sel(y,1,x)=1)$ is false in this reading: its
instance at $y=0$ is $\sel(0,1,0)=1$, which is true, and its instance at
$y=1$ is $\sel(1,1,0)=1$, which says $0=1$.  In $\Nat$ the rule
\rulename{DI} with the witness $1$ asks for $\sel(\Suc z,1,1)=1$, an
instance of \rulename{IS}.

(ii) Call a formula \emph{simple} if it is an equation or a formula
$\Q{y}E$ with $E$ an equation.  In $\Nat_{\mathrm v}$, every formula derived
over the empty scope from a context of simple formulas is simple.  This is
proved by induction on the derivation, the claim being made only for the
nodes whose scope is empty.  \rulename{Hyp} yields a member of the context.
\rulename{I0}, \rulename{IS} and \rulename{Tr} yield equations.  The major
premise $\Q{y}B$ of \rulename{QE} has the same scope and context and is
simple by the induction hypothesis, so $B$ is an equation and so is $B(t)$.
The major premise of \rulename{DE} would be a formula $\D{x}{y}B$ derived
over the empty scope from the same context, against the induction
hypothesis; hence \rulename{DE} does not occur at such a node.
\rulename{DI} needs a parameter of the empty scope as its witness and does
not occur.  The conclusion $\Q{y}E$ of \rulename{QI} is simple by the form
of the rule.  In $\Nat$ the rule \rulename{DI} with the witness $0$ derives
$\D{x}{y}(0=0)$, its premise being reflexivity.
\end{proof}

The Lean development contains both statements: the calculus with the witness
restricted by shape is \Lean{DWPrf}, and the two parts are
\Lean{dwprf\_zero\_incomplete} and \Lean{dwprf\_var\_incomplete} in the
module \Lean{NativeWitnessShape}.  Part (ii) explains why
\rulename{DI}$_\nu$ uses $0$ in the empty scope, and part (i) why the
parameters matter.  The proposition concerns the two shapes only; it is not a
statement about all proof calculi.

\subsection{The witness and the target cannot both be restricted}

The target $t$ of \rulename{QE} can be restricted to a variable or the
numeral $1$ without loss (Lean: \Lean{vqprf\_iff\_nprf}).  The two
restrictions do not combine.

\begin{theorem}[Finite-range obstruction]
\label{wit:thm:finite-range}
Let $\Nat_{\mathrm R}$ be $\Nat$ with the witness of \rulename{DI}
restricted to $0$ or a parameter in scope and the target of \rulename{QE}
restricted to $1$ or a parameter in scope.  Then $\Nat_{\mathrm R}$ does not
derive
\[
 \D{x}{y}\bigl(\sel(y,1,x)=1\bigr),
\]
which $\Nat$ derives.  The same holds when $\Nat$ is replaced by the
calculus of the reports, with transport in formula templates and with the
witnessed introduction by progress for literal equations, and for the
introduction by progress with a compound matrix.
\end{theorem}

\begin{proof}
Evaluate terms in the natural numbers.  For a formula $A$ and an assignment
$e$ of numbers to variables define a truth value $[A]_e$ by
\begin{align*}
 [s=t]_e &\iffdef \text{$s$ and $t$ have the same value under $e$},\\
 [\Q{y}B]_e &\iffdef [B]_{e[y\mapsto 0]}\Rightarrow[B]_{e[y\mapsto 1]},\\
 [\D{x}{y}B]_e &\iffdef
   [B]_{e[x\mapsto 0,\,y\mapsto 0]}\Rightarrow[B]_{e[x\mapsto 0,\,y\mapsto 1]} .
\end{align*}
These clauses define an invariant of derivations; they are not an ordinary
first-order interpretation.  By induction on $B$, the value
$[B(t)]_e$ equals $[B]_{e[y\mapsto n]}$, where $n$ is the value of $t$ under
$e$; bound variables are renamed away from $t$ before the substitution.
Call a judgment \emph{valid} if its conclusion is true under the assignment
$e_0$ that sends every parameter to $0$ whenever its context is true under
$e_0$.  Every rule of $\Nat_{\mathrm R}$ preserves validity.

\rulename{Hyp} is immediate.  The computation equations are true under every
assignment, and \rulename{Tr}, as well as transport in a formula template,
preserves truth by the substitution property.

\rulename{DI}: an allowed witness has the value $0$ under $e_0$, and the
eigenvariable $z$ has the value $0$, so $\Suc z$ has the value $1$.  The
premise therefore supplies exactly the implication that defines
$[\D{x}{y}B]_{e_0}$.  For the introduction by progress the premise has the
assumption $B(t,z)$, which has the same value at $z=0$, and the same
implication results; the matrix need not be an equation for this argument.
\rulename{QI} is the same computation with the clause for $\Q{}$.  The
context does not contain $z$ and is unaffected.
% ed.: the package states the invariant for the calculus of the reports
% (introduction by cases and by progress); the case of QI, needed for the
% eight rules of N, is added here and is the same verification.

\rulename{DE}: the fresh parameter $b$ has the value $0$.  The assumption
$\Q{y}B(b,y)$ then has the truth value of the major premise
$\D{x}{y}B$, so the validity of the second premise gives the conclusion.

\rulename{QE}: at a parameter, the conclusion $B(t)$ has the truth value of
the premise $B(0)$.  At the target $1$ the rule is the elimination of the
implication that defines $[\Q{y}B]_{e_0}$.

By induction on derivations every derivable judgment is valid.  For the
displayed formula the instance at $x=0$, $y=0$ is $\sel(0,1,0)=1$, which is
true, and the instance at $x=0$, $y=1$ is $\sel(1,1,0)=1$, which is false.
In $\Nat$, \rulename{DI} with the witness $1$ derives the formula from
$\sel(\Suc z,1,1)=1$.
\end{proof}

\begin{theorem}[Finitely many numeral targets do not help]
\label{wit:thm:finite-range-general}
Fix $N\ge 1$ and let $W=\{0,\dots,N\}$.  Let $\Nat_{\mathrm R,N}$ be $\Nat$
with the witness of \rulename{DI} restricted to $0$ or a parameter and the
target of \rulename{QE} restricted to a parameter or a numeral in $W$.
Then $\Nat_{\mathrm R,N}$ does not derive
\[
 \D{x}{y}\bigl(\sel(y,N+1,x)=N+1\bigr),
\]
which $\Nat$ derives with the witness $N+1$.
\end{theorem}

\begin{proof}
Evaluate terms in the natural numbers.  For an arbitrary assignment $e$ of
numbers to variables define $[s=t]_e$ as before and
\begin{align*}
 [\Q{y}B]_e &\iffdef [B]_{e[y\mapsto 0]}\Rightarrow
   \bigwedge_{j\in W}[B]_{e[y\mapsto j]},\\
 [\D{x}{y}B]_e &\iffdef \bigvee_{i\in W}\Bigl([B]_{e[x\mapsto i,\,y\mapsto 0]}
   \Rightarrow\bigwedge_{j\in W}[B]_{e[x\mapsto i,\,y\mapsto j]}\Bigr).
\end{align*}
The bound variables range over $W$; the values of terms, and hence the
values at which a matrix is evaluated after a substitution, may lie outside
$W$, and the substitution property holds as before.  A judgment is
\emph{valid} if its conclusion is true under every assignment that sends
its parameters into $W$ and makes its context true.  \rulename{DI} by cases
preserves validity: the witness has a value
$i\in W$, every positive $j\le N$ is $\Suc(j-1)$ with $j-1\in W$, and the
case $j=0$ is the assumption.  The introductions by progress, and
\rulename{QI}, preserve validity by the finite induction from $0$ to $N$.
\rulename{DE} uses a witness $i\in W$ and the validity of its second premise
at the assignment $b\mapsto i$.  \rulename{QE} has a target with a value in
$W$.  Transport and the computation equations are as before.  The displayed
formula is false: for every $i\in W$ the instance at $y=0$ is true, and the
instance at $y=1$, which belongs to the conjunction because $1\in W$, says
$i=N+1$.
\end{proof}

\begin{remark}[Scope of the obstruction]
\label{wit:rem:obstruction-scope}
The two theorems exclude the equality of the theorem sets in the identity
language under the simultaneous restriction.  They do not show that a
differently interpreted theory with both restrictions fails to be
bi-interpretable with $\iPRA$.  The invariant of
\cref{wit:thm:finite-range} is not the case $N=1$ of
\cref{wit:thm:finite-range-general}: it fixes the witness at $0$ and gives a
smaller separating formula.  Both invariants remain sound for the rule of
\pkg{constructed\_computation} that supplies a variable from a known
instance at zero,
\[
 \infrule{\Pi,w;\ \Gamma,H(w)\proves A\qquad\Pi;\Gamma\proves H(0)}
         {\Pi;\Gamma\proves A}\qquad w\notin\FV(\Gamma,A).
\]
In the first invariant $H(w)$ and $H(0)$ have the same truth value.  In the
second, extend the assignment by $w\mapsto 0$, which is allowed because
$0\in W$; the second premise makes $H(w)$ true there and the first supplies
$A$.  No rule that names an arbitrary term by a witness outside $W$ is
covered.
\end{remark}

% ---------------------------------------------------------------------
\section{Names}
\label{wit:sec:names}

This section draws on \pkg{deterministic\_introduction} for the naming
matrix, on \pkg{no\_witness} and \pkg{parameter\_witness} for the derived
rules of the calculus with variable witnesses, and on the Lean module
\Lean{NativeNamedWitness} for the organization.

A restricted introduction cannot introduce $\D{x}{y}B$ at a compound term
$t$.  It can introduce it at a parameter $b$ that is provably equal to $t$.
A \emph{name} of $t$ is a formula whose opening at a fresh $b$ yields the
equation $b=t$.

\subsection{Naming matrices}

\begin{definition}[Gate, naming matrix, alias matrix]
\label{wit:def:naming}
For an equation $E\equiv(s=t)$ and a term $v$ the \emph{gate} of $E$ is
\[
 \gamma_v(E)\defeq\bigl(\sel(v,0,s)=\sel(v,0,t)\bigr).
\]
For a term $r$, and variables $x,y$ not free in $r$, the \emph{naming
matrix} and the \emph{alias matrix} of $r$ are
\[
 \mathrm{Nm}_r(x,y)\defeq\gamma_y(x=r)
   =\bigl(\sel(y,0,x)=\sel(y,0,r)\bigr),
 \qquad
 \mathrm{Al}_r(x,y)\defeq\bigl(\sel(y,r,x)=r\bigr).
\]
The formulas $\D{x}{y}\mathrm{Nm}_r$ and $\D{x}{y}\mathrm{Al}_r$ are the
\emph{names} of $r$.
\end{definition}

By the selector equations, $\gamma_0(E)$ is derivable in every context, both
sides being equal to $0$, and $\gamma_{\Suc z}(E)$ is interderivable with
$E$.  The ordinary reading of either name is equivalent to
$\exists x\,(x=r)$: the instance of the matrix at $y=0$ is provable and its
instance at $y=1$ is equivalent to $x=r$.  The gate cannot be omitted.  The
formula $\D{x}{v}(x=r)$, whose matrix does not contain $v$, reads
$\exists x\,(x=r\to x=r)$ and says nothing about $r$.

\begin{lemma}[A name is derivable with the witness it names]
\label{wit:lem:name-in-N}
For every term $r$ over $\Pi$, $\Nat$ derives
$\Pi;\Gamma\proves\D{x}{y}\mathrm{Nm}_r$ and
$\Pi;\Gamma\proves\D{x}{y}\mathrm{Al}_r$ by \rulename{DI} with the witness
$r$.
\end{lemma}

\begin{proof}
The premise for the naming matrix asks for
$\sel(\Suc z,0,r)=\sel(\Suc z,0,r)$, which is reflexivity.  The premise for
the alias matrix asks for $\sel(\Suc z,r,r)=r$, an instance of
\rulename{IS}.  Neither uses its assumption.
\end{proof}

\begin{lemma}[Opening a name]
\label{wit:lem:open-name}
Let $\Gamma$ contain $\Q{y}\mathrm{Nm}_r(b,y)$ or
$\Q{y}\mathrm{Al}_r(b,y)$.  Then $\Gamma\proves b=r$, by a derivation that
uses \rulename{Hyp}, \rulename{I0}, \rulename{IS}, \rulename{Tr} and one
\rulename{QE} with the target $1$.
\end{lemma}

\begin{proof}
For the naming matrix, the instance at $0$ is
$\sel(0,0,b)=\sel(0,0,r)$; both sides equal $0$ by \rulename{I0}.
\rulename{QE} at $1$ gives $\sel(1,0,b)=\sel(1,0,r)$, and \rulename{IS}
turns the two sides into $b$ and $r$.  For the alias matrix, the instance at
$0$ is $\sel(0,r,b)=r$, an instance of \rulename{I0}.  \rulename{QE} at $1$
gives $\sel(1,r,b)=r$, and \rulename{IS} gives $\sel(1,r,b)=b$.
\end{proof}

Opening the naming matrix uses reflexivity where opening the alias matrix
uses one successor computation; otherwise the two are interchangeable.  The
naming matrix is the one that \pkg{deterministic\_introduction} recognizes
syntactically (\cref{wit:sec:guarded}); the alias matrix is the one of the
Lean development.

\subsection{Derived rules with variable witnesses}

Let $\Nat_{\mathrm v}^{+}$ be $\Nat_{\mathrm v}$ together with the axiom
scheme
\[
 \Pi;\Gamma\proves\D{x}{y}\mathrm{Al}_a\qquad\text{for every term $a$ over
 $\Pi$.}
\]
A renaming sends an instance of the scheme to an instance, so
\cref{wit:lem:structural-0v} holds for $\Nat_{\mathrm v}^{+}$ with the same
proof.  The next two lemmas are stated for $\Nat_{\mathrm v}^{+}$ and for
$\Nat_{0\mathrm v}$ at once.

\begin{lemma}[Closure by cases]
\label{wit:lem:qc-0v}
If $\Pi,z;\ \Gamma,B(0)\proves B(\Suc z)$, then
$\Pi;\Gamma\proves\Q{y}B(y)$.  Here $B$ is an arbitrary native formula.
\end{lemma}

\begin{proof}
Choose $r$ not free in $B$.  In $\Nat_{0\mathrm v}$, \rulename{DI} with the
witness $0$ derives $\D{r}{y}B(y)$ from the given premise.  \rulename{DE} at
a fresh $b$ adds the assumption $\Q{y}B(y)$, which is the conclusion and
does not contain $b$.  In $\Nat_{\mathrm v}^{+}$ a variable is needed as
the witness.  Open the name $\D{x}{y}\mathrm{Al}_0$ of $0$ at a fresh
parameter $c$ and introduce $\D{r}{y}B(y)$ at $c$; its premise is the given
one, renamed along the inclusion of scopes and weakened by
\cref{wit:lem:structural-0v}.  Eliminate $\D{r}{y}B(y)$ as before, and
discharge $c$ by \rulename{DE}: the conclusion $\Q{y}B(y)$ does not contain
it.
\end{proof}

\begin{lemma}[Leibniz rule]
\label{wit:lem:leibniz-0v}
For every native formula $F(v)$ with a distinguished variable $v$: if
$\Pi;\Gamma\proves s=t$ and $\Pi;\Gamma\proves F(s)$, then
$\Pi;\Gamma\proves F(t)$.
\end{lemma}

\begin{proof}
By induction on the number of binders of $F$.  The statement is taken for
all scopes, contexts and terms at once, because the induction hypothesis is
applied to new instances of the matrix; no term is substituted into a
derivation.

If $F(v)$ is an equation $p(v)=q(v)$, congruence gives $p(s)=p(t)$ and
$q(s)=q(t)$, and symmetry and transitivity give $p(t)=q(t)$.

Let $F(v)=\Q{y}B(v,y)$.  By \cref{wit:lem:qc-0v} it suffices to derive
$B(t,\Suc z)$ over $\Pi,z$ from $\Gamma,B(t,0)$.  The equation $t=s$ and the
induction hypothesis for the matrix $B(v,0)$ give $B(s,0)$.  \rulename{QE}
at $\Suc z$, applied to the weakened premise $\Q{y}B(s,y)$, gives
$B(s,\Suc z)$.  The equation $s=t$ and the induction hypothesis for
$B(v,\Suc z)$ give $B(t,\Suc z)$.

Let $F(v)=\D{x}{y}B(v,x,y)$.  Apply \rulename{DE} to the premise at a fresh
parameter $b$, which adds $\Q{y}B(s,b,y)$.  Introduce
$\D{x}{y}B(t,x,y)$ at the variable $b$.  The premise of this introduction
has the assumption $B(t,b,0)$ and asks for $B(t,b,\Suc z)$: the induction
hypothesis gives $B(s,b,0)$, \rulename{QE} at $\Suc z$ gives
$B(s,b,\Suc z)$, and the induction hypothesis gives $B(t,b,\Suc z)$.  The
conclusion does not contain $b$, and \rulename{DE} discharges it.  Where a
premise is used below a new parameter it is renamed and weakened by
\cref{wit:lem:structural-0v}.
\end{proof}

\subsection{The introduction at a name}

\begin{theorem}[Variable witnesses and a name for every term]
\label{wit:thm:named}
For every scope $\Pi$, context $\Gamma$ and formula $A$ over $\Pi$,
\[
 \Pi;\Gamma\proves_{\Nat_{\mathrm v}^{+}}A
 \quad\Longleftrightarrow\quad
 \Pi;\Gamma\proves_{\Nat}A .
\]
\end{theorem}

\begin{proof}
From left to right, an introduction at a variable is an instance of
\rulename{DI}, and the axioms are derivable in $\Nat$ by
\cref{wit:lem:name-in-N}.

From right to left it suffices to show that \rulename{DI} at an arbitrary
term $t$ over $\Pi$ is admissible in $\Nat_{\mathrm v}^{+}$; the claim then
follows by induction on the derivation in $\Nat$.  Let a derivation of
\[
 \Pi,z;\ \Gamma,B(t,0)\proves B(t,\Suc z)
\]
be given.  Take the axiom $\D{x}{y}\mathrm{Al}_t$ and apply \rulename{DE} to
it at a fresh parameter $b$.  The branch has the scope $\Pi,b$ and the
context $\Gamma_b=\Gamma,\Q{y}\mathrm{Al}_t(b,y)$, and
\cref{wit:lem:open-name} gives $\Gamma_b\proves b=t$.  In this branch
introduce $\D{x}{y}B$ at the variable $b$.  The premise of the introduction
has the scope $\Pi,b,z$ and the context $\Gamma_b,B(b,0)$.  The Leibniz rule
gives $B(t,0)$ from $B(b,0)$ and $b=t$.  Rename the given derivation along
the inclusion of $\Pi,z$ into $\Pi,b,z$, weaken it to the context
$\Gamma_b,B(b,0),B(t,0)$, and cut $B(t,0)$; the result is $B(t,\Suc z)$.
The Leibniz rule with $t=b$ gives $B(b,\Suc z)$.  The introduction at $b$
yields $\Pi,b;\Gamma_b\proves\D{x}{y}B$, a formula over $\Pi$, and
\rulename{DE} discharges $b$.
\end{proof}

\begin{corollary}[Completeness is the derivability of the names]
\label{wit:cor:names-iff}
$\Nat_{0\mathrm v}$ derives the same judgments as $\Nat$ if and only if it
derives $\Pi;\varnothing\proves\D{x}{y}\mathrm{Al}_a$ for every scope $\Pi$
and every term $a$ over $\Pi$.
\end{corollary}

\begin{proof}
If the two calculi agree, the names are derivable by
\cref{wit:lem:name-in-N}.  Conversely, suppose the names are derivable in
$\Nat_{0\mathrm v}$.  A derivation of $\Nat_{\mathrm v}^{+}$ becomes a
derivation of $\Nat_{0\mathrm v}$ when each axiom leaf is replaced by the
derivation of that name, weakened to the context of the leaf by
\cref{wit:lem:structural-0v}.  With \cref{wit:thm:named} and
\eqref{wit:eq:sandwich} this gives
$\Nat\subseteq\Nat_{\mathrm v}^{+}\subseteq\Nat_{0\mathrm v}\subseteq\Nat$.
\end{proof}

\begin{remark}[Lean status]
\label{wit:rem:lean-names}
\Cref{wit:thm:named} is \Lean{dnprf\_iff\_nprf} for the calculus
\Lean{DNPrf}, and \cref{wit:cor:names-iff} is
\Lean{dwprf\_complete\_iff\_names}; both are in \Lean{NativeNamedWitness}.
The steps of the Lean proof are the ones above: the name, its opening, the
introduction at the name, and a Leibniz rule proved over one context without
substitution into derivations; a dummy witness in a scope without variables
is a name of $0$, as in \cref{wit:lem:qc-0v}.  The derivation of the names
inside the restricted calculus, which is the subject of the next section, is
formalized as well (\cref{wit:rem:lean-complete}).
\end{remark}

% ---------------------------------------------------------------------
\section{The coverage packet}
\label{wit:sec:packet}

This section draws on \pkg{no\_witness} (its restriction note),
\pkg{parameter\_witness} (its appendix on restricted witnesses) and
\pkg{newest\_parameter}.

By \cref{wit:cor:names-iff} it remains to derive the name of an arbitrary
term $t$ with variable witnesses.  The name cannot be introduced at $t$.
The route is indirect: a formula $X(y)$ that says ``$y$ is $0$ or a
successor'' is proved for all $y$ by numerical cases, where in the successor
case the witness is the case variable; it is then instantiated at
$y=\Suc t$ by \rulename{QE}, whose target is unrestricted; and the
injectivity of the successor turns the opened instance into a name of $t$.

\subsection{The predecessor}

Let $\mathcal E$ be the equational fragment of \cref{base:def:E}:
\rulename{Hyp}, the two computation equations, \rulename{Tr} and induction
for an equation, on judgments whose assumptions and conclusions are
equations.  By \cref{base:thm:pred}, whose equational derivation replaces
every numerical case distinction by the interpolation and the two-point
rule of \cref{base:lem:interp}, $\mathcal E$ derives, for the predecessor
term written $P$ there and $\operatorname{pred}$ in this chapter,
\begin{equation}
 \operatorname{pred}(0)=0,\qquad\operatorname{pred}(\Suc n)=n .
 \label{wit:eq:pred}
\end{equation}
$\mathcal E$ is closed under substitution (\cref{base:lem:E}), so for every
term $r$ it derives $\operatorname{pred}(\Suc r)=r$.  In a calculus that has
\rulename{QI} and \rulename{QE}, a derivation of $\mathcal E$ is a
derivation without any introduction of $\D{}{}$.  Hence in $\Nat_{0}$,
$\Nat_{0\mathrm v}$ and $\Nat_{\nu}$, for all $\Pi$, $\Gamma$ and terms
$r,s$ over $\Pi$,
\begin{equation}
 \Pi;\Gamma\proves\operatorname{pred}(\Suc r)=r,
 \qquad
 \Pi;\ \Gamma,\ \Suc r=\Suc s\proves r=s ,
 \label{wit:eq:inj}
\end{equation}
the second by congruence with $\operatorname{pred}$.  The order of the
argument matters.  The predecessor is obtained before any term is named;
obtaining it from the bridge of \cref{base:thm:bridge}, which is proved for
the unrestricted calculus, would be circular.  That the derivation needs no
case distinction under a hypothesis, hence no \rulename{DI}, is the
contribution of the four packages recorded in \cref{base:rem:predsources}.
The predecessor is formalized in this fragment: in Lean the fragment is
\Lean{EqPrf} (hypothesis, the two computations, transport in one template,
equation induction), with \Lean{EqPrf.pred\_zero}, \Lean{EqPrf.pred\_succ}
and \Lean{EqPrf.succ\_inj}, and \Lean{EqPrf.toDW} embeds it in the kernel
under any restriction of the witness.

\subsection{The packet and the names}

\begin{definition}[Coverage packet]
\label{wit:def:packet}
For a term $y$ put
\[
 M_y(x,v)\defeq\gamma_v\bigl(\sel(y,0,\Suc x)=y\bigr)
  =\bigl(\sel(v,0,\sel(y,0,\Suc x))=\sel(v,0,y)\bigr),
 \qquad
 X(y)\defeq\D{x}{v}M_y(x,v).
\]
\end{definition}

The ordinary reading of $X(y)$ is equivalent to
$\exists x\,(\sel(y,0,\Suc x)=y)$, which holds at $y=0$ for every $x$ and
says $y=\Suc x$ at a successor.  The formula never asks for an introduction
of $\exists x\,(x=t)$ at the witness $t$.

\begin{lemma}[The packet with variable witnesses]
\label{wit:lem:packet-0v}
In $\Nat_{0\mathrm v}$, for every term $t$ over $\Pi$,
$\Pi;\Gamma\proves X(t)$.
\end{lemma}

\begin{proof}
$X(0)$: introduce at the witness $0$.  The premise asks for
$M_0(0,\Suc z)$, that is, for
$\sel(\Suc z,0,\sel(0,0,1))=\sel(\Suc z,0,0)$.  By \rulename{IS} and
\rulename{I0} the left side equals $\sel(0,0,1)$ and then $0$, and the right
side equals $0$.

$X(\Suc z)$ over the scope $\Pi,z$: introduce at the variable $z$.  The
premise asks for $M_{\Suc z}(z,\Suc z')$, that is, for
$\sel(\Suc z',0,\sel(\Suc z,0,\Suc z))=\sel(\Suc z',0,\Suc z)$; both sides
equal $\Suc z$ by \rulename{IS}.  Neither premise uses its assumption.

By \cref{wit:lem:qc-0v}, with the assumption $X(0)$ unused,
$\Pi;\Gamma\proves\Q{y}X(y)$.  \rulename{QE} with the base $X(0)$ and the
target $t$ gives $X(t)$.  This is an application of \rulename{QE}, not a
substitution into the derivation of $X(\Suc z)$.
\end{proof}

\begin{lemma}[The alias]
\label{wit:lem:alias-0v}
Let $t$ be over $\Pi$ and $b\notin\Pi$.  If $\Gamma$ contains
$\Q{v}M_{\Suc t}(b,v)$, then $\Gamma\proves b=t$, by a derivation in which
no $\D{}{}$ is introduced.
\end{lemma}

\begin{proof}
The instance $M_{\Suc t}(b,0)$ is a gate at $0$ and is derivable.
\rulename{QE} at $1$ gives
$\sel(1,0,\sel(\Suc t,0,\Suc b))=\sel(1,0,\Suc t)$, and \rulename{IS},
applied twice on the left and once on the right, gives $\Suc b=\Suc t$.
By \eqref{wit:eq:inj}, $b=t$.
\end{proof}

\begin{theorem}[The literal $0$ and variables suffice]
\label{wit:thm:zero-or-variable}
For every scope $\Pi$, context $\Gamma$ and formula $A$ over $\Pi$,
\[
 \Pi;\Gamma\proves_{\Nat_{0\mathrm v}}A
 \quad\Longleftrightarrow\quad
 \Pi;\Gamma\proves_{\Nat}A .
\]
Consequently $\Nat_{0\mathrm v}$ is closed under substitution, and the
bridge of \cref{base:thm:bridge} holds for it with the same translations.
\end{theorem}

\begin{proof}
By \cref{wit:cor:names-iff} it suffices to derive
$\Pi;\varnothing\proves\D{x}{y}\mathrm{Al}_t$ for $t$ over $\Pi$.  By
\cref{wit:lem:packet-0v}, $X(\Suc t)$ is derivable.  Apply \rulename{DE} to
it at a fresh $b$; by \cref{wit:lem:alias-0v} the branch derives $b=t$.
Introduce $\D{x}{y}\mathrm{Al}_t$ at the variable $b$: the premise asks for
$\sel(\Suc z,t,b)=t$, and $\sel(\Suc z,t,b)=b$ by \rulename{IS}.  The name
does not contain $b$, and \rulename{DE} discharges it.  The two calculi have
the same terms and formulas and the comparison is the identity, so every
property of the consequence relation of $\Nat$ transfers.
\end{proof}

\begin{remark}[Versions of the packet and of the proof]
\label{wit:rem:packet-versions}
The formula $X(y)$ is $H(y)$ in \pkg{no\_witness}, $X(y)$ in
\pkg{newest\_parameter} and $A(r)$ in \pkg{uniform\_raw}.
\pkg{parameter\_witness} uses the ungated packet
\[
 T(v)\defeq\D{x}{y}\bigl(\sel(y,v,\sel(v,0,\Suc x))=v\bigr),
\]
whose instance at $y=0$ is $v=v$ and whose instances at successors are
$\sel(v,0,\Suc x)=v$; it serves in a calculus where the mode of the
introduction depends on the witness (\cref{wit:thm:R}).  The restriction
note of \pkg{no\_witness} and the appendix of \pkg{parameter\_witness} prove
\cref{wit:thm:zero-or-variable} directly, with the naming matrix
$\D{x}{v}\gamma_v(x=t)$ and a derived rule
\[
 \infrule{\Pi,b;\ \Gamma,b=t\proves F}{\Pi;\Gamma\proves F}\qquad
 b\notin\FV(\Gamma,t,F),
\]
obtained by opening $X(\Suc t)$ and cutting the alias.  The proof above
factors the same argument through \cref{wit:thm:named}.  The theorem
answers the question whether \rulename{DI} at the literal $0$ or at a
variable is enough.
\end{remark}

\begin{remark}[Lean status]
\label{wit:rem:lean-complete}
\Cref{wit:thm:zero-or-variable} is formalized:
\Lean{dwprf\_iff\_nprf}, in the module \Lean{NativeWitnessComplete}, states
that \Lean{DWPrf true true} derives $\Gamma\proves A$ if and only if
\Lean{NPrf} does, that is, the kernel with \rulename{DI} by cases at the
literal $0$ or at a variable, and at no other term, derives exactly what
the kernel derives.  The names are derived there (\Lean{DWPrf.names}) from
the coverage packet in a form in which the predecessor
stands inside the matrix: the name
$\D{x}{v}\mathrm{Al}_{\operatorname{pred}(y)}$ of
$\operatorname{pred}(y)$ holds at $y=0$ with the witness $0$ and at
$y=\Suc z$ with the witness $z$, by the closure by cases with the dummy
witness $0$ (\Lean{DWPrf.namePQ}); it is used at $y=\Suc a$, and opening
that name gives a variable equal to $\operatorname{pred}(\Suc a)=a$
(\Lean{DWPrf.nameOf}).  The packet $X(y)$ applies the predecessor after the
opening instead; the two differ in this order only.  The deterministic
selection of \cref{wit:sec:newest}, with the regeneration of derivations
under an extension of the scope, and the calculi built on it are not
formalized (\cref{wit:q:lean-newest}).
\end{remark}

\subsection{A mode selected by the witness}
\label{wit:sec:R}

\pkg{parameter\_witness} states the restriction for a calculus with one
introduction whose mode is read off the witness.  Let $\mathsf{R}$ be the
seven-rule calculus $\mathsf{W}$ of \cref{kern:def:W}, with scopes, in which
the witness $t$ of its introduction is the literal $0$ or a parameter of
$\Pi$:
\[
 \infrule{\Pi,z;\ \Gamma,B(t,\eta_{B,t}(z))\proves B(t,\Suc z)}
         {\Pi;\Gamma\proves\D{x}{y}B(x,y)}\ \rulename{RDI},
\]
where $\eta_{B,t}(z)$ is $z$ if $B$ is a literal equation and $t$ is the
literal $0$, and $0$ otherwise.  There is no rule \rulename{QI}.  Every
introduction at a variable is by cases.

\begin{theorem}[The calculus with a witness-selected mode]
\label{wit:thm:R}
$\mathsf{R}$ and $\Nat$ derive the same judgments.
\end{theorem}

\begin{proof}
Every derivation of $\mathsf{R}$ is a derivation of $\mathsf{W}$, which is
equivalent to $\Nat$ by \cref{kern:thm:witness}.  For the converse, the
proof of \cref{wit:thm:zero-or-variable} is repeated with the following
changes.

\emph{Structural rules.}  \Cref{wit:lem:structural-0v} holds for
$\mathsf{R}$: a renaming sends variables to variables, fixes the literal $0$
and preserves the root of a matrix, hence the mode.

\emph{Closure by progress.}  If $\Pi,z;\Gamma,E(z)\proves E(\Suc z)$ for a
literal equation $E$, then $\Pi;\Gamma\proves\Q{y}E(y)$: \rulename{RDI} at
the witness $0$ on $\D{x}{y}E(y)$ with $x$ not free in $E$ is in the mode by
progress, and \rulename{DE} returns the opened assumption.  With
\rulename{QE} this is equation induction, so every derivation of
$\mathcal E$ is a derivation of $\mathsf{R}$ rule by rule, and
\eqref{wit:eq:inj} holds in $\mathsf{R}$.

\emph{Closure by cases.}  The introduction of $\D{r}{y}B(y)$ at $0$ is by
progress when $B$ is an equation, so a variable witness is used for every
$B$.  The formula $\D{r}{w}(0=0)$ is derivable by \rulename{RDI} at $0$, its
premise being reflexivity.  Open it at a fresh $b$ and introduce
$\D{r}{y}B(y)$ at $b$; this is by cases.  Eliminate the unused coordinate to
obtain $\Q{y}B(y)$ and discharge $b$.  This replaces
\cref{wit:lem:qc-0v}.

\emph{Leibniz rule.}  \Cref{wit:lem:leibniz-0v} holds with the same proof;
the reintroduction in the case of $\D{}{}$ is at a variable and is by cases.

\emph{Packet.}  With $T$ as in \cref{wit:rem:packet-versions}: $T(0)$ is
introduced at $0$ by progress, its premise asking for
$\sel(\Suc z,0,\sel(0,0,1))=0$, which holds by computation; $T(\Suc z)$ is
introduced at $z$ by cases, its premise asking for
$\sel(\Suc z',\Suc z,\sel(\Suc z,0,\Suc z))=\Suc z$, which holds by
computation.  Closure by cases gives $\Q{v}T(v)$ and \rulename{QE} gives
$T(t)$ for every $t$.  Opening $T(\Suc t)$ at a fresh $b$ gives a closure
whose instance at $0$ is $\sel(0,\Suc t,\cdot)=\Suc t$, an instance of
\rulename{I0}; \rulename{QE} at $1$ and two computations give
$\Suc b=\Suc t$, hence $b=t$.

\emph{Recovery of the rules of $\Nat$.}  \rulename{QI} is closure by
progress.  For \rulename{DI} at an arbitrary term $t$, with an arbitrary
matrix, open $T(\Suc t)$ at $b$, introduce at $b$, which is by cases, and
transport the assumption and the conclusion as in the proof of
\cref{wit:thm:named}.  This includes the introduction by cases of an
equation at the literal $0$, which \rulename{RDI} does not offer directly.
\end{proof}

% ---------------------------------------------------------------------
\section{The newest parameter}
\label{wit:sec:newest}

This section draws on \pkg{no\_witness} and \pkg{newest\_parameter} for the
eight-rule calculus, on \pkg{no\_witness} for the seven-rule diagonal
variant, on \pkg{parameter\_witness} for the one-bit rule and its clause
calculus, and on \pkg{local\_introduction} for the Alias/Back calculus.

In $\Nat_{\nu}$ the introduction takes no term.  The proof of
\cref{wit:thm:zero-or-variable} uses renaming, weakening to a larger scope
and cut, none of which is available for the derivations of $\Nat_{\nu}$
(\cref{wit:ex:weakening}).  The remedy is to work with families of
derivations that are constructed anew in every scope.  A local derivation
whose witness is a freshly chosen case variable $z$ is used while $z$ is the
newest parameter; it is never moved to a scope in which $z$ is older.

\subsection{Stable derivations and hereditary families}

\begin{definition}[Stable derivations]
\label{wit:def:stable}
An application of \rulename{DI}$_\nu$ with the matrix $B(x,y)$ is
\emph{dummy} if $x$ is not free in $B$.  A derivation of $\Nat_{\nu}$ is
\emph{stable} if all its introductions are dummy.
\end{definition}

The premise of a dummy introduction does not depend on the selected witness.

\begin{lemma}[Stability]
\label{wit:lem:stable}
Let $\Pi;\Gamma\proves A$ have a stable derivation.  Let $\theta$ assign to
each parameter of $\Pi$ a term over a scope $\Pi'$, and let $\Delta$ be a
context over $\Pi'$ that contains $\Gamma[\theta]$.  Then
$\Pi';\Delta\proves A[\theta]$ has a stable derivation.  If
$\Pi;\Gamma\proves A$ and $\Pi;\Gamma,A\proves B$ have stable derivations,
so has $\Pi;\Gamma\proves B$.
\end{lemma}

\begin{proof}
The first statement is proved as the substitution lemma of
\cref{base:sec:structural}, by induction on the derivation for all $\theta$
and $\Delta$.  At a dummy introduction the premise is
$\Gamma,B(0)\proves B(\Suc z)$ whatever the witness is, and the rule applies
in the scope $\Pi'$ to the matrix $B[\theta]$, in which $x$ is still not
free.  The second statement is the usual induction on the derivation of
$\Pi;\Gamma,A\proves B$; the inserted derivation of $A$ is moved to the
scope of each leaf by the first statement.
\end{proof}

Derivations without any introduction of $\D{}{}$ are stable.  In particular
the judgments \eqref{wit:eq:inj} and the derivation of
\cref{wit:lem:alias-0v} are stable in $\Nat_{\nu}$.

\begin{definition}[Hereditary derivability]
\label{wit:def:hereditary}
Write $\Pi;\Gamma\Vdash A$ if for every scope $\Pi'$, every injective map
$\rho$ from $\Pi$ into $\Pi'$ and every context $\Delta$ over $\Pi'$ that
contains $\rho\Gamma$, the calculus $\Nat_{\nu}$ derives
$\Pi';\Delta\proves\rho A$.
\end{definition}

Hereditary derivability is a property of a family of derivations, one for
each $(\Pi',\rho,\Delta)$; it is not a rule of the calculus.  The member of
the family at the identity is a derivation of $\Pi;\Gamma\proves A$.  The
families below are given by constructions that take the receiving scope and
context as input and choose their eigenvariables there.

\begin{lemma}[Closure properties]
\label{wit:lem:hereditary-closure}
\begin{enumerate}[label=\textup{(\roman*)}]
\item If $A\in\Gamma$, then $\Pi;\Gamma\Vdash A$.
\item If $\Pi;\Gamma\Vdash A$, $\rho$ is an injective map from $\Pi$ into
 $\Pi'$ and $\Delta\supseteq\rho\Gamma$ is over $\Pi'$, then
 $\Pi';\Delta\Vdash\rho A$.
\item If $\Pi;\Gamma\proves A$ has a stable derivation, then
 $\Pi;\Gamma\Vdash A$.
\item The rules \rulename{I0}, \rulename{IS}, \rulename{Tr} and
 \rulename{QE} hold with $\Vdash$ in place of $\proves$.
\item If $\Pi;\Gamma\Vdash\D{x}{y}B$ and
 $\Pi,b;\ \Gamma,\Q{y}B(b,y)\Vdash A$ with $A$ over $\Pi$, then
 $\Pi;\Gamma\Vdash A$.
\item If $\Pi,z;\ \Gamma,E(z)\Vdash E(\Suc z)$ for a literal equation $E$,
 then $\Pi;\Gamma\Vdash\Q{y}E(y)$.
\end{enumerate}
\end{lemma}

\begin{proof}
(i) is \rulename{Hyp} in every receiving context.  (ii) holds because
injective maps compose.  (iii) is \cref{wit:lem:stable}.  For (iv), apply
the rule to the members of the given families at $(\Pi',\rho,\Delta)$, with
the terms renamed by $\rho$.  For (v), let $(\Pi',\rho,\Delta)$ be given.
Choose $b'\notin\Pi'$ and extend $\rho$ by $b\mapsto b'$.  The member of the
second family at $(\Pi',b')$ and the context $\Delta,\Q{y}\rho B(b',y)$
derives $\rho A$, and \rulename{DE} applies to it and to the member of the
first family.  (vi) is proved in the same way with \rulename{QI}.
\end{proof}

\begin{lemma}[Uniform cut]
\label{wit:lem:uniform-cut}
If $\Pi;\Gamma\Vdash A$ and $\Pi;\Gamma,A\proves B$ in $\Nat_{\nu}$, then
$\Pi;\Gamma\proves B$ in $\Nat_{\nu}$.  Consequently, if
$\Pi;\Gamma\Vdash A$ and $\Pi;\Gamma,A\Vdash B$, then $\Pi;\Gamma\Vdash B$.
\end{lemma}

\begin{proof}
By induction on the derivation of $\Pi;\Gamma,A\proves B$, for all $\Pi$ and
$\Gamma$.  At \rulename{Hyp} for $A$, take the member of the family at the
identity; any other \rulename{Hyp} is kept.  A rule whose premises have the
same scope and context is reapplied to the results of the induction
hypothesis.  At a rule with an eigenvariable, say \rulename{DI}$_\nu$ with
the premise $\Pi,z;\ \Gamma,A,B(w,0)\proves B(w,\Suc z)$ and $w=\nu(\Pi)$,
\cref{wit:lem:hereditary-closure}(ii) gives
$\Pi,z;\ \Gamma,B(w,0)\Vdash A$, the induction hypothesis gives
$\Pi,z;\ \Gamma,B(w,0)\proves B(w,\Suc z)$, and the rule is reapplied in the
scope $\Pi$, where it selects the same $w$.  \rulename{DE} and
\rulename{QI} are treated in the same way.  Every node of the receiving
derivation keeps its scope, and no derivation is moved to another scope:
below a new parameter the family supplies a new member.

For the second statement let $(\Pi',\rho,\Delta)$ be given.  The second
family has a member deriving $\Pi';\Delta,\rho A\proves\rho B$, and
$\Pi';\Delta\Vdash\rho A$ by \cref{wit:lem:hereditary-closure}(ii).
\end{proof}

\subsection{Closure by cases, the Leibniz rule and equation cases}

\begin{lemma}[Closure by cases]
\label{wit:lem:qc}
In $\Nat_{\nu}$, if $\Pi,z;\ \Gamma,B(0)\proves B(\Suc z)$, then
$\Pi;\Gamma\proves\Q{y}B(y)$; if the given derivation is stable, so is the
result.  If $\Pi,z;\ \Gamma,B(0)\Vdash B(\Suc z)$, then
$\Pi;\Gamma\Vdash\Q{y}B(y)$.
\end{lemma}

\begin{proof}
Choose $r$ not free in $B$ and apply \rulename{DI}$_\nu$ to $\D{r}{y}B(y)$.
The introduction is dummy, and its premise is the given one whatever
$\nu(\Pi)$ is.  \rulename{DE} at a fresh $b$ adds the assumption
$\Q{y}B(y)$, which is the conclusion.  For the hereditary statement, the
member of the given family at $(\Pi',z')$, $\rho[z\mapsto z']$ and
$\Delta,\rho B(0)$ is the premise needed in the scope $\Pi'$.
\end{proof}

\begin{lemma}[Hereditary Leibniz rule]
\label{wit:lem:hl}
For every native formula $F(v)$: if $\Pi;\Gamma\Vdash s=t$ and
$\Pi;\Gamma\Vdash F(s)$, then $\Pi;\Gamma\Vdash F(t)$.
\end{lemma}

\begin{proof}
By induction on the number of binders of $F$, for all scopes, contexts and
terms.  For an equation the claim follows from
\cref{wit:lem:hereditary-closure}(iv), as in \cref{wit:lem:leibniz-0v}.

Let $F(v)=\Q{y}B(v,y)$.  By \cref{wit:lem:qc} it suffices to show
$\Theta\Vdash B(t,\Suc z)$, where $\Theta$ stands for the scope $\Pi,z$ with
the context $\Gamma,B(t,0)$.  By \cref{wit:lem:hereditary-closure},
$\Theta\Vdash t=s$, $\Theta\Vdash B(t,0)$ and $\Theta\Vdash\Q{y}B(s,y)$.
The induction hypothesis gives $\Theta\Vdash B(s,0)$, \rulename{QE} at
$\Suc z$ gives $\Theta\Vdash B(s,\Suc z)$, and the induction hypothesis
gives $\Theta\Vdash B(t,\Suc z)$.

Let $F(v)=\D{x}{y}B(v,x,y)$, and let $(\Pi',\rho,\Delta)$ be given; write
$s',t',B'$ for the images under $\rho$.  The first family gives
$\Pi';\Delta\proves\D{x}{y}B'(s',x,y)$.  Apply \rulename{DE} at a fresh $b$.
In the scope $\Pi',b$ the parameter $b$ is the newest, so
\rulename{DI}$_\nu$ for $\D{x}{y}B'(t',x,y)$ selects $b$; its premise has
the scope $\Pi',b,z$ and the context
$\Delta,\Q{y}B'(s',b,y),B'(t',b,0)$, denoted $\Theta$.  As in the previous
case, $\Theta\Vdash B'(s',b,0)$ by the induction hypothesis,
$\Theta\Vdash B'(s',b,\Suc z)$ by \rulename{QE}, and
$\Theta\Vdash B'(t',b,\Suc z)$ by the induction hypothesis.  The member at
the identity is the premise of the introduction.  The conclusion does not
contain $b$, and \rulename{DE} discharges it.  Whatever the receiving scope
is, the parameter $b$ is generated after all its parameters, so the same
recipe always introduces at the newest parameter.
\end{proof}

The derivation of \eqref{wit:eq:pred} needs only equation induction.
\pkg{newest\_parameter} observes in addition that case distinctions on a
term of the context are available in $\Nat_{\nu}$ for equational
conclusions, without \rulename{DI} at a chosen witness.  The statement is
\cref{guard:lem:np-imp} for hereditary families, with the implication
$U\Rightarrow_{=}V$ between equations defined there.

\begin{lemma}[Contextual equation cases]
\label{wit:lem:eqcases}
In $\Nat_{\nu}$, for equations $U$, $V$, $F$ and a term $e$ over $\Pi$:
\begin{enumerate}[label=\textup{(\roman*)}]
\item if $\Pi;\Gamma,U\Vdash V$, then
 $\Pi;\Gamma\Vdash U\Rightarrow_{=}V$;
\item if $\Pi;\Gamma\Vdash U\Rightarrow_{=}V$ and
 $\Pi;\Gamma\Vdash U$, then $\Pi;\Gamma\Vdash V$;
\item if $\Pi;\Gamma,e=0\Vdash F$ and
 $\Pi,z;\ \Gamma,e=\Suc z\Vdash F$, then $\Pi;\Gamma\Vdash F$.
\end{enumerate}
\end{lemma}

\begin{proof}
The proof of \cref{guard:lem:np-imp} applies, with \cref{wit:lem:qc} for
the closure by cases and \cref{wit:lem:uniform-cut} for the cut; in (iii)
the two hypotheses give $(e=y)\Rightarrow_{=}F$ at $y=0$ and at
$y=\Suc z$ by (i), the closure by cases and \rulename{QE} give it at $y=e$,
and (ii) with $e=e$ gives $F$.
\end{proof}

All introductions in the constructions of
\cref{wit:lem:qc,wit:lem:eqcases} are dummy; applied to stable derivations
they yield stable derivations.

\subsection{The packet in every scope}

\begin{lemma}[The packet with the newest parameter]
\label{wit:lem:packet}
In $\Nat_{\nu}$, for every scope $\Pi$, every context $\Gamma$ and every
term $t$ over $\Pi$: $\Pi;\Gamma\proves X(\Suc t)$.  Hence
$\Pi;\Gamma\Vdash X(\Suc t)$.
\end{lemma}

\begin{proof}
$X(0)$: the rule selects $w=\nu(\Pi)$, and its premise asks for
$M_0(w,\Suc z)$, that is, for
$\sel(\Suc z,0,\sel(0,0,\Suc w))=\sel(\Suc z,0,0)$.  Both sides equal $0$ by
computation, whatever $w$ is.  So $X(0)$ is derivable in the empty scope and
also when $w$ is an unrelated parameter.

$\D{r}{y}X(y)$ with $r$ unused: the premise of this introduction has the
scope $\Pi,z$ and asks for $X(\Suc z)$ from $X(0)$.  In the scope $\Pi,z$
the newest parameter is $z$, so the introduction of $X(\Suc z)$ selects $z$;
its premise asks for $M_{\Suc z}(z,\Suc z')$, both sides of which equal
$\Suc z$ by computation.  The assumptions of both premises are unused.

\rulename{DE} on $\D{r}{y}X(y)$ gives $\Q{y}X(y)$, and \rulename{QE} with
the base $X(0)$ and the target $\Suc t$ gives $X(\Suc t)$.  The construction
works in every scope and every context, and
$\rho X(\Suc t)=X(\Suc\rho t)$; this is the hereditary statement.
\end{proof}

The construction is applied to the whole closure.  It is not claimed that
$\Pi,z;\Gamma\Vdash X(\Suc z)$: a derivation that introduces $X(\Suc z)$ at
$z$ is not a derivation in a scope in which $z$ is no longer the newest
parameter.  In each receiving scope the case variable is chosen last.

\begin{lemma}[The hereditary alias]
\label{wit:lem:alias}
Let $\Theta$ be a scope and context such that the context contains
$\Q{v}M_{\Suc t}(b,v)$.  Then $\Theta\Vdash b=t$.
\end{lemma}

\begin{proof}
The derivation of \cref{wit:lem:alias-0v} is stable.
\end{proof}

\subsection{The equivalence theorem}

\begin{theorem}[No witness in the eight-rule calculus]
\label{wit:thm:newest}
For every scope $\Pi$, context $\Gamma$ and formula $A$ over $\Pi$,
\[
 \Pi;\Gamma\proves_{\Nat_{\nu}}A
 \quad\Longleftrightarrow\quad
 \Pi;\Gamma\proves_{\Nat}A .
\]
Both directions are effective transformations of derivations, and the terms
and formulas of the judgment are unchanged.
\end{theorem}

\begin{proof}
From left to right every rule of $\Nat_{\nu}$ is an instance of a rule of
$\Nat$.  For the converse the following is proved by induction on a
derivation of $\Nat$:
\[
 \text{if }\Pi;\Gamma\proves_{\Nat}A\text{, then }\Pi;\Gamma\Vdash A .
\]
\rulename{Hyp}, \rulename{I0}, \rulename{IS}, \rulename{Tr} and
\rulename{QE} are \cref{wit:lem:hereditary-closure}(i) and (iv).
\rulename{DE} and \rulename{QI} are
\cref{wit:lem:hereditary-closure}(v) and (vi), applied to the families that
the induction hypothesis gives for the premises.

Consider \rulename{DI} with the witness $t$ and the matrix $B$.  The
induction hypothesis is
\begin{equation}
 \Pi,z;\ \Gamma,B(t,0)\Vdash B(t,\Suc z).
 \label{wit:eq:ih}
\end{equation}
Let $(\Pi',\rho,\Delta)$ be given and write $t'=\rho t$, $B'=\rho B$.
By \cref{wit:lem:packet}, $\Pi';\Delta\proves X(\Suc t')$.  Apply
\rulename{DE} to it at a fresh $b$; the branch has the scope $\Pi',b$ and
the context $\Delta_b=\Delta,\Q{v}M_{\Suc t'}(b,v)$.  In this branch apply
\rulename{DI}$_\nu$ to $\D{x}{y}B'$.  The newest parameter is $b$, so the
rule selects $b$, and its eigenvariable $z'$ is appended after this
selection.  Let $\Theta$ be the scope $\Pi',b,z'$ with the context
$\Delta_b,B'(b,0)$; the premise to be derived is
$\Theta\proves B'(b,\Suc z')$.
\begin{enumerate}[label=\textup{(\alph*)}]
\item $\Theta\Vdash b=t'$ by \cref{wit:lem:alias}, and
 $\Theta\Vdash B'(b,0)$ by membership.  By \cref{wit:lem:hl},
 $\Theta\Vdash B'(t',0)$.
\item The map $\rho[z\mapsto z']$ from $\Pi,z$ into $\Pi',b,z'$ is
 injective, and the context of $\Theta$ extended by $B'(t',0)$ contains the
 image of $\Gamma,B(t,0)$.  By \eqref{wit:eq:ih} and
 \cref{wit:lem:hereditary-closure}(ii),
 $\Theta,B'(t',0)\Vdash B'(t',\Suc z')$.
\item By (a), (b) and \cref{wit:lem:uniform-cut},
 $\Theta\Vdash B'(t',\Suc z')$.
\item By \cref{wit:lem:hl} with the equation $t'=b$,
 $\Theta\Vdash B'(b,\Suc z')$.
\end{enumerate}
The member of the last family at the identity is the premise.  The
introduction yields $\Pi',b;\Delta_b\proves\D{x}{y}B'$, a formula over
$\Pi'$, and \rulename{DE} discharges $b$.  This was done for an arbitrary
$(\Pi',\rho,\Delta)$, so $\Pi;\Gamma\Vdash\D{x}{y}B$.

The construction is well founded.  The induction hypothesis is used for the
premises of the given derivation only.  \Cref{wit:lem:hl} is an induction on
formulas.  The packet, the alias and the predecessor are fixed
constructions that precede the theorem.  No derivation of $\Nat_{\nu}$ is
moved to another scope and no term is substituted into one.  Finally the
member at the identity of $\Pi;\Gamma\Vdash A$ is the required derivation.
\end{proof}

\begin{corollary}[Structural rules and the three calculi]
\label{wit:cor:structural}
$\Nat_{\nu}$, $\Nat_{0\mathrm v}$ and $\Nat$ derive the same judgments.
Substitution, weakening to a larger scope and cut are admissible in
$\Nat_{\nu}$ as transformations of derivable judgments.  The ordinary
realization of $\Nat_{\nu}$ is bi-interpretable and equiconsistent with
$\iPRA$ by the translations of \cref{base:thm:bridge}, with whole domains,
actual equality and the identity graph $x=y$ in both numerical composites.
\end{corollary}

\begin{proof}
The first statement is \cref{wit:thm:newest} and \eqref{wit:eq:sandwich}.
For the structural rules translate to $\Nat$, apply its rule
(\cref{base:sec:structural}) and translate back; they are consequences of
the theorem, not steps of its proof, and they do not say that an unchanged
derivation tree is a derivation after an extension of the scope.  The
comparison with $\Nat$ is the identity on terms, formulas and contexts, so
the translations of the bridge, the derivations of both formula round trips
and the preservation of $0=\Suc 0$ transfer unchanged.
\end{proof}

\begin{remark}[The three formulations of the proof families]
\label{wit:rem:families}
\pkg{no\_witness} calls the families \emph{portable}: a procedure that, for
every extension of the scope by fresh parameters and of the context by
assumptions, constructs a derivation of the unchanged formula.
\pkg{newest\_parameter} calls them \emph{persistent builders}, with the same
definition.  Both state the induction on the source derivation for an
arbitrary substitution $\sigma$ of terms over a target scope for the source
parameters, a target context $\Delta$, and a family for $E[\sigma]$ for each
source assumption $E$; the result is a family for $A[\sigma]$.  The
substitution takes the place of the renaming $\rho$, and stability under
extension takes the place of stability under injective maps.
\pkg{parameter\_witness} and \pkg{uniform\_raw} use hereditary derivability
as in \cref{wit:def:hereditary}, and the proof above follows them.
\Cref{wit:lem:uniform-cut} is the persistent cut of
\pkg{newest\_parameter} and the uniform cut of \pkg{parameter\_witness}.
\pkg{no\_witness} also derives the standalone name
$\D{x}{v}\gamma_v(x=t)$ in $\Nat_{\nu}$: while $b=t$ is available and $b$ is
the newest parameter, introduce the name, and then discharge the opening of
$X(\Suc t)$.  None of the packages claims a Lean proof of
\cref{wit:thm:newest}; a formalization has to encode the families, or an
equivalent protected translation, and cannot assume the weakening of
restricted derivations.
\end{remark}

\subsection{The seven-rule diagonal variant}
\label{wit:sec:diagonal}

The diagonal calculus is the seven-rule calculus $\Nat_{\mathrm{diag}}$ of
\cref{guard:thm:diagonal}: the calculus $\Nat_P$ of
\cref{guard:def:selector} for the selector $P$ that holds of the matrices
diagonal in $y$ (\cref{guard:def:shapes}).  Its one introduction
$\rulename{DI}_P$ takes a witness $t$ and reads its mode off the raw matrix,
before the witness is inserted; write $\eta_B(z)$ for $\eta^{P}_B(z)$.  It
derives the same judgments as $\Nat$ (\cref{guard:thm:diagonal}; Lean:
\Lean{gprf\_iff\_nprf}).  Let $\Nat_{\mathrm{diag},\nu}$ be
$\Nat_{\mathrm{diag}}$ with scopes and with $t=\nu(\Pi)$:
\[
 \infrule{\Pi,z;\ \Gamma,B(\nu(\Pi),\eta_B(z))\proves B(\nu(\Pi),\Suc z)}
         {\Pi;\Gamma\proves\D{x}{y}B(x,y)} .
\]
The matrix selects the mode and the scope selects the witness.  No guard,
axiom or auxiliary syntax is added.

\begin{theorem}[No witness in the seven-rule calculus]
\label{wit:thm:diagonal}
$\Nat_{\mathrm{diag},\nu}$, $\Nat_{\mathrm{diag}}$ and $\Nat$ derive the
same judgments.
\end{theorem}

\begin{proof}
Every rule of $\Nat_{\mathrm{diag},\nu}$ is a rule of
$\Nat_{\mathrm{diag}}$.  For the converse,
\cref{wit:def:stable,wit:def:hereditary} and
\cref{wit:lem:stable,wit:lem:hereditary-closure,wit:lem:uniform-cut} are
taken for $\Nat_{\mathrm{diag},\nu}$, with item (vi) of
\cref{wit:lem:hereditary-closure} omitted.  Their proofs are unchanged,
because the diagonal shape is preserved and reflected by the substitution
of terms for parameters (\cref{guard:thm:diagonal}(iii)).  The following
steps differ from the proof of \cref{wit:thm:newest}; they are the steps of
\cref{guard:thm:selector} carried out for families of derivations, since
the structural rules of $\Nat_{\mathrm{diag}}$ are not available under the
restriction.

\emph{Closure by cases.}  Let $\Pi,z;\Gamma,B(0)\proves B(\Suc z)$ and $r$
not free in $B$.  If $B$ is not diagonal, the introduction of
$\D{r}{y}B(y)$ is by cases and \cref{wit:lem:qc} applies.  If $B$ is
diagonal, $B(0)$ is derivable in every scope and context by a stable
derivation; this is condition (S1) of \cref{guard:thm:diagonal}(i).
\Cref{wit:lem:uniform-cut} removes the assumption $B(0)$ from the given
premise, and \cref{wit:lem:ctxweak}, which holds for
$\Nat_{\mathrm{diag},\nu}$ with the same proof, adds the assumption $B(z)$.
The result is the premise of the introduction by progress.  In both cases
\rulename{DE} yields $\Q{y}B(y)$.

\emph{Equation induction.}  There is no rule \rulename{QI}.  For a term
$s(y)$ let $L_s(y)\defeq\It{y}{y}{u}{s(y)}$ with $u$ not free in $s$, so
that $L_s(y)=L_t(y)$ is the diagonal lift $\Delta_yE$ of
\cref{guard:def:shapes} for $E(y)\equiv(s(y)=t(y))$; then $L_s(0)=0$ and
$L_s(\Suc z)=s(\Suc z)$ by computation.  The decoder equation
\begin{equation}
 \sel\bigl(y,s(0),L_s(y)\bigr)=s(y)
 \label{wit:eq:dec}
\end{equation}
holds at $y=0$ by \rulename{I0} and at $y=\Suc z$ by two successor
computations; closure by cases and \rulename{QE} give it at every term, in
every scope and context.  Let $E(y)\equiv(s(y)=t(y))$ and suppose
$\Pi,z;\Gamma,E(z)\Vdash E(\Suc z)$.  Then $\Pi;\Gamma\Vdash\Q{y}E(y)$.
Indeed, let $(\Pi',\rho,\Delta)$ be given, and drop $\rho$ from the
notation.  By closure by cases it suffices to derive $E(\Suc z)$ in the
scope $\Pi',z$ from $\Delta_0=\Delta,E(0)$.  The matrix
$L_s(y)=L_t(y)$ is diagonal.  Introduce $\D{r}{y}(L_s(y)=L_t(y))$ with $r$
unused; this is by progress, and the premise has a further eigenvariable
$z_1$ and the context $\Delta_0,L_s(z_1)=L_t(z_1)$, denoted $\Theta$.  By
\eqref{wit:eq:dec} for $s$ and for $t$, the assumption and the base $E(0)$,
\[
 \Theta\Vdash s(z_1)=\sel(z_1,s(0),L_s(z_1))
       =\sel(z_1,t(0),L_t(z_1))=t(z_1).
\]
The step family, at the map $z\mapsto z_1$, gives
$\Theta,E(z_1)\Vdash E(\Suc z_1)$; uniform cut gives
$\Theta\Vdash E(\Suc z_1)$, and computation gives
$\Theta\Vdash L_s(\Suc z_1)=L_t(\Suc z_1)$.  After the introduction,
\rulename{DE} adds the closure $\Q{y}(L_s(y)=L_t(y))$.  Its instance at $0$
is derivable, \rulename{QE} at $\Suc z$ gives
$L_s(\Suc z)=L_t(\Suc z)$, and computation gives $E(\Suc z)$.  With
\rulename{QE} this is equation induction for hereditary families.  Every
rule of the fragment $\mathcal E$ therefore preserves $\Vdash$, and by
induction on a derivation of $\mathcal E$ the judgments \eqref{wit:eq:inj}
hold with $\Vdash$.  The construction uses closure by cases, computations
and the equality rules; it uses no predecessor and no chosen witness.

\emph{Leibniz rule.}  In the case $F(v)=\D{x}{y}B(v,x,y)$ of
\cref{wit:lem:hl} the target is introduced at the opened parameter $b$.  If
its matrix is not diagonal, the proof is unchanged.  If it is diagonal, the
premise has the assumption $B'(t',b,z)$ in place of $B'(t',b,0)$; the
latter is derivable outright, and the same three steps give
$B'(t',b,\Suc z)$ without the assumption.

\emph{Packet.}  The matrices $M_y(x,v)$ have outer iterators with the seed
$0$, and the diagonal shape requires the seed to be the bound variable; the
matrix of $\D{r}{y}X(y)$ is not an equation.  All introductions in the
proof of \cref{wit:lem:packet} are therefore by cases, and the proof is
unchanged.

\emph{Translation.}  The induction is on a derivation of $\Nat_{\mathrm{diag}}$.
The mode of an introduction is the same in source and target, because the
alias changes the selected witness and not the raw matrix.  In the mode by
progress, step (a) of the proof of \cref{wit:thm:newest} transports
$B'(b,z')$ to $B'(t',z')$, step (b) uses the induction hypothesis with the
assumption $B(t,z)$, and steps (c) and (d) are unchanged.

The equivalence with $\Nat$ is the cited one for $\Nat_{\mathrm{diag}}$.
\end{proof}

\subsection{One bit}
\label{wit:sec:bit}

Let $\mathsf{R}_{\nu}$ be the calculus $\mathsf{R}$ of \cref{wit:sec:R}
with the variable witnesses restricted to $\nu(\Pi)$ in a nonempty scope;
the literal $0$ remains available, and in the empty scope it is the only
witness.  For an implementation both choices are made total.  The calculus
$\mathsf{R}_{\beta}$ has, in place of \rulename{RDI}, the rule
\begin{gather*}
 \infrule{\Pi,z;\ \Gamma,B(t_\beta,\eta)\proves B(t_\beta,\Suc z)}
         {\Pi;\Gamma\proves\D{x}{y}B(x,y)}\ \rulename{Bit-DI}
 \\[1.2ex]
 t_\beta=\begin{cases}0,&\beta=0,\\ \nu(\Pi),&\beta=1,\end{cases}
 \qquad
 \eta=\begin{cases}
  z,&\beta=0\text{ and $B$ is a literal equation,}\\
  0,&\text{otherwise,}
 \end{cases}
\end{gather*}
whose only datum is the bit $\beta\in\{0,1\}$.  The mode is selected by the
bit and not by an inspection of the selected term: with $\beta=1$ in the
empty scope the rule is an introduction by cases at $0$, also for an
equation.

\begin{theorem}[Zero or the newest parameter]
\label{wit:thm:R-newest}
$\mathsf{R}_{\nu}$, $\mathsf{R}_{\beta}$ and $\Nat$ derive the same
judgments.
\end{theorem}

\begin{proof}
Every rule of $\mathsf{R}_{\nu}$ is an instance of \rulename{Bit-DI}: an
introduction at the literal $0$ is the bit $0$, with the same mode, and an
introduction at the newest parameter is the bit $1$.  Every instance of
\rulename{Bit-DI} is an introduction by cases or an introduction by progress
at the witness $0$, and both are rules of $\Nat$, the second by
\cref{base:cor:combined}.  It remains to show
$\Nat\subseteq\mathsf{R}_{\nu}$.
Hereditary derivability and
\cref{wit:lem:hereditary-closure,wit:lem:uniform-cut} are taken for
$\mathsf{R}_{\nu}$; a renaming preserves the mode of an introduction, and
the proof of uniform cut is unchanged because each node keeps its scope.
The following steps differ from the proof of \cref{wit:thm:newest}.

\emph{Closure by progress and equation induction.}  If
$\Pi,z;\Gamma,E(z)\Vdash E(\Suc z)$, then $\Pi;\Gamma\Vdash\Q{y}E(y)$, by
the introduction of $\D{x}{y}E(y)$ at the literal $0$ in each receiving
scope.  This replaces \cref{wit:lem:hereditary-closure}(vi), and with
\rulename{QE} it gives \eqref{wit:eq:inj} with $\Vdash$ by induction on a
derivation of $\mathcal E$.

\emph{Closure by cases.}  Let $\Pi,z;\Gamma,B(0)\Vdash B(\Suc z)$ and let
$(\Pi',\rho,\Delta)$ be given.  The formula $\D{r}{w}(0=0)$ is derivable at
the witness $0$ by progress.  Open it at a fresh $b$, which is then the
newest parameter, and introduce $\D{r}{y}\rho B(y)$ at $b$, by cases.  Its
premise is the member of the given family in the scope $\Pi',b,z'$; it is
not obtained by renaming a derivation.  Eliminate the two unused
coordinates.

\emph{Leibniz rule.}  As \cref{wit:lem:hl}.  The reintroduction is at the
parameter just opened, which is the newest, and it is by cases.

\emph{Packet.}  With $T$ as in \cref{wit:rem:packet-versions}, $T(0)$ is
derivable in every scope by progress at $0$.  For $\Q{v}T(v)$, expand the
closure by cases: open $\D{r}{w}(0=0)$ at a fresh $b$, introduce the dummy
$\D{r}{v}T(v)$ at $b$, and in its premise, where the eigenvariable $z$ has
just been appended and is the newest parameter although $b$ was added
before, introduce $T(\Suc z)$ at $z$ by cases; the instance asked for holds
by computation.  Eliminate the two unused coordinates.  \rulename{QE} gives
$T(t)$ for every $t$, in every scope and context.  It is not claimed that
$T(\Suc z)$ has a hereditary derivation.  Opening $T(\Suc t)$ at a fresh
$b$ gives $b=t$ hereditarily, as in the proof of \cref{wit:thm:R}.

\emph{Translation.}  The induction is on a derivation of $\Nat$.
\rulename{QI} is closure by progress.  For \rulename{DI} at $t$ open
$T(\Suc\rho t)$ at a fresh $b$ and introduce at $b$, by cases; steps (a) to
(d) of the proof of \cref{wit:thm:newest} are unchanged.

Finally $\Nat\subseteq\mathsf{R}_{\nu}\subseteq\mathsf{R}_{\beta}\subseteq
\Nat$ on derivable judgments.
\end{proof}

\subsection{The five-family clause calculus}
\label{wit:sec:clause}

The checker of \pkg{parameter\_witness} uses the raw clause presentation
without the binder $\Q{}$, and the restriction is proved directly for it: a
translation of $\Nat_{\nu}$ into clauses may insert parameters and need not
preserve which parameter is the newest.  The comparison calculus
$\mathsf{C}$ is the five-rule clause calculus $\mathsf{C}$ of
\cref{claus:def:c3}, that is, $\mathsf{C}_3$ with the two-case evaluation
$\Phi_2$: rules \rulename{Refl}, \rulename{Eval}, \rulename{Use},
\rulename{Elim} and \rulename{Intro}, the last at an arbitrary witness $t$
and by progress exactly when $B$ is a literal equation and $t$ is the
literal $0$.  Its formulas are equations and $\D{x}{y}B$, its contexts are
lists of clauses (\cref{base:sec:clauses}), and an assumption $F$ is the
clause $\mathsf{a}(F)\defeq[0=0;\,v.F]$, used by \rulename{Refl} at $0$
and \rulename{Use}.  Here the scope is written explicitly
(\cref{wit:def:scope}): \rulename{Elim} appends its parameter $b$ and the
introduction its eigenvariable $z$, and the witness of \rulename{Intro} is
over $\Pi$.  The calculus $\mathsf{C}_{\beta}$ has the rules
\rulename{Refl}, \rulename{Eval}, \rulename{Use} and \rulename{Elim} of
$\mathsf{C}$ and, in place of \rulename{Intro},
\[
 \infrule{\Pi,z;\ \Gamma,\mathsf{a}(B(t_\beta,\eta))\proves
          B(t_\beta,\Suc z)}
         {\Pi;\Gamma\proves\D{x}{y}B}\ \rulename{Bit-Intro}
\]
with $t_\beta$ and $\eta$ as in \rulename{Bit-DI}.  There are no other
axioms or rules.

By \cref{claus:lem:eq}, \rulename{Refl} and the two cases of
\rulename{Eval} derive, over every clause context, symmetry, transitivity,
the two computation equations, congruence in every term template and
transport in a term and in an equation.  These derivations contain no
\rulename{Elim} and no introduction, so they are derivations of
$\mathsf{C}_{\beta}$ as well and are valid in every scope.

\begin{theorem}[The one-bit clause calculus]
\label{wit:thm:clause-bit}
$\mathsf{C}_{\beta}$ and $\mathsf{C}$ derive the same judgments
$\Pi;\Gamma\proves A$, for all clause contexts $\Gamma$.
\end{theorem}

\begin{proof}
\emph{From $\mathsf{C}_{\beta}$ to $\mathsf{C}$.}  An introduction with the
bit $0$ is the introduction of $\mathsf{C}$ at the literal $0$, with the
same mode.  With the bit $1$ in a nonempty scope it is the introduction of
$\mathsf{C}$ at a variable, by cases.  With the bit $1$ in the empty scope
it is by cases at $0$.  For a compound matrix this is the introduction of
$\mathsf{C}$ at $0$.  For an equation matrix, $\mathsf{C}$ introduces at
the term $w=\It{0}{0}{u}{0}$, which is not the literal $0$, so the mode is
by cases; the assumption $B(w,0)$ is transported along $w=0$ to $B(0,0)$,
the given premise is applied, and its conclusion is transported back to
$B(w,\Suc z)$.  The calculus $\mathsf{C}$ is closed under injective
renaming and weakening and admits the cut of an assumption
(\cref{claus:lem:struct}), so this composition is a derivation.
% ed.: the package states the alias and the transport; the structural
% facts of C that compose them are made explicit here.

\emph{From $\mathsf{C}$ to $\mathsf{C}_{\beta}$.}  Write
$\Pi;\Gamma\Vdash_{c}A$ if for every injective map $\rho$ from $\Pi$ into a
scope $\Pi'$ and every clause context $\Delta$ over $\Pi'$ that contains
$\rho\Gamma$, $\mathsf{C}_{\beta}$ derives $\Pi';\Delta\proves\rho A$.  The
renaming acts on the antecedent and on the body of each clause, and the
receiving context retains the clauses themselves.

\emph{Cut of an assumption.}  If
$\Pi;\Gamma,\mathsf{a}(A)\proves F$ and $\Pi;\Gamma\Vdash_{c}A$, then
$\Pi;\Gamma\proves F$; this is proved as \cref{wit:lem:uniform-cut}.  A
\rulename{Use} of the clause $\mathsf{a}(A)$ is replaced by the member of
the family in the scope of that leaf, its premise $0=0$ being dropped; every
other \rulename{Use} keeps its clause.  \rulename{Elim} and
\rulename{Bit-Intro} add their clauses in both contexts.  No clause other
than an assumption is cut.

\emph{Rules.}  \rulename{Refl}, \rulename{Eval} and \rulename{Use} hold
with $\Vdash_{c}$, the clause of a \rulename{Use} being a member of the
renamed context; \rulename{Elim} holds with $\Vdash_{c}$ by sending its
eigenvariable to a fresh parameter of the receiving scope and retaining the
new clause there.

\emph{Closures.}  Write $q_yA(y)$ for $\D{x}{y}A(y)$ with $x$ not free in
$A$.  Its introduction by cases is \rulename{Bit-Intro} with the bit $1$:
the witness is unused, and the fallback in the empty scope is still by
cases.  For an equation $E$ its introduction by progress is the bit $0$.
Its elimination is \rulename{Elim} followed by \rulename{Use} of the clause
$[A(0);\,y.A(y)]$.  With hereditary premises these three operations hold
for $\Vdash_{c}$, as in \cref{wit:lem:qc}.  In particular equation
induction holds for $\Vdash_{c}$, and \eqref{wit:eq:inj} holds with
$\Vdash_{c}$ by induction on a derivation of $\mathcal E$.

\emph{Leibniz rule.}  If $\Pi;\Gamma\Vdash_{c}s=t$ and
$\Pi;\Gamma\Vdash_{c}F(s)$, then $\Pi;\Gamma\Vdash_{c}F(t)$, by induction on
$F$.  For an equation this is the equality rules.  For
$\D{x}{y}B(s,x,y)$, open the formula at a fresh $b$ in the receiving scope,
retaining the clause $[B(s,b,0);\,y.B(s,b,y)]$, and introduce the target
with the bit $1$, which selects $b$.  From the assumption $B(t,b,0)$ the
induction hypothesis gives $B(s,b,0)$, \rulename{Use} of the retained
clause at $\Suc z$ gives $B(s,b,\Suc z)$, and the induction hypothesis
gives $B(t,b,\Suc z)$; the cut of an assumption composes these steps.
\rulename{Elim} discharges $b$.

\emph{Packet.}  With $T$ as in \cref{wit:rem:packet-versions}, $T(0)$ is
introduced with the bit $0$ in every scope: the matrix is an equation, the
introduction is by progress, and the instance asked for holds by
computation.  For $q_vT(v)$ use the bit $1$ on the outer formula; in its
premise the eigenvariable $z$ is the newest parameter, and $T(\Suc z)$ is
introduced with the bit $1$, which selects $z$.  The construction is
repeated in every scope.  \rulename{Elim} and \rulename{Use} give
$\Pi;\Gamma\Vdash_{c}T(t)$ for every $t$.  Opening $T(\Suc t)$ at a fresh
$b$ and using its clause at $1$ gives $\Suc b=\Suc t$, hence $b=t$, with
$\Vdash_{c}$ from the opened clause context.

\emph{Translation.}  By induction on a derivation of $\mathsf{C}$ its
conclusion is shown to be hereditarily derivable.  \rulename{Refl},
\rulename{Eval}, \rulename{Use} and \rulename{Elim} are treated above.  An
introduction by progress at the literal $0$ is the bit $0$.  For an
introduction by cases at a term $t$, open $T(\Suc\rho t)$ at a fresh $b$,
retain $\mathsf{a}(b=\rho t)$ for the moment, and introduce the formula with
the bit $1$.  In the premise the Leibniz rule turns the assumption
$B(b,0)$ into $B(\rho t,0)$.  The family given by the induction hypothesis
is taken at the map that keeps the images of the source parameters, sends
the source eigenvariable to the new one and skips $b$; the cut of an
assumption supplies its assumption, and the Leibniz rule gives
$B(b,\Suc z)$.  After the introduction, the cut of an assumption replaces
$\mathsf{a}(b=\rho t)$ by its hereditary derivation from the opened packet,
and \rulename{Elim} discharges $b$.
\end{proof}

\begin{corollary}
\label{wit:cor:clause-bridge}
The ordinary realizations of $\mathsf{R}_{\nu}$, $\mathsf{R}_{\beta}$,
$\Nat_{\mathrm{diag},\nu}$ and $\mathsf{C}_{\beta}$ are bi-interpretable and
equiconsistent with $\iPRA$, with whole numerical domains, actual equality
and the identity graph in both numerical composites.
\end{corollary}

\begin{proof}
Each calculus derives the same judgments, on the same terms and formulas, as
a calculus for which the bridge of \cref{base:thm:bridge} is established:
$\Nat$ for the first three, and the unrestricted clause calculus
$\mathsf{C}$ for the last, which derives the same sequents as $\Nat$ under
the translation of \cref{claus:thm:five}.  The translations and the
derivations of both formula round trips transfer unchanged.
\end{proof}

\subsection{The newest parameter in the Alias/Back calculus}
\label{wit:sec:aliasback}

This subsection draws on \pkg{local\_introduction}; its equality interface
is the Alias/Back basis of the chapter on discharge, from
\pkg{alias\_back}.

\Cref{wit:thm:diagonal} removes the witness from the diagonal introduction
in a calculus with the primitive transport \rulename{Tr}, and it recovers
an arbitrary witness $t$ from the packet $X(\Suc t)$, whose opening supplies
a parameter $b$ with $b=t$.  The Alias/Back basis has no primitive
transport: equality is given by a fresh alias, discharged by substitution,
and by a backward substitution at the newest parameter.  There the alias
itself supplies the parameter $b$ with $b=t$, and the witness is removed
without the packet, without names, without a predecessor and without
normalization.  The mode of the introduction is the diagonal mode of
\cref{wit:sec:diagonal}.  The proof below has the shape of the proof of
\cref{wit:thm:newest} and differs from it in three respects: equality on
families is built from \rulename{Alias} and \rulename{Back}
(\cref{wit:lem:ab-eq}); the source clauses are given by realizers and not by
literal membership, because the transported assumption of an introduction
is not the image of the source assumption (this replaces the cut of an
assumption in the proof of \cref{wit:thm:clause-bit}); and one induction
removes, together with the witness, the arbitrary terms of \rulename{Use}
and the arbitrary sources of \rulename{Back}.

\begin{definition}[The calculi $\mathsf{A}_{\nu}$ and $\mathsf{A}_{\ast}$]
\label{wit:def:ab}
Terms are those of \cref{base:sec:terms}; formulas are equations and
$\D{x}{y}B$, without $\Q{}$.  Substitution is capture avoiding and
simultaneous, inserts each replacement once and performs no computation.
A context is a list of clauses $[A;\,y.B(y)]$ and $[\ ;\,y.B(y)]$, read as
$A\to\forall y\,B(y)$ and $\forall y\,B(y)$; an assumption $F$ is the clause
$[\ ;\,v.F]$ with $v$ not in $F$.  Judgments are $\Pi;\Gamma\proves A$ as in
\cref{wit:def:scope}, with the clauses of $\Gamma$ over $\Pi$.  Put
\[
 \epsilon_B(z)\defeq
 \begin{cases}
  z,&\text{$B$ is diagonal in $y$, that is, literally
        $\It{y}{y}{u}{h}=\It{y}{y}{v}{k}$,}\\
  0,&\text{otherwise,}
 \end{cases}
\]
the diagonal test of \cref{guard:def:shapes} and \cref{claus:def:k}, made on
the matrix $B(x,y)$ before a witness is substituted.  The calculus
$\mathsf{A}_{\nu}$ has exactly the following seven rule patterns.
\begin{enumerate}[label=\textup{(\roman*)}]
\item \rulename{I0}: $\Pi;\Gamma\proves\It{0}{a}{u}{h}=a$.
\item \rulename{IS}: $\Pi;\Gamma\proves\It{\Suc n}{a}{u}{h}
 =h\subst{\It{n}{a}{u}{h}}{u}$.
\item \rulename{Use}: for a clause $[A;\,y.B]$ of $\Gamma$, from
 $\Pi;\Gamma\proves A$ infer $\Pi;\Gamma\proves B(\nu(\Pi))$; for a clause
 $[\ ;\,y.B]$ of $\Gamma$, infer $\Pi;\Gamma\proves B(\nu(\Pi))$ outright.
\item \rulename{Back}: if $\Pi$ is not empty and $b=\nu(\Pi)$, from
 $\Pi;\Gamma\proves b=t$ and $\Pi;\Gamma\proves F\subst{t}{b}$ infer
 $\Pi;\Gamma\proves F$.  The term $t$ may contain $b$.
\item \rulename{Alias}: if $b\notin\Pi$ and $t$ is over $\Pi$, from
 $\Pi,b;\ \Gamma,b=t\proves F$ infer $\Pi;\Gamma\proves F\subst{t}{b}$.  The
 formula $F$ may contain $b$.
\item \rulename{Open}: if $b\notin\Pi$ and $A$ is over $\Pi$, from
 $\Pi;\Gamma\proves\D{x}{y}B$ and
 $\Pi,b;\ \Gamma,[B(b,0);\,y.B(b,y)]\proves A$ infer $\Pi;\Gamma\proves A$.
\item \rulename{Close}: if $z\notin\Pi$, from
 $\Pi,z;\ \Gamma,B(\nu(\Pi),\epsilon_B(z))\proves B(\nu(\Pi),\Suc z)$ infer
 $\Pi;\Gamma\proves\D{x}{y}B$.
\end{enumerate}
The calculus $\mathsf{A}_{\ast}$ has the same syntax, \rulename{I0},
\rulename{IS}, \rulename{Alias} and \rulename{Open}, a \rulename{Use} at an
arbitrary term $r$ over $\Pi$ in place of $\nu(\Pi)$, a \rulename{Back} at
an arbitrary parameter $b$ of $\Pi$, and a \rulename{Close} at an arbitrary
witness $t$ over $\Pi$ in place of $\nu(\Pi)$, with the same mode
$\epsilon_B$.
\end{definition}

\rulename{Close} of $\mathsf{A}_{\nu}$ takes no term.  There is no
reflexivity rule, transport rule, normalizer, primitive specialization at
a chosen term or logical connective.  The rules other than \rulename{Close}
are those of the Alias/Back basis, whose \rulename{Use} and \rulename{Back}
are already at the newest parameter; its introduction takes an arbitrary
witness and is by progress exactly when the matrix is a literal equation
and the witness is literally $0$.  $\mathsf{A}_{\nu}$ changes only the
introduction: the witness is the newest parameter, and the mode is read off
the matrix alone.  Every rule of $\mathsf{A}_{\nu}$ is a rule of
$\mathsf{A}_{\ast}$.  Every rule that opens a branch appends one parameter
and one clause, so along a derivation of a closed judgment the length of
the scope equals the length of the context.

\begin{lemma}[Soundness]
\label{wit:lem:ab-sound}
Every rule of $\mathsf{A}_{\ast}$ is valid in $\Teq$ under the ordinary
reading of \cref{base:sec:reading} and of the clauses.
\end{lemma}

\begin{proof}
The computations are axioms of $\Teq$ (\cref{base:sec:teq}).
\rulename{Use} is the instantiation of the universal quantifier of the
clause, after modus ponens for a clause with an antecedent.
\rulename{Back} is the substitution of equals.  For \rulename{Alias}, the
ordinary derivation of the branch is closed under the substitution of $t$
for the variable $b$, which occurs neither in $\Gamma$ nor in $t$; the
assumption becomes $t=t$, which reflexivity proves.  \rulename{Open} is the
elimination of the existential quantifier of
$\exists x\,(B(x,0)\to\forall y\,B(x,y))$ at the fresh $b$.  For
\rulename{Close} at the witness $t$: if $B$ is diagonal it is a literal
equation, its premise is a progressive step, and induction for literal
equations proves $B(t,0)\to\forall y\,B(t,y)$; otherwise the premise is a
case step, and coverage proves the same implication.  The existential
quantifier is introduced at $t$.
\end{proof}

\begin{definition}[Persistent families and realizers]
\label{wit:def:ab-family}
A \emph{world} is a pair $w=(\Pi_w;\Gamma_w)$ of a scope and a context over
it.  A world $w'$ \emph{extends} $w$, written $w'\succeq w$, if $\Pi_{w'}$
is $\Pi_w$ followed by further parameters and $\Gamma_{w'}$ is $\Gamma_w$
followed by further clauses over $\Pi_{w'}$.  For a formula $A$ over
$\Pi_w$, write $w\Vdash A$ if there is a \emph{persistent family} for $A$
at $w$: an effective procedure that assigns to every $w'\succeq w$ a
derivation of $\mathsf{A}_{\nu}$ of $\Pi_{w'};\Gamma_{w'}\proves A$.  A
\emph{persistent realizer} at $w$ of a clause $[A;\,y.B]$ over $\Pi_w$ is
an effective procedure that assigns to every $w'\succeq w$, every term $r$
over $\Pi_{w'}$ and every persistent family for $A$ at $w'$ a persistent
family for $B(r)$ at $w'$; for a clause $[\ ;\,y.B]$ the family argument is
absent.
\end{definition}

A family or realizer at $w$ is one at every $w'\succeq w$, by the same
procedure.  A family for $E$ at $w$ is a realizer at $w$ of the assumption
$[\ ;\,v.E]$: it ignores the term $r$, since $E\subst{r}{v}\equiv E$.
Families are the hereditary derivability of \cref{wit:def:hereditary}
restricted to extensions; the injective map is carried by the induction
statement of \cref{wit:thm:ab-newest} instead.  They and the realizers are
constructions of the metatheory and are not rules of the calculus.

\begin{lemma}[Clauses and equality on families]
\label{wit:lem:ab-eq}
Let $w$ be a world, $s$, $t$ terms over $\Pi_w$ and $H(v)$ a formula with a
distinguished variable $v$, which may occur any number of times and below
binders.
\begin{enumerate}[label=\textup{(\roman*)}]
\item Every clause of $\Gamma_w$ has a persistent realizer at $w$.
\item $w\Vdash t=t$.
\item If $w\Vdash s=t$ and $w\Vdash H(t)$, then $w\Vdash H(s)$.
\item If $w\Vdash s=t$, then $w\Vdash t=s$; if moreover $w\Vdash H(s)$,
 then $w\Vdash H(t)$.
\end{enumerate}
\end{lemma}

\begin{proof}
In each item let $w'\succeq w$ be given, and let $c$ be a parameter not in
$\Pi_{w'}$ and not in $s$, $t$, $H$ or the terms named below.

(i) Let $[A;\,y.B]$ be a clause of $\Gamma_w$, $w'\succeq w$, $r$ a term
over $\Pi_{w'}$, and a family for $A$ at $w'$ be given; let $w''\succeq w'$.
The world $w''_c=(\Pi_{w''},c;\ \Gamma_{w''},c=r)$ extends $w'$, so the
family gives $A$ there.  The clause is a member of $\Gamma_{w''_c}$ and $c$
is the newest parameter, so \rulename{Use} gives $B(c)$, and
\rulename{Alias} gives $B(c)\subst{r}{c}\equiv B(r)$ at $w''$.  For a clause
without antecedent the family is not needed.

(ii) In $(\Pi_{w'},c;\ \Gamma_{w'},c=t)$, \rulename{Use} of the last clause
gives $c=t$, and \rulename{Alias} gives $(c=t)\subst{t}{c}\equiv(t=t)$ at
$w'$.

(iii) Work in $w'_c=(\Pi_{w'},c;\ \Gamma_{w'},c=s)$, which extends $w$.
\rulename{Use} of the last clause gives $c=s$, and the first family gives
$s=t\equiv(c=t)\subst{s}{c}$; \rulename{Back} at $c$ gives $c=t$.  The
second family gives $H(t)\equiv H(c)\subst{t}{c}$, and \rulename{Back} at $c$
with the major premise $c=t$ gives $H(c)$.  \rulename{Alias} gives
$H(c)\subst{s}{c}\equiv H(s)$ at $w'$.

(iv) Apply (iii) to the template $t=v$, whose instance $t=t$ holds by
(ii); the result is $t=s$.  For the second statement apply (iii) to $t=s$
and $H(s)$.

Both applications of \rulename{Back} in (iii) are at $c$, which is the
newest parameter of $w'_c$; the derivations of their premises are the
members of the given families at $w'_c$ itself, and no derivation is moved
to another world.  These constructions use no computation, induction,
coverage, weakening of a completed derivation or substitution into one.
\end{proof}

\begin{theorem}[No witness in the Alias/Back calculus]
\label{wit:thm:ab-newest}
For every scope $\Pi$, clause context $\Gamma$ and formula $A$ over $\Pi$,
\[
 \Pi;\Gamma\proves_{\mathsf{A}_{\nu}}A
 \quad\Longleftrightarrow\quad
 \Pi;\Gamma\proves_{\mathsf{A}_{\ast}}A ,
\]
by effective transformations of derivations that leave the terms, formulas
and clauses of the judgment unchanged.
\end{theorem}

\begin{proof}
From left to right every rule of $\mathsf{A}_{\nu}$ is a rule of
$\mathsf{A}_{\ast}$.  For the converse the following \emph{induction
statement} is proved by induction on a derivation of $\mathsf{A}_{\ast}$ of
$\Pi;\Gamma\proves A$: for every world $w$, every injective map $\theta$
from $\Pi$ into $\Pi_w$ and every assignment of a persistent realizer at
$w$ of $\theta C$ to each clause $C$ of $\Gamma$, $w\Vdash\theta A$.
Here $\theta$ acts on free parameters; since it is injective,
$\theta(F\subst{t}{b})\equiv(\theta F)\subst{\theta t}{\theta b}$, and
$\epsilon_{\theta B}=\epsilon_B$, because the diagonal shape concerns the
bound variable $y$ and the literal identity of two iterator terms, both
preserved and reflected by an injective renaming.  The required derivation
of $\Pi;\Gamma\proves A$ is the member at $w=(\Pi;\Gamma)$ of the family
that the induction statement gives for the identity map and the realizers of
\cref{wit:lem:ab-eq}(i).  In the cases below, $w'\succeq w$ is the world at
which a member is constructed.

\emph{Computations.}  The image of an instance is an instance, in every
world.

\emph{Use} of a clause $C=[A;\,y.B]$ at the term $r$.  The induction
hypothesis gives a family for $\theta A$ at $w$; the realizer of $\theta C$
at the term $\theta r$ turns it into a family for
$\theta B(\theta r)\equiv\theta(B(r))$.  For a clause without antecedent no
family is needed.

\emph{Back} at a parameter $b$ of $\Pi$, with premises $b=t$ and
$F\subst{t}{b}$.  Let $H(v)\defeq(\theta F)\subst{v}{\theta b}$ with $v$
fresh; then $H(\theta t)\equiv\theta(F\subst{t}{b})$ and
$H(\theta b)\equiv\theta F$.  The induction hypothesis gives families for
$\theta b=\theta t$ and for $H(\theta t)$, and \cref{wit:lem:ab-eq}(iii)
gives $w\Vdash\theta F$.

\emph{Alias} with the parameter $b$, the term $t$ and the branch
$\Pi,b;\Gamma,b=t\proves F$.  Choose $b'\notin\Pi_{w'}$ and let $w''$ be
$(\Pi_{w'},b';\ \Gamma_{w'},b'=\theta t)$.  The map
$\theta'=\theta[b\mapsto b']$ is injective; the clauses of $\Gamma$ keep
their realizers, restricted to $w''$, and the new clause $b=t$, whose
image is $b'=\theta t$, is realized by \cref{wit:lem:ab-eq}(i).  The
induction hypothesis gives $\theta'F$ at $w''$, and \rulename{Alias} gives
$(\theta'F)\subst{\theta t}{b'}\equiv\theta(F\subst{t}{b})$ at $w'$.

\emph{Open} with the major premise $\D{x}{y}B$ and the parameter $b$.  The
induction hypothesis gives $\theta\D{x}{y}B$ at $w'$.  Choose
$b'\notin\Pi_{w'}$ and extend $\theta$ by $b\mapsto b'$; the opened clause
is realized in
$w''=(\Pi_{w'},b';\ \Gamma_{w'},[\theta B(b',0);\,y.\theta B(b',y)])$ by
\cref{wit:lem:ab-eq}(i).  The induction hypothesis gives $\theta A$ at
$w''$, since $A$ does not contain $b$, and \rulename{Open} applies.

\emph{Close} at the witness $t$ with the matrix $B$ and the premise
$\Pi,z;\ \Gamma,B(t,\epsilon_B(z))\proves B(t,\Suc z)$.  Choose
$b\notin\Pi_{w'}$ and apply \rulename{Alias} with the defining term
$\theta t$; its branch is the world
$w_b=(\Pi_{w'},b;\ \Gamma_{w'},b=\theta t)$.  In $w_b$ apply
\rulename{Close} to $\D{x}{y}\theta B$.  The newest parameter is $b$, so the
rule selects $b$; let $z'$ be its eigenvariable.  Its premise is to be
derived in
\[
 w_z=\bigl(\Pi_{w'},b,z';\ \Gamma_{w'},\ b=\theta t,\
 \theta B(b,\epsilon_B(z'))\bigr).
\]
Put $\theta'=\theta[z\mapsto z']$, an injective map from $\Pi,z$ into
$\Pi_{w_z}$; the parameter $b$ is not in its image.  The clauses of
$\Gamma$ keep their realizers, restricted to $w_z$.  The source assumption
$B(t,\epsilon_B(z))$ has the image $\theta B(\theta t,\epsilon_B(z'))$,
which is not a clause of $w_z$.  By \cref{wit:lem:ab-eq}(i) for the last
two clauses of $w_z$, which are assumptions,
$w_z\Vdash b=\theta t$ and $w_z\Vdash\theta B(b,\epsilon_B(z'))$, and
\cref{wit:lem:ab-eq}(iv) with the template $\theta B(v,\epsilon_B(z'))$
gives $w_z\Vdash\theta B(\theta t,\epsilon_B(z'))$; this family is the
realizer of the source assumption.  The induction hypothesis for the
premise gives $w_z\Vdash\theta B(\theta t,\Suc z')$, and
\cref{wit:lem:ab-eq}(iii) with the template $\theta B(v,\Suc z')$ and the
equation $b=\theta t$ gives $w_z\Vdash\theta B(b,\Suc z')$.  Its member at
$w_z$ is the premise of the target \rulename{Close}, whose mode is
$\epsilon_{\theta B}=\epsilon_B$.  The conclusion $\D{x}{y}\theta B$ does not
contain $b$, so \rulename{Alias} yields
$(\D{x}{y}\theta B)\subst{\theta t}{b}\equiv\theta\D{x}{y}B$ at $w'$.

The eigenvariable $z'$ is newer than $b$; this does no harm, because every
transport of \cref{wit:lem:ab-eq} opens its own auxiliary parameter, which
is then the newest.  The construction is well founded.  The realizer of the
source assumption in the case of \rulename{Close} is built from the actual
clauses of $w_z$ and does not call the translation of the premise.  Every
call of the induction hypothesis concerns a proper subderivation, although
it may be made in several worlds; \rulename{Use} calls realizers on
families of proper subderivations; the constructions of
\cref{wit:lem:ab-eq} call the given families finitely often.  No derivation
of $\mathsf{A}_{\nu}$ is moved to another world, and no term is substituted
into one.
\end{proof}

\begin{corollary}[The Alias/Back calculi and $\iPRA$]
\label{wit:cor:ab-bridge}
For a clause context $\Gamma$ let $\Gamma^{\circ}$ be $\Gamma$ with every
clause $[\ ;\,y.B]$ replaced by $[0=0;\,y.B]$.  Then
\[
 \Pi;\Gamma\proves_{\mathsf{A}_{\nu}}A
 \iff \Pi;\Gamma\proves_{\mathsf{A}_{\ast}}A
 \iff \Gamma^{\circ}\proves_{\mathsf{K}}A ,
\]
where $\mathsf{K}$ is the diagonal clause calculus of \cref{claus:def:k}.
The ordinary realizations of $\mathsf{A}_{\nu}$ and $\mathsf{A}_{\ast}$ are
bi-interpretable and equiconsistent with $\iPRA$, with whole numerical
domains, actual equality and the identity graph in both numerical
composites.
\end{corollary}

\begin{proof}
The first equivalence is \cref{wit:thm:ab-newest}.  In $\mathsf{A}_{\ast}$,
injective renaming and weakening by appended parameters and clauses are
admissible, by induction on derivations after the eigenvariables are made
fresh: no rule of $\mathsf{A}_{\ast}$ refers to the order of the scope, the
source of a \rulename{Back} is renamed to a parameter, and $\epsilon_B$ is
preserved.  A derivation of $\mathsf{A}_{\ast}$ therefore gives a
derivation in every extension, and the constructions of
\cref{wit:lem:ab-eq}, which are derivations of $\mathsf{A}_{\ast}$, give
reflexivity, symmetry, transitivity and transport in every formula over
every scope and context.  In $\mathsf{A}_{\ast}$ the clauses $[\ ;\,y.B]$
and $[0=0;\,y.B]$ are interchangeable: a \rulename{Use} of the second drops
its premise when it becomes a \rulename{Use} of the first, and a
\rulename{Use} of the first becomes a \rulename{Use} of the second with the
premise $0=0$ by reflexivity.

\emph{From $\mathsf{K}$ to $\mathsf{A}_{\ast}$.}  By induction on a
derivation of $\Gamma^{\circ}\proves_{\mathsf{K}}A$ in the scope $\Pi$.
First its eigenvariables are made fresh for $\Pi$
(\cref{claus:lem:struct}, which holds in $\mathsf{K}$), and every other
variable that occurs in it and is not in $\Pi$ is replaced by $0$
(\cref{claus:thm:ksubst}); the judgment is unchanged, and every term of the
derivation is then over the scope of its node.
\rulename{Refl} is reflexivity.  The zero row of \rulename{Eval}, from
$a=t$ and $\It{0}{a}{u}{h}=r$, gives $r=t$ by \rulename{I0}, symmetry and
transitivity; the successor row, from $\It{n}{a}{u}{h}=t$ and
$\It{\Suc n}{a}{u}{h}=r$, gives $r=h\subst{t}{u}$ by \rulename{IS},
transport in the template $\It{\Suc n}{a}{u}{h}=h\subst{v}{u}$, symmetry and
transitivity.  \rulename{Use}, \rulename{Elim} and \rulename{Intro} are
\rulename{Use}, \rulename{Open} and \rulename{Close} at the same term, the
same parameter and the same witness, with the same mode, since
$\epsilon_B$ is the mode $\eta_B$ of \cref{claus:def:k}.  The result is a
derivation of $\Pi;\Gamma^{\circ}\proves A$, hence of $\Pi;\Gamma\proves A$.

\emph{From $\mathsf{A}_{\ast}$ to $\mathsf{K}$.}  By induction on a
derivation of $\Pi;\Gamma\proves A$, with the context $\Gamma^{\circ}$.  The
computation equations, symmetry and transport in every formula hold in
$\mathsf{K}$ (\cref{claus:lem:eq} and \cref{claus:lem:k}(iii)).
\rulename{Use} is \rulename{Use}, with the premise $0=0$ by \rulename{Refl}
for a clause without antecedent.  \rulename{Back} at $b$ is transport in
the template $F\subst{v}{b}$ along $t=b$.  For \rulename{Alias}, the branch
gives $\Gamma^{\circ},[0=0;\,v.\,b=t]\proves_{\mathsf{K}}F$; substitution of
$t$ for $b$ (\cref{claus:thm:ksubst}) gives
$\Gamma^{\circ},[0=0;\,v.\,t=t]\proves_{\mathsf{K}}F\subst{t}{b}$, since $b$
is not in $\Gamma$, and the cut of the ordinary assumption against
\rulename{Refl} (\cref{claus:lem:struct}) removes the clause.
\rulename{Open} is \rulename{Elim} and \rulename{Close} is \rulename{Intro}
at the same witness with the same mode; the side conditions of $\mathsf{K}$
on fresh variables follow from those on scopes.

By \cref{claus:thm:k}, $\mathsf{K}$ derives, for $Q$-free formulas and the
translation $\tau$ of native formulas, the same judgments as $\Nat$
(\Lean{kprf\_iff\_nprf} and \Lean{kprf\_iff\_nprf\_of\_qfree},
\cref{claus:rem:kpkgs}).  The bridge of \cref{base:thm:bridge} applies as in
\cref{wit:cor:clause-bridge}: the comparisons are the identity on terms and
on $Q$-free formulas, $\Q{y}B$ is represented by $\D{x}{y}B$ with $x$ unused,
and the translations, both formula round trips and the preservation of
$0=\Suc 0$ transfer unchanged.
\end{proof}

\begin{remark}[What the proof uses of the mode]
\label{wit:rem:ab-mode}
The proof of \cref{wit:thm:ab-newest} uses the mode only in the case of
\rulename{Close}, and there only through $\epsilon_{\theta B}=\epsilon_B$
for injective renamings $\theta$.  The same proof therefore removes the
witness from every introduction of the form
$\Gamma,B(t,\epsilon_B(z))\proves B(t,\Suc z)$ over $\Gamma\proves\D{x}{y}B$,
with $\epsilon_B(z)\in\{0,z\}$ determined by the uninstantiated matrix and
preserved and reflected by injective renaming, in a calculus with the
same \rulename{Alias}, \rulename{Use} and \rulename{Back}.  It is a theorem
of relative derivability; soundness and arithmetic adequacy of the selected
mode are separate obligations.  The mirror guard of
\cref{claus:prop:mirror} qualifies: an injective renaming preserves and
reflects the shape $a=\It{y}{a}{u}{h}$, although an arbitrary substitution
can create it, and the proof uses no invariance under substitution.

The mode of the introduction of the Alias/Back basis does not qualify: it
tests whether the witness is literally $0$, and the alias $b$ is a
parameter.  A source introduction at $0$ in the progress mode would become
a target introduction at $b$ in the case mode, whose assumption
$\theta B(b,0)$ does not transport to the source assumption
$\theta B(0,z')$, and the proof does not apply.  Replacing the witness by
the newest parameter while keeping that mode is not justified by it; the
one-bit rule of \cref{wit:sec:bit} is a different repair.  The proof uses
\rulename{Back} on compound formulas, in the two transports of the case of
\rulename{Close}, when $B$ has a binder; \cref{wit:thm:eqback} removes this
use.
\end{remark}

\begin{remark}[The package]
\label{wit:rem:ab-package}
\pkg{local\_introduction} states the theorem for the temporary calculus
$\mathsf{A}_{\ast}$ with persistent families over worlds and realizers, as
above, and the corollary as the equality of the open consequences of the
two diagonal alias calculi together with the inherited bridge; the
comparison with $\mathsf{K}$ is made explicit by the adapters for
\rulename{Refl}, the two rows of \rulename{Eval} and the clauses without
antecedent, which are those of the proof above.  The package includes an
independent reading of the argument, which records the six assumptions
the proof needs: \rulename{Back} on compound formulas, the substitutive
export of \rulename{Alias}, the mode tested on the raw matrix, the
retention of old parameters in an extension, effective realizers quantified
over all extensions, and capture-avoiding simultaneous substitution.  The
theorem is not formalized in Lean or Rocq.  The first assumption is removed
by \cref{wit:thm:eqback}; the second is reduced to the export of equations
and one naming scheme by \cref{wit:thm:ns-complete}.
\end{remark}

\subsection{Backward substitution in equations only}
\label{wit:sec:eqback}

This subsection draws on \pkg{equational\_transport}.

The proof of \cref{wit:thm:ab-newest} obtains transport in every formula
from \rulename{Back} on compound formulas (\cref{wit:lem:ab-eq}(iii)).  If
\rulename{Back} acts on equations only, formula transport has to be
recovered first, and two shortcuts are unavailable.  A completed derivation
cannot be moved to an extension of its world, since its next \rulename{Use}
or \rulename{Close} would select a different newest parameter; and the
transport needed in the case of \rulename{Close} cannot be obtained by
translating a derivation of $\mathsf{A}_{\ast}$ that uses compound
\rulename{Back}, which would be circular.  The construction below proceeds
in the order: realizers of clauses, transport in equations, zero rows of
diagonal matrices, transport in every formula by induction on the number of
binders, and only then the translation of derivations.

\begin{definition}[The calculi $\mathsf{A}^{=}_{\nu}$ and
$\mathsf{A}^{\times}_{\nu}$]
\label{wit:def:eqback}
$\mathsf{A}^{=}_{\nu}$ and $\mathsf{A}^{\times}_{\nu}$ have the syntax,
judgments and mode $\epsilon_B$ of \cref{wit:def:ab}, and the rules
\rulename{I0}, \rulename{IS}, \rulename{Use}, \rulename{Alias},
\rulename{Open} and \rulename{Close} of $\mathsf{A}_{\nu}$, with
\rulename{Use} and \rulename{Close} at the newest parameter.  In place of
\rulename{Back}, $\mathsf{A}^{=}_{\nu}$ has
\begin{itemize}
\item \rulename{EqBack}: if $\Pi$ is not empty, $b=\nu(\Pi)$ and $E$ is an
 equation, from $\Pi;\Gamma\proves b=t$ and $\Pi;\Gamma\proves E\subst{t}{b}$
 infer $\Pi;\Gamma\proves E$.  The term $t$ may contain $b$;
\end{itemize}
and $\mathsf{A}^{\times}_{\nu}$ has
\begin{itemize}
\item \rulename{XBack}: if $\Pi$ is not empty, $b=\nu(\Pi)$ and $h$, $r$, $t$
 are terms over $\Pi$,
 \[
  \infrule{\Pi;\Gamma\proves b=t\qquad \Pi;\Gamma\proves h\subst{t}{b}=r}
          {\Pi;\Gamma\proves r=h}.
 \]
 Only $h$ is substituted; $r$ and $t$ may contain $b$.
\end{itemize}
Each calculus has seven rule patterns.
\end{definition}

\rulename{EqBack} is the instance of \rulename{Back} of $\mathsf{A}_{\nu}$
whose conclusion is an equation, so $\mathsf{A}^{=}_{\nu}$ is a subcalculus of
$\mathsf{A}_{\nu}$.  \rulename{XBack} has a term template in place of a
formula template; it is \rulename{Tr} of \cref{base:sec:kernel} with the
major premise reversed, the source restricted to the newest parameter and
the template $h\subst{q}{b}$.  It is not an instance of \rulename{Back} when
$r$ contains $b$, and it is compared with $\mathsf{A}_{\ast}$ through derived
rules below.  Both restrictions concern backward substitution only:
\rulename{Alias} still exports arbitrary formulas through substitution.
\Cref{wit:sec:ns} restricts \rulename{Alias} as well.

For $\mathsf{C}$ one of $\mathsf{A}^{=}_{\nu}$ and $\mathsf{A}^{\times}_{\nu}$,
\emph{persistent families} and \emph{persistent realizers} of $\mathsf{C}$,
and $w\Vdash_{\mathsf{C}}A$, are defined as in \cref{wit:def:ab-family} with
derivations of $\mathsf{C}$ in place of derivations of $\mathsf{A}_{\nu}$;
the remarks after that definition hold unchanged.  The index $\mathsf{C}$ is
omitted when it is fixed.  In this subsection $q$ denotes a distinguished
template variable, distinct from every parameter and bound variable named.

\begin{lemma}[Equality in equations on families]
\label{wit:lem:eqb-atomic}
Let $\mathsf{C}$ be $\mathsf{A}^{=}_{\nu}$ or $\mathsf{A}^{\times}_{\nu}$, $w$ a
world, $s$, $t$ terms over $\Pi_w$, and $E(q)$ an equation over $\Pi_w$ and
$q$, in which $q$ may occur any number of times.
\begin{enumerate}[label=\textup{(\roman*)}]
\item Every clause of $\Gamma_w$ has a persistent realizer at $w$.
\item $w\Vdash t=t$.
\item If $w\Vdash s=t$ and $w\Vdash E(t)$, then $w\Vdash E(s)$.
\item If $w\Vdash s=t$, then $w\Vdash t=s$; if moreover $w\Vdash E(s)$, then
 $w\Vdash E(t)$.
\end{enumerate}
\end{lemma}

\begin{proof}
(i) and (ii) are proved as \cref{wit:lem:ab-eq}(i) and (ii), whose
constructions use \rulename{Use} at the newest parameter and
\rulename{Alias} only.

For $\mathsf{A}^{=}_{\nu}$, (iii) and (iv) are proved as
\cref{wit:lem:ab-eq}(iii) and (iv) with $H\defeq E$: the two applications of
\rulename{Back} in (iii) are at the newest parameter $c$ of $w'_c$ and have
the conclusions $c=t$ and $E(c)$, which are equations, so they are
applications of \rulename{EqBack}; (iv) applies (iii) to the equation $t=q$.

For $\mathsf{A}^{\times}_{\nu}$, write $E(q)$ as $a(q)=d(q)$ with terms
$a(q)$, $d(q)$.  For (iii) let $w'\succeq w$, choose a parameter $b$ not in
$\Pi_{w'}$, and work in $w'_b=(\Pi_{w'},b;\ \Gamma_{w'},b=s)$, which extends
$w$ and has the newest parameter $b$.  Since $b$ occurs in none of $s$, $t$,
$a(q)$, $d(q)$, the substitution identities $b\subst{s}{b}\equiv s$,
$t\subst{s}{b}\equiv t$, $a(b)\subst{t}{b}\equiv a(t)$ and
$d(b)\subst{t}{b}\equiv d(t)$ hold literally.  The following derivations are
made in $w'_b$:
\begin{enumerate}[label=\textup{(\arabic*)}]
\item $b=s$, by \rulename{Use} of the last clause;
\item $t=b$, by \rulename{XBack} with the major premise (1), $h\defeq b$, and
 the minor premise $s=t$, the member at $w'_b$ of the first family;
\item $b=t$, by \rulename{XBack} with the major premise (1), $h\defeq t$, and
 the minor premise (2); here $r\equiv b$ contains $b$;
\item $d(t)=a(b)$, by \rulename{XBack} with the major premise (3),
 $h\defeq a(b)$, and the minor premise $a(t)=d(t)$, the member at $w'_b$ of
 the second family;
\item $a(b)=d(b)$, by \rulename{XBack} with the major premise (3),
 $h\defeq d(b)$, and the minor premise (4); here $r\equiv a(b)$ contains $b$
 and is not substituted.
\end{enumerate}
\rulename{Alias} gives $(a(b)=d(b))\subst{s}{b}\equiv(a(s)=d(s))$ at $w'$.
For the first statement of (iv), (1), (2) and \rulename{Alias} give
$(t=b)\subst{s}{b}\equiv(t=s)$ at $w'$; the second statement is (iii)
applied to $t=s$ and $E(s)$.

In both calculi every backward substitution is at the newest parameter of
the auxiliary world, its premises are members of the given families at that
world itself, and no derivation is moved to another world or receives a
substitution.  No computation, induction or \rulename{Close} is used.
\end{proof}

\begin{lemma}[Diagonal zero rows]
\label{wit:lem:eqb-zero}
Let $\mathsf{C}$ be as in \cref{wit:lem:eqb-atomic}, and let $B(x,y)$ be
diagonal in $y$, that is, literally $\It{y}{y}{u}{h}=\It{y}{y}{v}{k}$.  For
every world $w$ and every term $r$ over $\Pi_w$ such that $B(r,0)$ is over
$\Pi_w$, $w\Vdash B(r,0)$.
\end{lemma}

\begin{proof}
$B(r,0)$ is $\It{0}{0}{u}{h'}=\It{0}{0}{v}{k'}$, where $h'$ and $k'$ are $h$
and $k$ with $r$ for $x$ and $0$ for $y$.  The instances
$\It{0}{0}{u}{h'}=0$ and $\It{0}{0}{v}{k'}=0$ of \rulename{I0} are derivable in
every world, so they have persistent families.  \Cref{wit:lem:eqb-atomic}(iv)
gives $w\Vdash 0=\It{0}{0}{v}{k'}$, and \cref{wit:lem:eqb-atomic}(iii) with
$s\defeq\It{0}{0}{u}{h'}$, $t\defeq 0$ and the equation
$q=\It{0}{0}{v}{k'}$ gives $w\Vdash B(r,0)$.  Neither formula transport nor
\rulename{Close} is used.
\end{proof}

\begin{lemma}[Transport in every formula]
\label{wit:lem:eqb-transport}
Let $\mathsf{C}$ be as in \cref{wit:lem:eqb-atomic}.  For every world $w$,
all terms $s$, $t$ over $\Pi_w$ and every formula $H(q)$ over $\Pi_w$ and $q$,
in which $q$ may occur any number of times and below binders:
\begin{enumerate}[label=\textup{(\roman*)}]
\item if $w\Vdash s=t$ and $w\Vdash H(t)$, then $w\Vdash H(s)$;
\item if $w\Vdash s=t$ and $w\Vdash H(s)$, then $w\Vdash H(t)$.
\end{enumerate}
\end{lemma}

\begin{proof}
(ii) follows from (i) and the symmetry of \cref{wit:lem:eqb-atomic}(iv).
Statement (i) is proved, simultaneously for all worlds and terms, by
induction on the number $d(H)$ of binders $\D{}{}$ in $H$; substituting a
term for $q$ does not change this number.  If $d(H)=0$, $H$ is an equation
and (i) is \cref{wit:lem:eqb-atomic}(iii).

Let $H(q)=\D{x}{y}B(q,x,y)$, so that $d(B)=d(H)-1$, and let $w'\succeq w$.
The member at $w'$ of the family for $H(t)$ is a derivation of
$\D{x}{y}B(t,x,y)$.  Choose $b\notin\Pi_{w'}$ and apply \rulename{Open} with
this major premise; its branch is to be derived in
\[
 w_b=\bigl(\Pi_{w'},b;\ \Gamma_{w'},[B(t,b,0);\,y.B(t,b,y)]\bigr)
\]
with the goal $H(s)$, which is over $\Pi_{w'}$.  Put $C(x,y)\defeq B(s,x,y)$,
so that $H(s)\equiv\D{x}{y}C$.  In $w_b$ apply \rulename{Close} to
$\D{x}{y}C$; the newest parameter of $w_b$ is $b$, so the rule selects $b$.
With its eigenvariable $z\notin\Pi_{w_b}$ the premise is $C(b,\Suc z)\equiv
B(s,b,\Suc z)$, to be derived in
\[
 w_z=\bigl(\Pi_{w'},b,z;\ \Gamma_{w'},[B(t,b,0);\,y.B(t,b,y)],\
 C(b,\epsilon_C(z))\bigr).
\]
The world $w_z$ extends $w$, so $w_z\Vdash s=t$ by restriction.
\begin{enumerate}[label=\textup{(\alph*)}]
\item $w_z\Vdash B(s,b,0)$.  If $\epsilon_C(z)=0$, the last clause of $w_z$
 is the assumption $C(b,0)\equiv B(s,b,0)$, and its realizer of
 \cref{wit:lem:eqb-atomic}(i), at any term, is a family for it (the remark
 after \cref{wit:def:ab-family}).  If $\epsilon_C(z)=z$, the
 matrix $C$ is diagonal and \cref{wit:lem:eqb-zero} gives $w_z\Vdash C(b,0)$;
 the assumption $C(b,z)$ is not used.
\item $w_z\Vdash B(t,b,0)$, by the induction hypothesis (ii) for the template
 $B(q,b,0)$, which has $d(H)-1$ binders, applied to (a).
\item $w_z\Vdash B(t,b,\Suc z)$: the opened clause $[B(t,b,0);\,y.B(t,b,y)]$
 is a clause of $\Gamma_{w_z}$, and its realizer of
 \cref{wit:lem:eqb-atomic}(i), at the term $\Suc z$ and with the family of
 (b), gives this family.
\item $w_z\Vdash B(s,b,\Suc z)$, by the induction hypothesis (i) for the
 template $B(q,b,\Suc z)$, applied to (c).
\end{enumerate}
The member at $w_z$ of the family (d) is the premise of the
\rulename{Close}, which yields $H(s)$ in $w_b$.  $H(s)$ does not contain $b$,
and \rulename{Open} yields $H(s)$ at $w'$.  The construction is uniform in
$w'\succeq w$, so $w\Vdash H(s)$.

Every backward substitution occurs in \cref{wit:lem:eqb-atomic}, on
equations.  The recursive calls concern templates with fewer binders.  Only
the mode of the target matrix $C$ is used, and only through
\cref{wit:lem:eqb-zero}; the clause opened by \rulename{Open} has the same
shape in either mode, so the modes of $B(t,x,y)$ and $B(s,x,y)$ are never
compared.  The auxiliary parameters opened by \cref{wit:lem:eqb-atomic} above
$w_b$ and $w_z$ are discharged by \rulename{Alias} before a member at the
requested world is returned, so they do not change the parameter that the
\rulename{Close} in $w_b$ selects.  No derivation is moved to another world
or receives a substitution, and the lemma does not use
\cref{wit:thm:eqback}.
\end{proof}

\begin{theorem}[Equation-only and one-term \rulename{Back}]
\label{wit:thm:eqback}
For every scope $\Pi$, clause context $\Gamma$ and formula $A$ over $\Pi$,
\[
 \Pi;\Gamma\proves_{\mathsf{A}^{=}_{\nu}}A
 \iff \Pi;\Gamma\proves_{\mathsf{A}^{\times}_{\nu}}A
 \iff \Pi;\Gamma\proves_{\mathsf{A}_{\nu}}A
 \iff \Pi;\Gamma\proves_{\mathsf{A}_{\ast}}A ,
\]
by effective transformations of derivations that leave the terms, formulas
and clauses of the judgment unchanged.
\end{theorem}

\begin{proof}
The last equivalence is \cref{wit:thm:ab-newest}.  Let $\mathsf{C}$ be
$\mathsf{A}^{=}_{\nu}$ or $\mathsf{A}^{\times}_{\nu}$.

\emph{From $\mathsf{C}$ to $\mathsf{A}_{\ast}$.}  Every rule of
$\mathsf{A}^{=}_{\nu}$ is a rule of $\mathsf{A}_{\nu}$, hence of
$\mathsf{A}_{\ast}$.  Every rule of $\mathsf{A}^{\times}_{\nu}$ other than
\rulename{XBack} is a rule of $\mathsf{A}_{\ast}$.  By the proof of
\cref{wit:cor:ab-bridge}, reflexivity, symmetry, transitivity and transport in
every formula are derived rules of $\mathsf{A}_{\ast}$ over every scope and
context.  Given derivations of $\mathsf{A}_{\ast}$ of $\Pi;\Gamma\proves b=t$
and $\Pi;\Gamma\proves h\subst{t}{b}=r$, let $q$ be a variable occurring in
none of $h$, $t$, $r$, $\Pi$, and put
$G(q)\defeq\bigl(h\subst{q}{b}=h\subst{t}{b}\bigr)$.  Then
$G(t)\equiv(h\subst{t}{b}=h\subst{t}{b})$, an instance of reflexivity, and
$G(b)\equiv(h=h\subst{t}{b})$; transport along $b=t$ gives
$h=h\subst{t}{b}$, transitivity with the minor premise gives $h=r$, and
symmetry gives $r=h$.  Replacing each \rulename{XBack} by this derivation
translates $\mathsf{A}^{\times}_{\nu}$ into $\mathsf{A}_{\ast}$.

\emph{From $\mathsf{A}_{\ast}$ to $\mathsf{C}$.}  The proof of
\cref{wit:thm:ab-newest} is repeated with families and realizers of
$\mathsf{C}$.  That proof applies \rulename{Alias}, \rulename{Open},
\rulename{Close} and the computations directly, and \rulename{Back} only
through \cref{wit:lem:ab-eq}: part (i) for the realizers of clauses, part
(iii) in the case of \rulename{Back}, with the template
$(\theta F)\subst{q}{\theta b}$, and parts (iii) and (iv) in the case of
\rulename{Close}, with the templates $\theta B(q,\epsilon_B(z'))$ and
$\theta B(q,\Suc z')$, which are compound when $B$ has a binder.  Part (i)
is replaced by \cref{wit:lem:eqb-atomic}(i), and parts (iii) and (iv) by
\cref{wit:lem:eqb-transport}(i) and (ii), which hold for $\mathsf{C}$.  All
other steps are unchanged; in particular the case of \rulename{Back} at a
source $b$ with a replacement $t$ that contains $b$ uses
$H(\theta t)\equiv\theta(F\subst{t}{b})$, which holds because substitution
is simultaneous and inserts each replacement once.  The induction is well
founded as before: \cref{wit:lem:eqb-atomic,wit:lem:eqb-zero,wit:lem:eqb-transport}
are constructed before and independently of it, by recursion on the number
of binders of a template, and the realizer of the source assumption in the
case of \rulename{Close} is built from the actual clauses of $w_z$ by
\cref{wit:lem:eqb-transport}, without a call of the translation of the
premise.  The derivation of $\Pi;\Gamma\proves A$ is the member at
$(\Pi;\Gamma)$ of the family for the identity map and the realizers of
\cref{wit:lem:eqb-atomic}(i).
\end{proof}

\begin{remark}[Two steps of \rulename{XBack} for one \rulename{EqBack}]
\label{wit:rem:eqback-sim}
A translation from $\mathsf{A}^{=}_{\nu}$ to $\mathsf{A}^{\times}_{\nu}$ that
does not pass through families replaces an \rulename{EqBack} with the
conclusion $a=d$, the major premise $b=t$ and the minor premise
$a\subst{t}{b}=d\subst{t}{b}$ by two applications of \rulename{XBack} in the
same world and with the same major premise: with $h\defeq a$ and
$r\defeq d\subst{t}{b}$ it gives $d\subst{t}{b}=a$, and then with $h\defeq d$
and $r\defeq a$ it gives $a=d$.  This holds when $t$ contains $b$, because
only $h$ is substituted.  The derivation of the major premise is repeated.
\end{remark}

\begin{corollary}[The restricted Alias/Back calculi and $\iPRA$]
\label{wit:cor:eqback-bridge}
With $\Gamma^{\circ}$ as in \cref{wit:cor:ab-bridge},
\[
 \Pi;\Gamma\proves_{\mathsf{A}^{=}_{\nu}}A
 \iff \Pi;\Gamma\proves_{\mathsf{A}^{\times}_{\nu}}A
 \iff \Gamma^{\circ}\proves_{\mathsf{K}}A .
\]
The ordinary realizations of $\mathsf{A}^{=}_{\nu}$ and
$\mathsf{A}^{\times}_{\nu}$ are bi-interpretable and equiconsistent with
$\iPRA$, with whole numerical domains, actual equality, the identity graph
$x=y$ in both numerical composites, provable formula round trips in both
directions, and the contradiction $0=\Suc 0$.
\end{corollary}

\begin{proof}
By \cref{wit:thm:eqback} both calculi derive exactly the judgments of
$\mathsf{A}_{\ast}$, on the same terms, formulas and clause contexts, by
effective translations.  \Cref{wit:cor:ab-bridge} gives the comparison with
$\mathsf{K}$ and the bridge for $\mathsf{A}_{\ast}$; its translations, the
derivations of both formula round trips and the preservation of $0=\Suc 0$
transfer unchanged.
\end{proof}

The equivalence is one of open consequence relations, on every scope and
clause context, and not only of closed theorems or of consistency.  The
recursion on formulas in \cref{wit:lem:eqb-transport} is a recursion of the
metatheory that constructs derivations; no induction on compound arithmetic
formulas enters the calculi.

\begin{remark}[What the proof uses of the mode]
\label{wit:rem:eqback-mode}
\Cref{wit:lem:eqb-transport} uses the mode only through the target matrix of
a \rulename{Close}, and only in this form: every matrix whose mode selects
$z$ has a zero row that is derivable on families without formula transport.
The translation of \cref{wit:thm:eqback} uses in addition
$\epsilon_{\theta B}=\epsilon_B$ for injective renamings $\theta$
(\cref{wit:rem:ab-mode}).  The same proof therefore restricts
\rulename{Back} to equations, or to \rulename{XBack}, for every mode
$\epsilon_B(z)\in\{0,z\}$ determined by the uninstantiated matrix, preserved
and reflected by injective renaming, such that every matrix of mode $z$ is a
literal equation whose zero row follows from the computations and
\cref{wit:lem:eqb-atomic}.  The mirror guard of \cref{claus:prop:mirror}
qualifies: the zero row of $a=\It{y}{a}{u}{h}$ is $a'=\It{0}{a'}{u}{h'}$,
which is \rulename{I0} and symmetry.  As in \cref{wit:rem:ab-mode} this is a
statement of relative derivability; soundness and arithmetic adequacy of the
mode are separate obligations.  It does not extend to modes that select $z$
for arbitrary progressive equations, whose zero rows need not be derivable.
\end{remark}

\begin{remark}[The package]
\label{wit:rem:eqback-package}
\pkg{equational\_transport} proves \cref{wit:thm:eqback} in the order of
this subsection, states \cref{wit:lem:eqb-transport} with the selector
condition of \cref{wit:rem:eqback-mode}, and gives both the two-step
simulation of \cref{wit:rem:eqback-sim} and the direct construction of
\cref{wit:lem:eqb-atomic}(iii) for $\mathsf{A}^{\times}_{\nu}$, with four
applications of \rulename{XBack}.  The proof above uses the direct
construction, so that \cref{wit:thm:eqback} holds for
$\mathsf{A}^{\times}_{\nu}$ without the simulation; the simulation is the
step of the package's compiler from $\mathsf{A}^{=}_{\nu}$ to
$\mathsf{A}^{\times}_{\nu}$.  An independent mathematical review in
the package checks the order of the lemmas and the absence of a weakening of
completed derivations.  The theorem is not formalized in Lean or Rocq.
\end{remark}

\subsection{Equation-only aliases and one naming scheme}
\label{wit:sec:ns}

This subsection draws on \pkg{names\_and\_scopes}.

\Cref{wit:thm:eqback} restricts \rulename{Back} to equations and keeps the
export of an arbitrary formula by \rulename{Alias}.  \Cref{disch:thm:alias-atomic}
restricts both \rulename{Back} and \rulename{Alias} to equations and keeps a
supplied witness of \rulename{Close}, which derives the name of every term
(\cref{disch:lem:alias-name}(a)).  The two restrictions together with the
introduction at the newest parameter lose completeness: the closed formula
$\mathrm{Name}(1)$ below is no longer derivable
(\cref{wit:thm:ns-obstruction}).  One axiom scheme of rank one, the name of
a term by its alias matrix $\mathrm{Al}_t$ of \cref{wit:def:naming},
restores every open consequence (\cref{wit:thm:ns-complete}); finitely many
naming templates instantiated at parameters do not
(\cref{wit:thm:ns-templates}).

\begin{definition}[The calculi $\mathsf{A}^{\mathrm e}_{\nu}$ and
$\mathsf{A}^{\mathrm n}_{\nu}$]
\label{wit:def:ns}
$\mathsf{A}^{\mathrm e}_{\nu}$ has the syntax, judgments and mode
$\epsilon_B$ of \cref{wit:def:ab} and exactly the following seven rule
patterns: \rulename{I0}, \rulename{IS}, \rulename{Use} at the newest
parameter, \rulename{Open} and \rulename{Close} at the newest parameter, as
in $\mathsf{A}_{\nu}$; \rulename{EqBack} of \cref{wit:def:eqback}; and, in
place of \rulename{Alias},
\begin{itemize}
\item \rulename{EqAlias}: if $b\notin\Pi$, $t$ is over $\Pi$ and $E$ is an
 equation over $\Pi,b$, from $\Pi,b;\ \Gamma,b=t\proves E$ infer
 $\Pi;\Gamma\proves E\subst{t}{b}$.
\end{itemize}
$\mathsf{A}^{\mathrm n}_{\nu}$ is $\mathsf{A}^{\mathrm e}_{\nu}$ together
with the axiom scheme
\[
 \Pi;\Gamma\proves\mathrm{Name}(t)\quad\text{for every term $t$ over $\Pi$},
 \qquad
 \mathrm{Name}(t)\defeq\D{x}{y}\bigl(\sel(y,t,x)=t\bigr),
\]
called \rulename{Names}; the bound variables $x$, $y$ and the bound variable
of the selector are not free in $t$.  For a parameter $b$ not in $t$ put
\[
 \mathcal N_{b,t}\defeq[\sel(0,t,b)=t;\ y.\ \sel(y,t,b)=t].
\]
\end{definition}

The lists are exhaustive: there is no reflexivity, transport, weakening,
cut or induction rule, and \rulename{Use} and \rulename{Close} take no term.
\rulename{EqAlias} is $\rulename{Alias}^{=}$
of \cref{disch:def:alias}, and $\mathsf{A}^{\mathrm e}_{\nu}$ is the calculus
$\mathcal A$ of \cref{disch:def:alias-atomic} with \rulename{Close} at the
newest parameter in the mode $\epsilon_B$.  $\mathrm{Name}(t)$ is the name
$\D{x}{y}\mathrm{Al}_t$ of \cref{wit:def:naming} and the formula
$\mathrm{Name}(t)$ of \cref{disch:lem:alias-name}; its rank is one, and its
opening by \rulename{Open} at a fresh $b$ adds the clause
$\mathcal N_{b,t}$.  Its matrix $\sel(y,t,x)=t$ is not diagonal, since its
right side $t$ does not contain $y$; an introduction of a name is in the
case mode.  Its ordinary reading is equivalent to $\exists x\,(x=t)$
(\cref{wit:def:naming}).  Every rule of $\mathsf{A}^{\mathrm e}_{\nu}$ is a
rule of $\mathsf{A}^{=}_{\nu}$, since \rulename{EqAlias} is the instance of
\rulename{Alias} with an equation as its conclusion.

\begin{lemma}[One compound export]
\label{wit:lem:ns-shape}
Let $\mathsf{A}^{\mathrm e+}_{\nu}$ be $\mathsf{A}^{\mathrm e}_{\nu}$ with
the instance of \rulename{Alias} whose premise concludes a name of the
alias:
\[
 \infrule{\Pi,b;\ \Gamma,b=t\proves\mathrm{Name}(b)}
         {\Pi;\Gamma\proves\mathrm{Name}(t)}\ \rulename{NameAlias}
 \qquad b\notin\Pi,\ t\text{ over }\Pi .
\]
Then $\mathsf{A}^{\mathrm e+}_{\nu}$ and $\mathsf{A}^{\mathrm n}_{\nu}$
derive the same judgments.
\end{lemma}

\begin{proof}
Replace each \rulename{NameAlias} of a derivation of
$\mathsf{A}^{\mathrm e+}_{\nu}$ by the axiom $\Pi;\Gamma\proves\mathrm{Name}(t)$
and delete the derivation of its premise.  Conversely, an instance
$\Pi;\Gamma\proves\mathrm{Name}(t)$ of \rulename{Names} is derived in
$\mathsf{A}^{\mathrm e+}_{\nu}$ as follows.  In the scope $\Pi,b$ the
parameter $b$ is the newest, so \rulename{Close} for the matrix
$\sel(y,b,x)=b$ selects $b$; the matrix is not diagonal, and the premise asks,
in the scope $\Pi,b,z$ and from the unused assumption $\sel(0,b,b)=b$, for
$\sel(\Suc z,b,b)=b$, which is \rulename{IS} with the constant body $b$.  This
derives $\Pi,b;\ \Gamma,b=t\proves\mathrm{Name}(b)$, and \rulename{NameAlias}
gives $\mathrm{Name}(b)\subst{t}{b}\equiv\mathrm{Name}(t)$.  Neither direction
uses another compound export or a compound backward substitution.
\end{proof}

Thus the scheme \rulename{Names} is the same as allowing \rulename{Alias} to
export compound formulas of one fixed shape.  It restricts the shape of the
conclusions an alias may export, not the number of rule patterns.

\begin{theorem}[Equation-only aliases and the newest witness]
\label{wit:thm:ns-obstruction}
$\mathsf{A}^{\mathrm e}_{\nu}$ does not derive
$\varnothing;\varnothing\proves\mathrm{Name}(1)$, which $\mathsf{A}_{\ast}$
derives.  More generally, let $K\ge 0$, write $\overline{k}$ for the numeral
$\Suc^k0$, and let $\mathsf{A}^{\mathrm e}_{K}$ have the rules
\rulename{I0}, \rulename{IS}, \rulename{Open} and \rulename{EqAlias}, and
\begin{itemize}
\item \rulename{Use} at a parameter of $\Pi$ or at a numeral $\overline{k}$
 with $k\le K$;
\item transport in every formula: for terms $s$, $t$ over $\Pi$ and a formula
 $F$ with a distinguished variable $q$, from $\Pi;\Gamma\proves s=t$ and
 $\Pi;\Gamma\proves F\subst{t}{q}$ infer $\Pi;\Gamma\proves F\subst{s}{q}$, and
 from $\Pi;\Gamma\proves s=t$ and $\Pi;\Gamma\proves F\subst{s}{q}$ infer
 $\Pi;\Gamma\proves F\subst{t}{q}$;
\item \rulename{Close} at a witness $w$ that is a parameter of $\Pi$ or a
 numeral $\overline{k}$ with $k\le K$, for every matrix with either of the
 premises $\Pi,z;\ \Gamma,B(w,0)\proves B(w,\Suc z)$ (cases) and
 $\Pi,z;\ \Gamma,B(w,z)\proves B(w,\Suc z)$ (progress).
\end{itemize}
Then $\mathsf{A}^{\mathrm e}_{K}$ does not derive
$\varnothing;\varnothing\proves\mathrm{Name}(\overline{K+1})$.  Every rule of
$\mathsf{A}^{\mathrm e}_{\nu}$ is a rule of $\mathsf{A}^{\mathrm e}_{0}$.
\end{theorem}

\begin{proof}
The last statement: $\nu(\Pi)$ is a parameter or the numeral $0$;
\rulename{EqBack} with the major premise $b=t$ is transport in the template
$F\defeq E\subst{q}{b}$ from $F\subst{t}{q}\equiv E\subst{t}{b}$ to
$F\subst{b}{q}\equiv E$; and \rulename{Close} in the mode $\epsilon_B$ has one
of the two premises.  In $\mathsf{A}_{\ast}$, \rulename{Close} at the witness
$1$ for the matrix $\sel(y,1,x)=1$ has the premise $\sel(\Suc z,1,1)=1$, an
instance of \rulename{IS}.

\emph{The evaluation.}  The \emph{worlds} are the integers $n\ge K$, ordered
by $\le$; the \emph{active domain} of $n$ is $A_n=\{0,1,\dots,n\}$.  An
environment $\rho$ assigns a natural number to every variable, and
$\llbracket t\rrbracket_\rho\in\mathbb{N}$ is the standard value of a term; it may
exceed $n$.  For a world $n$ and an environment $\rho$ define
\begin{align*}
 n,\rho\Vdash s=t &\iffdef \llbracket s\rrbracket_\rho=\llbracket t\rrbracket_\rho,\\
 n,\rho\Vdash\D{x}{y}B &\iffdef \text{for some $a\in A_n$ and all $m\ge n$:
   if $m,\rho[x\mapsto a,y\mapsto 0]\Vdash B$, then}\\
 &\qquad\ \ \ell,\rho[x\mapsto a,y\mapsto c]\Vdash B
   \text{ for all $\ell\ge m$ and $c\in A_\ell$},
\end{align*}
and for clauses
\begin{align*}
 n,\rho\Vdash[A;\,y.B] &\iffdef \text{for all $m\ge n$ with $m,\rho\Vdash A$,
   all $\ell\ge m$ and $c\in A_\ell$: } \ell,\rho[y\mapsto c]\Vdash B,\\
 n,\rho\Vdash[\ ;\,y.B] &\iffdef \text{for all $\ell\ge n$ and $c\in A_\ell$: }
   \ell,\rho[y\mapsto c]\Vdash B .
\end{align*}
These are the increasing-domain forcing clauses of the ordinary readings,
with existential witnesses from the current domain and universal
quantifiers over the domains of all later worlds, except that terms are
evaluated in the natural numbers and not in the active domain.  The
evaluation is an invariant of derivations and not a Kripke model of a
first-order theory, since a term with active arguments may have a value
outside the active domain, as $\Suc n$ at the world $n$.
The following facts hold, each by induction on formulas.
\begin{enumerate}[label=\textup{(\alph*)}]
\item Monotonicity: if $n\le n'$ and $n,\rho\Vdash A$, then $n',\rho\Vdash A$;
 the same for clauses.  Equations do not depend on the world; for
 $\D{x}{y}B$ the witness $a\in A_n\subseteq A_{n'}$ is kept, and the
 condition at $n'$ quantifies over fewer worlds $m$.
\item Forcing depends only on the values of the free variables.
\item Substitution: $n,\rho\Vdash F\subst{t}{b}$ if and only if
 $n,\rho[b\mapsto\llbracket t\rrbracket_\rho]\Vdash F$, the bound variables of
 $F$ being renamed away from $t$.  The forcing relation is defined for all
 environments, also those with values outside the active domain, so the
 induction needs no condition on $\llbracket t\rrbracket_\rho$.
\item By (a), $n,\rho\Vdash[\ ;\,v.F]$ if and only if $n,\rho\Vdash F$, for
 $v$ not in $F$.
\end{enumerate}
A judgment $\Pi;\Gamma\proves A$ is \emph{valid} if for every world $n$ and
every environment $\rho$ with $\rho(p)\in A_n$ for all $p\in\Pi$: if
$n,\rho$ forces every clause of $\Gamma$, then $n,\rho\Vdash A$.

\emph{Every rule of $\mathsf{A}^{\mathrm e}_{K}$ preserves validity.}  Let
$n$ and $\rho$ be as in the definition, with $\Gamma$ forced at $n,\rho$.
The values of all parameters of $\Pi$, and of all numerals $\overline{k}$
with $k\le K\le n$, lie in $A_n$.
\begin{itemize}
\item \rulename{I0} and \rulename{IS} are true equations of $\mathbb{N}$.
\item \rulename{Use} of $[A;\,y.B]$ at $w$: the premise gives $n,\rho\Vdash A$;
 with $m=\ell=n$ and $c=\llbracket w\rrbracket_\rho\in A_n$ the clause gives
 $n,\rho[y\mapsto c]\Vdash B$, which is $n,\rho\Vdash B(w)$ by (c).  The
 clause without antecedent is the same.
\item Transport: $\llbracket s\rrbracket_\rho=\llbracket t\rrbracket_\rho$, and
 by (c) $F\subst{s}{q}$ and $F\subst{t}{q}$ are forced together.
\item \rulename{EqAlias} with the premise $\Pi,b;\ \Gamma,b=t\proves E$: put
 $a=\llbracket t\rrbracket_\rho$, $m=\max(n,a)$ and $\rho'=\rho[b\mapsto a]$.
 Then $\rho'$ sends $\Pi,b$ into $A_m$; $\Gamma$ is forced at $m,\rho$ by (a)
 and at $m,\rho'$ by (b), because $b$ is not free in $\Gamma$; and
 $m,\rho'\Vdash b=t$, because $b$ is not in $t$.  Validity of the premise
 gives $m,\rho'\Vdash E$.  $E$ is an equation, so $n,\rho'\Vdash E$, and
 $n,\rho\Vdash E\subst{t}{b}$ by (c).  The world is changed only for an
 equation; a compound conclusion forced at $m$ need not be forced at $n$.
\item \rulename{Open} with the major premise $\D{x}{y}B$ and the fresh $b$:
 let $a\in A_n$ be a witness of the major premise and
 $\rho'=\rho[b\mapsto a]$, which sends $\Pi,b$ into $A_n$.  By (b) and (c),
 $n,\rho'$ forces $\Gamma$ and the opened clause $[B(b,0);\,y.B(b,y)]$, whose
 forcing condition is literally the condition on the witness $a$.  The branch
 gives $n,\rho'\Vdash A$, and $n,\rho\Vdash A$ since $b$ is not in $A$.
\item \rulename{Close} at $w$ with the eigenvariable $z$: put
 $a=\llbracket w\rrbracket_\rho\in A_n$; it is shown that $a$ witnesses
 $\D{x}{y}B$ at $n,\rho$.  Let $m\ge n$ with
 $m,\rho[x\mapsto a,y\mapsto 0]\Vdash B$, and let $\ell\ge m$ and
 $c\in A_\ell$.  For $j\in A_\ell$ let $\rho_j=\rho[z\mapsto j]$; it sends
 $\Pi,z$ into $A_\ell$, forces $\Gamma$ at $\ell$ by (a) and (b), and by (c)
 and (d) forces the assumption $B(w,e)$ at $\ell$ exactly when
 $\ell,\rho[x\mapsto a,y\mapsto\llbracket e\rrbracket_{\rho_j}]\Vdash B$.
 In the case mode: for $c=0$ the claim is monotonicity; for $c>0$ take
 $j=c-1\in A_\ell$, where the assumption $B(w,0)$ is forced by monotonicity,
 and the premise gives $\ell,\rho_j\Vdash B(w,\Suc z)$, that is,
 $\ell,\rho[x\mapsto a,y\mapsto c]\Vdash B$.  In the progress mode: by
 induction on $j\le c$ in the metatheory, $\ell,\rho[x\mapsto a,y\mapsto j]\Vdash B$;
 the base is monotonicity, and for $j<c$ the number $j$ lies in $A_\ell$, the
 assumption $B(w,z)$ is forced at $\ell,\rho_j$, and the premise gives the
 instance at $j+1$.  This holds for an arbitrary matrix $B$.
\end{itemize}
By induction on derivations every judgment derivable in
$\mathsf{A}^{\mathrm e}_{K}$ is valid.

\emph{The name of $K+1$.}  At the world $K$, with the empty scope and
context, let $a\in A_K$ be a proposed witness of
$\mathrm{Name}(\overline{K+1})$.  Its row at $y=0$, $\sel(0,K+1,a)=K+1$, is
true, so the condition applies at $m=K$; at $\ell=\max(K,1)$ and $c=1$ the
row $\sel(1,K+1,a)=K+1$ says $a=K+1$, which is false because $a\le K$.  So
$\mathrm{Name}(\overline{K+1})$ is not forced at $K$, the judgment is not
valid and not derivable.  For $K=0$ this is the statement on
$\mathsf{A}^{\mathrm e}_{\nu}$.
\end{proof}

\begin{remark}[What the obstruction uses]
\label{wit:rem:ns-obstruction}
The evaluation validates neither unrestricted operation of the complete
calculi.  \rulename{Close} of $\mathsf{A}_{\ast}$ may take a witness whose
value is outside the active domain, and \rulename{Alias} of
$\mathsf{A}_{\nu}$ may export a compound formula that is forced only after
the domain has grown: $\mathrm{Name}(b)$ is forced at the world
$\max(n,\llbracket t\rrbracket_\rho)$ under $b\mapsto\llbracket t\rrbracket_\rho$,
and its export $\mathrm{Name}(t)$ need not be forced at $n$.  Unlike the
invariants of \cref{wit:thm:finite-range,wit:thm:finite-range-general},
which bound the quantifiers by a fixed finite set and restrict the target of
the elimination, this one lets the domain grow and restricts which
conclusions may be pulled back to an earlier world.  The growing domains
separate the restricted rule systems only; they are not the domains of the
bridge of \cref{wit:cor:ns-bridge}, and the theorem does not exclude that a
differently interpreted theory with the rules of $\mathsf{A}^{\mathrm e}_{\nu}$
is bi-interpretable with $\iPRA$.
\end{remark}

For the positive result, persistent families, persistent realizers and
$w\Vdash A$ are defined for $\mathsf{A}^{\mathrm n}_{\nu}$ as in
\cref{wit:def:ab-family}, with derivations of $\mathsf{A}^{\mathrm n}_{\nu}$
in place of derivations of $\mathsf{A}_{\nu}$, and $q$ denotes a
distinguished template variable as in \cref{wit:sec:eqback}.  The
construction of \cref{disch:sec:alias-atomic} is carried out with the
introduction at the newest parameter; where the supplied witness of that
section introduces a name or an opened clause, a \rulename{Names} axiom is
opened, and the progress mode of a target introduction is handled by the
diagonal zero rows of \cref{wit:lem:eqb-zero}.

\begin{lemma}[The equational base and the names]
\label{wit:lem:ns-base}
In $\mathsf{A}^{\mathrm n}_{\nu}$, let $w$ be a world.
\begin{enumerate}[label=\textup{(\roman*)}]
\item If an assumption $[\ ;\,v.F]$ occurs in $\Gamma_w$, then $w\Vdash F$.
\item The statements (ii)--(iv) of \cref{wit:lem:eqb-atomic} and
 \cref{wit:lem:eqb-zero} hold.
\item $(\mathrm S_0)$ If $[A;\,y.E(y)]$ occurs in $\Gamma_w$ with $E$ an
 equation and $w\Vdash A$ (no condition for a clause without antecedent),
 then $w\Vdash E(r)$ for every term $r$ over $\Pi_w$.  Hence such a clause
 has a persistent realizer at $w$.
\item Every instance of \rulename{Names} over $\Pi_w$ has a persistent
 family at $w$.  If $\mathcal N_{b,t}$ occurs in $\Gamma_w$, at any position,
 and $b$ is a parameter of $\Pi_w$ not in $t$, then $w\Vdash b=t$.
\end{enumerate}
\end{lemma}

\begin{proof}
(i) In every $w'\succeq w$, \rulename{Use} of the clause at
$\nu(\Pi_{w'})$ gives $F$, because $v$ is not in $F$.  No \rulename{Alias}
is used.

(ii) The proofs for $\mathsf{A}^{=}_{\nu}$ in \cref{wit:lem:eqb-atomic}
apply \rulename{Alias} only with the conclusions $t=t$ in (ii), $E(s)$ in
(iii), where $E$ is an equation, and $t=s$ through (iii) in (iv); these are
applications of \rulename{EqAlias}.  Their other rules are \rulename{Use} at
the newest parameter and \rulename{EqBack}, and the alias assumption is used
through \rulename{Use} of the last clause.  The proof of
\cref{wit:lem:eqb-zero} uses \rulename{I0} and these statements only.

(iii) Let $w'\succeq w$, choose $c\notin\Pi_{w'}$ and work in
$w'_c=(\Pi_{w'},c;\ \Gamma_{w'},c=r)$, which extends $w$.  The member of the
family for $A$ at $w'_c$ is the premise of \rulename{Use} of the clause at
the newest parameter $c$, which gives $E(c)$; \rulename{EqAlias} gives
$E(c)\subst{r}{c}\equiv E(r)$ at $w'$.  The clause is a clause of every
$w'\succeq w$, so the construction at $w'$ is a realizer.

(iv) The instance is an axiom in every extension.  For the second statement
follow \cref{disch:lem:alias-name}(b).  The antecedent $\sel(0,t,b)=t$ is an
instance of \rulename{I0} and has a persistent family; (iii) at the term $1$
gives $w\Vdash\sel(1,t,b)=t$; \rulename{IS} gives $\sel(1,t,b)=b$ in every
world, the bound variable of the selector not being $b$; and the symmetry and
transport of (ii) in the template $q=b$ give $w\Vdash t=b$ and then
$w\Vdash b=t$.  The specialization of (iii) opens its own alias, so neither
the clause nor $b$ needs to be the newest.
\end{proof}

\begin{theorem}[Transport and specialization by rank]
\label{wit:thm:ns-rank}
For every $k\ge 0$ and every world $w$ the following hold in
$\mathsf{A}^{\mathrm n}_{\nu}$.  The rank of a formula is its number of
binders $\D{}{}$.
\begin{itemize}
\item[$(\mathrm T_k)$] For terms $s$, $t$ over $\Pi_w$ and a formula $H(q)$
 of rank at most $k$ over $\Pi_w$ and $q$: if $w\Vdash s=t$ and
 $w\Vdash H(s)$, then $w\Vdash H(t)$.
\item[$(\mathrm S_k)$] If $[A;\,y.F(y)]$ occurs in $\Gamma_w$ with $F$ of
 rank at most $k$, and $w\Vdash A$ (no condition for a clause without
 antecedent), then $w\Vdash F(r)$ for every term $r$ over $\Pi_w$.
\item[$(\mathrm L_k)$] For $k>0$ and $H(q)$ of rank exactly $k$: from a
 derivation of $\Pi_w;\Gamma_w\proves H(s)$ and $w\Vdash s=t$ a derivation
 of $\Pi_w;\Gamma_w\proves H(t)$ is constructed.
\end{itemize}
\end{theorem}

\begin{proof}
$(\mathrm T_0)$ is \cref{wit:lem:eqb-atomic}(iv) and $(\mathrm S_0)$ is
\cref{wit:lem:ns-base}(iii), both in $\mathsf{A}^{\mathrm n}_{\nu}$ by
\cref{wit:lem:ns-base}(ii).  Let $k>0$ and assume $(\mathrm T_{k-1})$ and
$(\mathrm S_{k-1})$.

\emph{$(\mathrm L_k)$.}  Write $H(q)=\D{x}{y}B(q,x,y)$, with $B$ of rank
$k-1$.  Choose $c\notin\Pi_w$.  The given derivation of $H(s)$ is the major
premise of \rulename{Open} at $c$, whose branch has the world
\[
 w_c=\bigl(\Pi_w,c;\ \Gamma_w,\mathcal C\bigr),\qquad
 \mathcal C\defeq[B(s,c,0);\,y.B(s,c,y)],
\]
and the goal $H(t)$, which is over $\Pi_w$.  The given derivation is used
once, in the world in which it is given.  In $w_c$ apply \rulename{Close} to
$\D{x}{y}C$ with $C(x,y)\defeq B(t,x,y)$; the newest parameter is $c$, so the
rule selects $c$.  With its eigenvariable $z$ the premise is to be derived in
\[
 w_z=\bigl(\Pi_w,c,z;\ \Gamma_w,\mathcal C,C(c,\epsilon_C(z))\bigr)
\]
with the goal $B(t,c,\Suc z)$.  The world $w_z$ extends $w$, so $w_z\Vdash s=t$
by restriction and $w_z\Vdash t=s$ by symmetry.
\begin{enumerate}[label=\textup{(\alph*)}]
\item $w_z\Vdash B(t,c,0)$.  If $\epsilon_C(z)=0$, the last clause of $w_z$
 is the assumption $B(t,c,0)$, and \cref{wit:lem:ns-base}(i) applies.  If
 $\epsilon_C(z)=z$, the matrix $C$ is diagonal, and \cref{wit:lem:eqb-zero}
 gives $w_z\Vdash C(c,0)$; the assumption $C(c,z)$ is not used.
\item $w_z\Vdash B(s,c,0)$, by $(\mathrm T_{k-1})$ for the template
 $B(q,c,0)$ and the equation $t=s$.
\item $w_z\Vdash B(s,c,\Suc z)$, by $(\mathrm S_{k-1})$ for the clause
 $\mathcal C$ of $\Gamma_{w_z}$, whose body has rank $k-1$, at the term
 $\Suc z$ and with the antecedent (b).
\item $w_z\Vdash B(t,c,\Suc z)$, by $(\mathrm T_{k-1})$ for the template
 $B(q,c,\Suc z)$ and the equation $s=t$.
\end{enumerate}
The member of (d) at $w_z$ is the premise of \rulename{Close}, which gives
$H(t)$ in $w_c$; $H(t)$ does not contain $c$, and \rulename{Open} gives
$\Pi_w;\Gamma_w\proves H(t)$.

\emph{$(\mathrm T_k)$.}  Templates of rank below $k$ are covered by
$(\mathrm T_{k-1})$.  For rank $k$ and $w'\succeq w$, the member at $w'$ of
the family for $H(s)$ is a derivation of $H(s)$ at $w'$, the family for
$s=t$ restricts to $w'$, and $(\mathrm L_k)$ at $w'$ gives the member for
$H(t)$.

\emph{$(\mathrm S_k)$.}  Bodies of rank below $k$ are covered by
$(\mathrm S_{k-1})$.  For rank $k$ and $w'\succeq w$: \rulename{Names} gives
$\Pi_{w'};\Gamma_{w'}\proves\mathrm{Name}(r)$; choose $b\notin\Pi_{w'}$ and
open the name at $b$, which gives the world
$w_b=(\Pi_{w'},b;\ \Gamma_{w'},\mathcal N_{b,r})$, where $w_b\Vdash b=r$ by
\cref{wit:lem:ns-base}(iv).  The member of the family for $A$ at $w_b$ is
the premise of \rulename{Use} of the clause at the newest parameter $b$, which
gives a derivation of $F(b)$ in $w_b$; $(\mathrm L_k)$ at $w_b$ turns it into
a derivation of $F(r)$ in $w_b$; and $F(r)$ does not contain $b$, so
\rulename{Open} gives $F(r)$ at $w'$.

The recursion is well founded: $(\mathrm L_k)$ uses
$(\mathrm T_{k-1})$, $(\mathrm S_{k-1})$ and
\cref{wit:lem:ns-base,wit:lem:eqb-zero}; $(\mathrm T_k)$ uses
$(\mathrm L_k)$; $(\mathrm S_k)$ uses $(\mathrm L_k)$ and the extraction of
the name, which uses rank $0$ only.  Every derivation is built in the world
where it is required; no derivation is moved to another world or receives a
substitution.
\end{proof}

By $(\mathrm S_k)$, applied in every $w'\succeq w$, every clause of
$\Gamma_w$ has a persistent realizer at $w$ in $\mathsf{A}^{\mathrm n}_{\nu}$,
whatever the rank of its body; and by $(\mathrm T_k)$ and the symmetry of
\cref{wit:lem:ns-base}(ii) the two statements of
\cref{wit:lem:eqb-transport} hold for $\mathsf{A}^{\mathrm n}_{\nu}$.  These
replace \cref{wit:lem:ab-eq}(i) and (iii), (iv) in the translation below.

\begin{theorem}[One naming scheme suffices]
\label{wit:thm:ns-complete}
For every scope $\Pi$, clause context $\Gamma$ and formula $A$ over $\Pi$,
\[
 \Pi;\Gamma\proves_{\mathsf{A}^{\mathrm n}_{\nu}}A
 \iff \Pi;\Gamma\proves_{\mathsf{A}^{=}_{\nu}}A
 \iff \Pi;\Gamma\proves_{\mathsf{A}_{\nu}}A
 \iff \Pi;\Gamma\proves_{\mathsf{A}_{\ast}}A ,
\]
by effective transformations of derivations that leave the terms, formulas
and clauses of the judgment unchanged.  In particular full \rulename{Alias},
\rulename{Back} on every formula, \rulename{Use} at every term and
\rulename{Close} at every witness are admissible in
$\mathsf{A}^{\mathrm n}_{\nu}$.
\end{theorem}

\begin{proof}
The last two equivalences are \cref{wit:thm:eqback}.

\emph{From $\mathsf{A}^{\mathrm n}_{\nu}$ to $\mathsf{A}^{=}_{\nu}$.}  Every
rule of $\mathsf{A}^{\mathrm e}_{\nu}$ is a rule of $\mathsf{A}^{=}_{\nu}$.
An instance $\Pi;\Gamma\proves\mathrm{Name}(t)$ of \rulename{Names} is
derived in $\mathsf{A}^{=}_{\nu}$ by the second half of the proof of
\cref{wit:lem:ns-shape}, with \rulename{Alias} in place of
\rulename{NameAlias}.

\emph{From $\mathsf{A}_{\ast}$ to $\mathsf{A}^{\mathrm n}_{\nu}$.}  The
induction statement of the proof of \cref{wit:thm:ab-newest} is proved with
families and realizers of $\mathsf{A}^{\mathrm n}_{\nu}$: for every world
$w$, every injective map $\theta$ from $\Pi$ into $\Pi_w$ and every
assignment of a persistent realizer at $w$ of $\theta C$ to each clause $C$
of $\Gamma$, $w\Vdash\theta A$.  The computations and \rulename{Use} are
treated as there.  \rulename{Back} is treated as there with
\cref{wit:lem:eqb-transport}(i), which holds for
$\mathsf{A}^{\mathrm n}_{\nu}$ by $(\mathrm T_k)$ and symmetry.
\rulename{Open} is treated as there, the opened clause being realized by
$(\mathrm S_k)$ instead of \cref{wit:lem:ab-eq}(i).  The cases of
\rulename{Alias} and \rulename{Close} in that proof end with an
\rulename{Alias} that exports a compound formula and are replaced as
follows; $w'\succeq w$ is the world at which a member is constructed.

\emph{Alias} with the parameter $b$, the term $t$ and the branch
$\Pi,b;\ \Gamma,b=t\proves F$.  \rulename{Names} gives
$\mathrm{Name}(\theta t)$ at $w'$.  Choose $c\notin\Pi_{w'}$ and open it at
$c$; the branch has the world
$w_c=(\Pi_{w'},c;\ \Gamma_{w'},\mathcal N_{c,\theta t})$, and
$w_c\Vdash c=\theta t$ by \cref{wit:lem:ns-base}(iv).  The map
$\theta'=\theta[b\mapsto c]$ is injective; the clauses of $\Gamma$ keep their
realizers, restricted to $w_c$, and the source assumption $b=t$, whose image
is $c=\theta t$, is realized by the family for $c=\theta t$ (a family for an
equation is a realizer of the assumption, after \cref{wit:def:ab-family}).
The induction hypothesis gives $w_c\Vdash\theta'F$.  With
$H(q)\defeq(\theta'F)\subst{q}{c}$ one has $H(c)\equiv\theta'F$ and
$H(\theta t)\equiv\theta(F\subst{t}{b})$, and $(\mathrm T_k)$ gives
$w_c\Vdash\theta(F\subst{t}{b})$.  Its member at $w_c$ does not contain $c$
in its conclusion, and \rulename{Open} with the major premise
$\mathrm{Name}(\theta t)$ gives $\theta(F\subst{t}{b})$ at $w'$.

\emph{Close} at the witness $t$ with the matrix $B$ and the premise
$\Pi,z;\ \Gamma,B(t,\epsilon_B(z))\proves B(t,\Suc z)$.  \rulename{Names}
gives $\mathrm{Name}(\theta t)$ at $w'$; choose $b\notin\Pi_{w'}$ and open it
at $b$, which gives $w_b=(\Pi_{w'},b;\ \Gamma_{w'},\mathcal N_{b,\theta t})$.
In $w_b$ apply \rulename{Close} to $\D{x}{y}\theta B$; the newest parameter is
$b$, so the rule selects $b$; let $z'$ be its eigenvariable.  Its premise is
to be derived in
\[
 w_z=\bigl(\Pi_{w'},b,z';\ \Gamma_{w'},\ \mathcal N_{b,\theta t},\
 \theta B(b,\epsilon_B(z'))\bigr).
\]
Put $\theta'=\theta[z\mapsto z']$.  The clauses of $\Gamma$ keep their
realizers, restricted to $w_z$.  By \cref{wit:lem:ns-base}(iv) and (i),
$w_z\Vdash b=\theta t$ and $w_z\Vdash\theta B(b,\epsilon_B(z'))$, and
$(\mathrm T_k)$ for the template $\theta B(q,\epsilon_B(z'))$ gives
$w_z\Vdash\theta B(\theta t,\epsilon_B(z'))$; this family realizes the source
assumption.  The induction hypothesis for the premise gives
$w_z\Vdash\theta B(\theta t,\Suc z')$, and $(\mathrm T_k)$ with the
symmetric equation $\theta t=b$ gives $w_z\Vdash\theta B(b,\Suc z')$.  Its
member at $w_z$ is the premise of the target \rulename{Close}, whose mode
is $\epsilon_{\theta B}=\epsilon_B$.  The conclusion $\D{x}{y}\theta B$ does
not contain $b$, and \rulename{Open} with the major premise
$\mathrm{Name}(\theta t)$ gives it at $w'$.

The induction is well founded as in \cref{wit:thm:eqback}:
\cref{wit:lem:ns-base,wit:thm:ns-rank} are constructed before and
independently of it, every call of the induction hypothesis concerns a
proper subderivation, and the realizers of the new assumptions in the cases
of \rulename{Alias} and \rulename{Close} are built from the actual clauses of
$w_c$ and $w_z$ without calling the translation.  The required derivation is
the member at $(\Pi;\Gamma)$ of the family for the identity map and the
realizers given by $(\mathrm S_k)$.
\end{proof}

\begin{corollary}[The naming calculus and $\iPRA$]
\label{wit:cor:ns-bridge}
With $\Gamma^{\circ}$ as in \cref{wit:cor:ab-bridge},
$\Pi;\Gamma\proves_{\mathsf{A}^{\mathrm n}_{\nu}}A$ if and only if
$\Gamma^{\circ}\proves_{\mathsf{K}}A$.  The ordinary realization of
$\mathsf{A}^{\mathrm n}_{\nu}$ is bi-interpretable and equiconsistent with
full first-order $\iPRA$, with whole numerical domains, actual equality, the
identity graph $x=y$ in both numerical composites, provable formula round
trips in both directions, and the contradiction $0=\Suc 0$.
\end{corollary}

\begin{proof}
By \cref{wit:thm:ns-complete}, $\mathsf{A}^{\mathrm n}_{\nu}$ derives exactly
the judgments of $\mathsf{A}_{\ast}$ on the same terms, formulas and clause
contexts, by effective translations; \cref{wit:cor:ab-bridge} gives the rest
unchanged.  The growing domains of \cref{wit:thm:ns-obstruction} play no role
here.
\end{proof}

\begin{remark}[\rulename{XBack} in place of \rulename{EqBack}]
\label{wit:rem:ns-xback}
The proofs of \cref{wit:lem:ns-base,wit:thm:ns-rank,wit:thm:ns-complete} use
\rulename{EqBack} only through \cref{wit:lem:eqb-atomic}(ii)--(iv).  The
construction of these statements for $\mathsf{A}^{\times}_{\nu}$ in the proof
of \cref{wit:lem:eqb-atomic} applies \rulename{Alias} only with the
equations $a(s)=d(s)$ and $t=s$ as conclusions.  Hence the calculus with
\rulename{XBack} in place of \rulename{EqBack}, \rulename{EqAlias} and
\rulename{Names} also derives the judgments of $\mathsf{A}_{\ast}$.  The
article draws this consequence; \pkg{names\_and\_scopes} states the theorem
for \rulename{EqBack} only.
\end{remark}

\begin{remark}[Comparison with the other naming results]
\label{wit:rem:ns-compare}
The scheme \rulename{Names} is the axiom scheme of $\Nat_{\mathrm v}^{+}$ in
\cref{wit:thm:named}, $\D{x}{y}\mathrm{Al}_a$ for every term $a$, which the
Lean module \Lean{NativeNamedWitness} uses for the calculus \Lean{DNPrf} with
variable witnesses and the native elimination (\Lean{dnprf\_iff\_nprf}).
There renaming, weakening and cut are admissible
(\cref{wit:lem:structural-0v}), transport is a primitive rule, and the name
serves only the introduction.  In $\mathsf{A}^{\mathrm n}_{\nu}$ none of these
is available: the same scheme also replaces the export of compound formulas
by \rulename{Alias} and the specialization of clauses at compound terms, and
the proof needs the persistent families and the induction on rank.
Compared with \cref{disch:thm:alias-rank}, the names are axioms instead of
introductions at the supplied witness $t$, and the target introduction of
$(\mathrm L_k)$ is in the mode $\epsilon_B$, so its zero row comes from
\cref{wit:lem:eqb-zero} when the matrix is diagonal.  The naming exception
of $\mathsf{G}_{\mathrm d}$ (\cref{wit:sec:guarded}) uses the naming matrix
$\mathrm{Nm}_r$ and a guarded calculus with a primitive transport.
\end{remark}

\begin{theorem}[Finitely many naming templates do not suffice]
\label{wit:thm:ns-templates}
Let $t_1,\dots,t_p$ be fixed raw terms, $t_j$ with free variables among
$x_1,\dots,x_{r_j}$.  Let $\mathsf{A}^{\mathrm e}_{\mathcal T}$ have the
rules \rulename{I0}, \rulename{IS}, \rulename{Open}, \rulename{EqAlias} and
the transport in every formula of \cref{wit:thm:ns-obstruction}, together
with
\begin{itemize}
\item \rulename{Use} at a parameter of $\Pi$ or at any numeral;
\item \rulename{Close} at a witness that is a parameter of $\Pi$ or a
 numeral, in the case mode for every matrix, and in the progress mode for
 every matrix that is a literal equation;
\item the axioms $\Pi;\Gamma\proves\mathrm{Name}\bigl(t_j(p_1,\dots,p_{r_j})\bigr)$
 for $j\le p$ and parameters $p_1,\dots,p_{r_j}$ of $\Pi$, not necessarily
 distinct.
\end{itemize}
Then there is a raw term $H(n,a)$ such that $\mathsf{A}^{\mathrm e}_{\mathcal T}$
does not derive
\[
 \varnothing;\varnothing\proves U_H,\qquad
 U_H\defeq\D{d}{z}\mathrm{Name}\bigl(H(z,z)\bigr),\quad d\text{ unused}.
\]
$\mathsf{A}^{\mathrm n}_{\nu}$ derives $U_H$, and $\iPRA$ proves its ordinary
reading.  Every rule of $\mathsf{A}^{\mathrm e}_{\nu}$ is a rule of
$\mathsf{A}^{\mathrm e}_{\mathcal T}$.
\end{theorem}

The instances of the templates are restricted to parameters.  With
arbitrary terms in place of the $p_i$ the single template $t(x)=x$ would
give every instance of \rulename{Names}.

\begin{proof}
The last statement holds as in \cref{wit:thm:ns-obstruction}; in the mode
$\epsilon_B$ the progress mode occurs only for diagonal matrices, which are
literal equations.  In $\mathsf{A}^{\mathrm n}_{\nu}$, \rulename{Close} in the
empty scope selects $0$ for the matrix $\mathrm{Name}(H(z,z))$, which is not
an equation, so the mode is $0$, and its premise
$z;\ \mathrm{Name}(H(0,0))\proves\mathrm{Name}(H(\Suc z,\Suc z))$ is an
instance of \rulename{Names}.  In $\iPRA$, $\mathrm{Name}(H(z,z))$ holds for
every $z$ with the witness $H(z,z)$, since $\sel(y,a,a)=a$ for all $y$ by
quantifier-free induction; hence the reading of $U_H$ holds.

\emph{The term $H$.}  Let $\bar t_j$ be the standard function of $t_j$ and
\[
 f(b)\defeq b+1+\sum_{j\le p}\max\bigl\{\bar t_j(a_1,\dots,a_{r_j}):
   a_1,\dots,a_{r_j}\le b\bigr\},
\]
with the single value for $r_j=0$.  The function $f$ is primitive recursive,
$f(b)>b$, $f$ is strictly increasing, and $\bar t_j(\vec a)<f(b)$ whenever
all $a_i\le b$.  By \cref{base:thm:primrec}, and the soundness of $\Teq$ in
$\mathbb{N}$, there is a raw term $F(u)$ whose standard function is $f$.  Put
$H(n,a)\defeq\It{n}{a}{u}{F(u)}$; then $H(n,a)\ge H(n',a)$ whenever $n\ge n'$,
and $H(k,a)$ is the $k$-fold iterate $f^k(a)$.

\emph{A nonstandard structure.}  Let $L^{*}$ have the order relation $<$ and,
for every raw term $t$ and every list $\vec x$ of distinct variables
containing its free variables, a function symbol interpreted in $\mathbb{N}$ as
the standard function of $t$ in the arguments $\vec x$.  Let $M$ be a proper
elementary extension of this standard $L^{*}$-structure, and let the value
$\llbracket t\rrbracket_\rho\in M$ of a term under an assignment $\rho$ in $M$
be the interpretation in $M$ of its symbol.  By elementarity:
\begin{enumerate}[label=\textup{(M\arabic*)}]
\item the computation equations, $\sel(y,a,a)=a$, and the identity
 $\llbracket t\subst{s}{b}\rrbracket_\rho=\llbracket t\rrbracket_{\rho[b\mapsto\llbracket s\rrbracket_\rho]}$
 hold in $M$, being universal sentences true in $\mathbb{N}$;
\item $<$ is a discrete linear order with least element $0$ in which every
 nonzero element is a successor, the numerals are its standard elements, and
 the properties of $f$, $H$ and $\bar t_j$ stated above, which are
 $L^{*}$-sentences, hold for the functions of $F$, $H$ and $t_j$ in $M$,
 written $f$, $H$ and $\bar t_j$ again;
\item for every equation $E(y)$ and every assignment, if $E$ holds at $y=0$
 and its truth at $y=j$ implies its truth at $y=j+1$ for all $j\in M$, then
 $E$ holds at every $j\in M$, because this instance of induction is true in
 $\mathbb{N}$.
\end{enumerate}
% ed.: the package takes a nonstandard model of PA and a raw term for f; the
% properties of that term in the model then require a PA proof that the term
% computes f.  An elementary extension of the standard structure with a
% symbol for every raw term gives (M1)--(M3) without that proof.
Fix a nonstandard $c\in M$ and put
\[
 J\defeq\{a\in M:\ a<f^k(c)\text{ for some standard }k\},
\]
with the iterates taken in the metatheory.  $J$ is an initial segment of
$M$; it contains every numeral and $c<f(c)$.  It is closed under successor
and under $f$: if $a<f^k(c)$, then $a+1\le f^k(c)<f^{k+1}(c)$ and
$f(a)<f^{k+1}(c)$.  It is closed under every $\bar t_j$: if all $a_i<f^k(c)$,
then $\bar t_j(\vec a)<f(f^k(c))=f^{k+1}(c)$.  But
$h\defeq\llbracket H(c,c)\rrbracket\notin J$: for every standard $k$, $k<c$,
so $f^k(c)=\llbracket H(\overline k,c)\rrbracket\le h$ by (M1) and (M2).

\emph{The evaluation.}  Take two worlds $J\le M$ with the active domains $J$
and $M$, evaluate all terms in $M$, also at the world $J$, and define forcing
and validity by the clauses of the proof of \cref{wit:thm:ns-obstruction},
with these worlds and domains.  Monotonicity, dependence on free variables
only, substitution (by (M1)) and the reading of assumptions hold as there.

\emph{Every rule of $\mathsf{A}^{\mathrm e}_{\mathcal T}$ preserves
validity.}  The computations hold by (M1), transport by substitution, and
\rulename{Use} and \rulename{Open} as before, since parameters take active
values and numerals are standard, hence in $J$.  For \rulename{EqAlias} move
to the world $M$, where the value of the defining term is active; the equation
obtained there pulls back, because its truth does not depend on the world.
\rulename{Close} in the case mode is sound as before, because both domains
contain $0$ and the predecessor of each of their nonzero elements ($J$ is an
initial segment of $M$, and every nonzero element of $M$ is a successor).
For \rulename{Close} in the progress mode the matrix is an equation $E(x,y)$;
let $a$ be the value of the witness, suppose that $E$ holds at $x=a$, $y=0$,
and that the context is forced at the current world.  The context is forced at
$M$, and for every $j\in M$ the assignment $z\mapsto j$ is active there; so
the premise turns the truth of $E$ at $y=j$ into its truth at $y=j+1$, for
every $j\in M$.  By (M3), $E$ holds at $x=a$ for every $y\in M$, and hence at
every later world and every active value.  This uses induction in $M$ for one
fixed equation, not induction over $J$ or over the forcing relation.  Finally
an axiom $\mathrm{Name}(t_j(\vec p))$ is forced: at $J$ the value
$v=\bar t_j(\rho\vec p)$ lies in $J$ by closure and witnesses the name,
because $\sel(y,v,v)=v$ for every $y\in M$ by (M1); at $M$ every value is
active.

\emph{$U_H$ is not forced at $J$.}  $\llbracket H(0,0)\rrbracket=0$ by (M1),
and $\mathrm{Name}(0)$ is forced at $J$ with the witness $0$.  So for every
witness $d\in J$ of $U_H$ the condition applies at the world $J$, and
$\mathrm{Name}(H(c,c))$, the instance at the active value $z\mapsto c$, would
have to be forced at $J$.  Its row at $y=0$ is true for every witness, and its
row at $y=1$ requires the witness to equal $h$, which is not in $J$.  Hence
$\varnothing;\varnothing\proves U_H$ is not valid and not derivable.
\end{proof}

\begin{remark}[Scope of the template theorem]
\label{wit:rem:ns-templates}
The theorem admits all numerals as targets and witnesses, and so covers
names of $0$ and of successors of parameters, and every fixed finite set of
primitive recursive operations.  It does not contradict
\cref{wit:thm:ns-complete}: the full scheme names the term $H(z,z)$ itself.
Its argument does not show soundness of the progress mode for compound
matrices in the two-world evaluation; that the matrix is an equation, whose
truth does not depend on the world, is essential to the step that uses (M3).
Like \cref{wit:thm:ns-obstruction} it separates the consequence relations of
the stated calculi and is not an obstruction to every interpretation of them.
\end{remark}

\begin{remark}[The package]
\label{wit:rem:ns-package}
\pkg{names\_and\_scopes} states the calculi as $\mathcal K_0$ and
$\mathcal K_1$, gives \cref{wit:lem:ns-shape}, the evaluation of
\cref{wit:thm:ns-obstruction} with the generalization to numerals up to $K$,
the proof of \cref{wit:thm:ns-complete} by the persistent families and the
induction on rank, which it credits to \pkg{equation\_aliases} and
\pkg{equational\_transport}, the bridge, and \cref{wit:thm:ns-templates}.  An
independent mathematical review is included.  None of the theorems is
formalized in Lean or Rocq.
\end{remark}

% ---------------------------------------------------------------------
\section{The raw guarded calculus}
\label{wit:sec:guarded}

This section draws on \pkg{deterministic\_introduction} for the naming
exception, and on \pkg{uniform\_raw} and the guarded companion of
\pkg{newest\_parameter} for the calculus without it.

\subsection{The calculus and its three introductions}

Terms are the raw terms of \cref{base:sec:terms}; no computation is built
into syntactic identity or into substitution.  A \emph{canonical formula} is
an equation or a formula $\D{x}{y}B$ in which $B$ is canonical and, if $B$
is an equation, it has the literal shape
\begin{equation}
 \It{y}{0}{u}{h}=\It{y}{0}{v}{k}.
 \label{wit:eq:guard}
\end{equation}
The bodies $h$ and $k$ may contain $x$, $y$ and parameters.  An equation at
the root of a formula is unrestricted.  Substitution of terms for
parameters preserves canonical formulas.  Contexts are lists of clauses
$[A;\,y.B(y)]$ of canonical formulas, with $y$ not free in $A$.  Put
$H\defeq(\It{0}{0}{u}{0}=0)$, an instance of \rulename{I0}; an assumption
$A$ is the clause $\mathsf{a}(A)\defeq[H;\,v.A]$ and is used by
\rulename{Use} with one computation.

The calculus $\mathsf{G}$ has the two computation equations as axioms in
every context and four inference rules:
\begin{gather*}
 \infrule{\Pi;\Gamma\proves A}{\Pi;\Gamma\proves B(t)}\ \rulename{Use}
 \ \bigl([A;\,y.B(y)]\in\Gamma\bigr)
 \qquad
 \infrule{\Pi;\Gamma\proves s=t\qquad\Pi;\Gamma\proves F(s)}
         {\Pi;\Gamma\proves F(t)}\ \rulename{Tr}
 \\[1.2ex]
 \infrule{\Pi;\Gamma\proves\D{x}{y}B\qquad
          \Pi,b;\ \Gamma,[B(b,0);\,y.B(b,y)]\proves A}
         {\Pi;\Gamma\proves A}\ \rulename{Open}
 \qquad
 \infrule{\Pi,z;\ \Gamma,\mathsf{a}(B(t,\eta_B(z)))\proves B(t,\Suc z)}
         {\Pi;\Gamma\proves\D{x}{y}B}\ \rulename{Intro}
\end{gather*}
with $\eta_B(z)=z$ if $B$ is an equation and $\eta_B(z)=0$ if $B$ is
compound.  In \rulename{Tr}, $F(v)$ is a canonical formula template.  The
mode of \rulename{Intro} depends on the root of the matrix and not on the
witness.  There is no primitive hypothesis rule, reflexivity, weakening,
cut, coverage or induction rule.  The ordinary reading of $\D{x}{y}B$ is
that of \cref{base:sec:reading}; the comparison of $\mathsf{G}$ with $\Nat$,
by padding equation matrices, and hence the bridge to $\iPRA$, is
established for $\mathsf{G}$ in \pkg{guarded\_core} and is inherited here.

\begin{definition}[Three selections of the witness]
\label{wit:def:guarded}
In $\mathsf{G}$ the witness $t$ of \rulename{Intro} is an arbitrary term
over $\Pi$.  In $\mathsf{G}_{\nu}$ it is $\nu(\Pi)$.  In
$\mathsf{G}_{\mathrm d}$ it is
\[
 \omega(B,\Pi)\defeq
 \begin{cases}
  r, & \text{$B$ is literally $\mathrm{Nm}_r$ with $r$ over $\Pi$,}\\
  \nu(\Pi), & \text{otherwise.}
 \end{cases}
\]
A match with a naming matrix has priority, and it is literal up to the names
of bound variables: $r$ contains neither $x$, $y$ nor the bound variables of
the two iterators.  It is not a match up to provable equality.
\end{definition}

Every naming matrix has the shape \eqref{wit:eq:guard}.  The derivations of
$\mathsf{G}_{\nu}$ and of $\mathsf{G}_{\mathrm d}$ are derivations of
$\mathsf{G}$ once the selected witnesses are written in.  In all three
calculi reflexivity and symmetry are derived from \rulename{I0} and
\rulename{Tr}: with $q=\It{0}{r}{u}{0}$, the equation $q=r$ used as both
premises of \rulename{Tr} in the template $v=r$ gives $r=r$, and from $s=t$
the transport of $s=s$ in the template $v=s$ gives $t=s$.

\subsection{The naming exception}

In $\mathsf{G}_{\mathrm d}$ the names are immediate, and the work is
structural: weakening to a larger scope can change the selected witness, and
the general cut needs the weakening of the inserted derivation.  The order
of the argument is a restricted cut, then weakening, then the general cut,
then completeness.

\begin{lemma}[Algebraic derivations and restricted cut]
\label{wit:lem:algebraic}
Call a derivation \emph{algebraic} if it uses only the computation axioms,
\rulename{Use} and \rulename{Tr}.
\begin{enumerate}[label=\textup{(\roman*)}]
\item An algebraic derivation of $\Pi;\Gamma\proves A$ yields an algebraic
 derivation of $\Pi';\Delta\proves\rho A$ for every injective map $\rho$
 from $\Pi$ into $\Pi'$ and every clause context $\Delta$ over $\Pi'$ into
 which $\rho\Gamma$ is embedded.
\item If $\Pi;\Gamma\proves A$ has an algebraic derivation and
 $\Pi;\Gamma,\mathsf{a}(A)\proves F$ in $\mathsf{G}_{\mathrm d}$, then
 $\Pi;\Gamma\proves F$ in $\mathsf{G}_{\mathrm d}$.
\end{enumerate}
\end{lemma}

\begin{proof}
(i) By induction on the derivation; its rules have no choice that depends
on the scope and no eigenvariable.  (ii) By induction on the derivation of
$F$.  A \rulename{Use} of the clause $\mathsf{a}(A)$ is replaced by the
algebraic derivation of $A$, moved to the scope and context of that leaf by
(i).  At \rulename{Intro} and \rulename{Open} the induction continues in
the extended branch.  Deleting a clause does not change the scope, so every
\rulename{Intro} selects the same witness before and after.  No weakening
of a derivation that contains an \rulename{Intro} is used.
\end{proof}

\begin{lemma}[Every term has a name]
\label{wit:lem:gd-name}
For every term $r$ over $\Pi$, $\mathsf{G}_{\mathrm d}$ derives
$\Pi;\Gamma\proves\D{x}{y}\mathrm{Nm}_r$.  If a clause context contains
$[\mathrm{Nm}_r(b,0);\,y.\mathrm{Nm}_r(b,y)]$, then $b=r$ and $r=b$ have
algebraic derivations from it.
\end{lemma}

\begin{proof}
The rule selects $r$ by the naming alternative.  The matrix is an equation,
so the premise asks for $\mathrm{Nm}_r(r,\Suc z)$ from
$\mathrm{Nm}_r(r,z)$; the conclusion is literally
$\sel(\Suc z,0,r)=\sel(\Suc z,0,r)$ and is a reflexivity.  For the second
statement, both sides of $\mathrm{Nm}_r(b,0)$ equal $0$ by \rulename{I0}, so
the antecedent of the clause is derivable; \rulename{Use} at $1$ gives
$\sel(1,0,b)=\sel(1,0,r)$, and \rulename{IS} with \rulename{Tr} gives $b=r$;
symmetry gives $r=b$.
\end{proof}

\begin{theorem}[Weakening by naming the old selection anew]
\label{wit:thm:weakening}
Let $\Pi;\Gamma\proves A$ be derivable in $\mathsf{G}_{\mathrm d}$.  For
every injective map $\rho$ from $\Pi$ into a scope $\Pi'$ and every clause
context $\Delta$ over $\Pi'$ into which $\rho\Gamma$ is embedded,
$\Pi';\Delta\proves\rho A$ is derivable in $\mathsf{G}_{\mathrm d}$, by an
effective construction.
\end{theorem}

\begin{proof}
By induction on the height of the given derivation, for all $\rho$ and
$\Delta$.  The computation axioms, \rulename{Use} and \rulename{Tr} are
immediate.  At \rulename{Open}, choose a fresh parameter of the receiving
scope, extend $\rho$ and the embedding of clauses, and apply the induction
hypothesis to the two premises.

Consider an \rulename{Intro} with the matrix $B$ and the selected witness
$w$, and let $B'=\rho B$, $w'=\rho w$.  A match with a naming matrix is
preserved and reflected by $\rho$: a bound variable does not become a
parameter or conversely, and a parameter does not become a numeral or a
compound term.  If $B=\mathrm{Nm}_r$, the witness selected in the receiving
scope is $\rho r=w'$, and the induction hypothesis for the premise gives the
rule directly.  The same holds if $B$ is not a naming matrix and
$\nu(\Pi')=w'$.

Suppose $B'$ is not a naming matrix and $\nu(\Pi')\neq w'$.  Derive
$\D{x}{y}\mathrm{Nm}_{w'}$ by \cref{wit:lem:gd-name} and open it at a fresh
parameter $e$, which is then the newest; its clause gives algebraic
derivations of $e=w'$ and $w'=e$.  The formula $\D{x}{y}B'$ will be
introduced in the scope $\Pi',e$, where the rule selects $e$.  Let $v$ be
the eigenvariable of that introduction and put
\[
 A_{\mathrm{old}}=B'(w',\eta_{B'}(v)),\qquad
 A_{\mathrm{new}}=B'(e,\eta_{B'}(v)).
\]
Apply the induction hypothesis to the premise of the given
\rulename{Intro}, with $\rho$ extended by sending its eigenvariable to $v$,
and with the receiving context consisting of $\Delta$, the clause of the
name, $\mathsf{a}(A_{\mathrm{old}})$ and $\mathsf{a}(A_{\mathrm{new}})$.
This is legitimate although $e$ and $v$ are new parameters of the receiving
scope, because the induction statement quantifies over all injective maps of
the smaller derivation.  The result derives $B'(w',\Suc v)$.  From
$A_{\mathrm{new}}$, \rulename{Tr} along $e=w'$ gives $A_{\mathrm{old}}$;
this derivation is algebraic.  \Cref{wit:lem:algebraic}(ii) removes the
clause $\mathsf{a}(A_{\mathrm{old}})$.  \rulename{Tr} along $w'=e$ gives
$B'(e,\Suc v)$.  Apply \rulename{Intro}, which selects $e$, and discharge
the name by \rulename{Open}; the conclusion does not contain $e$.  If the
old selection was $0$ in an empty scope, the name is that of $0$.

Every use of the induction hypothesis is for a proper subderivation of the
given derivation.  The naming and transport derivations built here are
never fed back into the induction, and the restricted cut was established
before.  The construction is therefore well founded and uses neither the
general cut nor weakening.
\end{proof}

\begin{corollary}[Cut of an assumption]
\label{wit:cor:gd-cut}
If $\Pi;\Gamma\proves A$ and $\Pi;\Gamma,\mathsf{a}(A)\proves F$ in
$\mathsf{G}_{\mathrm d}$, then $\Pi;\Gamma\proves F$ in
$\mathsf{G}_{\mathrm d}$.
\end{corollary}

\begin{proof}
By induction on the derivation of $F$.  A \rulename{Use} of
$\mathsf{a}(A)$ is replaced by the derivation of $A$, moved to the scope and
context of that leaf by \cref{wit:thm:weakening}.  Deleting a clause leaves
the scope unchanged, so the introductions of the receiving derivation keep
their selections.
\end{proof}

\begin{theorem}[The introduction with a naming exception is complete]
\label{wit:thm:gd}
For every scope $\Pi$, canonical clause context $\Gamma$ and canonical
formula $A$,
\[
 \Pi;\Gamma\proves_{\mathsf{G}_{\mathrm d}}A
 \quad\Longleftrightarrow\quad
 \Pi;\Gamma\proves_{\mathsf{G}}A .
\]
\end{theorem}

\begin{proof}
From left to right the selected witness is a witness of \rulename{Intro} in
$\mathsf{G}$.  For the converse induct on the derivation in $\mathsf{G}$.
At an \rulename{Intro} whose matrix is a naming matrix $\mathrm{Nm}_r$,
discard the supplied witness and the premise and use
\cref{wit:lem:gd-name}.  At any other \rulename{Intro}, with the witness
$t$, translate the premise by the induction hypothesis, derive
$\D{x}{y}\mathrm{Nm}_t$ and open it at a fresh $b$.  In the context with the
opened name and the assumption $B(b,\eta_B(z))$, \rulename{Tr} along $b=t$
gives $B(t,\eta_B(z))$.  \Cref{wit:thm:weakening,wit:cor:gd-cut} insert the
translated premise and give $B(t,\Suc z)$, and \rulename{Tr} along $t=b$
gives $B(b,\Suc z)$.  The newest parameter of the scope is $b$ and $B$ is
not a naming matrix, so the rule selects $b$.  Close the introduction and
the opening of the name.
\end{proof}

\begin{remark}[Why the qualifications matter]
\label{wit:rem:gd-qualifications}
The simulation changes the witness from $t$ to $b$ and must not change the
assumption of the premise.  In $\mathsf{G}$ the mode depends on the root of
the matrix only.  A rule that selects progress exactly when the witness is
the literal $0$ does not fit this argument as it stands: at $t=0$ it has the
assumption $B(0,z)$ and at $b$ the assumption $B(b,0)$; the calculus
$\mathsf{R}$ of \cref{wit:sec:R} is treated by recovering the two modes
separately.  The diagonal selection of \cref{wit:sec:diagonal} is
independent of the witness as well.  The recognition of a naming matrix is
a fixed syntactic match.  Replacing it by the condition that some instance
of the matrix names the witness would need additional data or a search, and
a rule about a provable naming instance would reintroduce a premise.  The
proof uses transport in arbitrary canonical formulas: the donor of the
restricted cut converts $B(e,\eta)$, which may be compound.  Through the
comparison, substitution is admissible in $\mathsf{G}_{\mathrm d}$ as a
transformation of derivable judgments; a certificate is not claimed to
remain valid after a substitution.
\end{remark}

\subsection{Uniform introduction: the stable fragment}

In $\mathsf{G}_{\nu}$ the name $\D{x}{y}\mathrm{Nm}_r$ is no longer
immediate, and the names are generated as in \cref{wit:sec:packet}.  The
guarded calculus has no rule for equation induction, so the predecessor has
to be obtained first, in a part of the calculus that is insensitive to the
selected witness.

\begin{definition}
\label{wit:def:guarded-stable}
An \rulename{Intro} of $\mathsf{G}_{\nu}$ with the matrix $B(x,y)$ is
\emph{dummy} if $x$ is not free in $B$; a derivation is \emph{stable} if all
its introductions are dummy.
\end{definition}

\begin{lemma}[Stability]
\label{wit:lem:guarded-stable}
Stable derivations of $\mathsf{G}_{\nu}$ are closed under the substitution
of terms over a scope $\Pi'$ for the parameters, with any clause context
over $\Pi'$ that contains the substituted clauses, and under the cut of an
assumption by a stable derivation.
\end{lemma}

\begin{proof}
As \cref{wit:lem:stable}.  The premise of a dummy introduction does not
depend on the witness; a substitution does not change whether the matrix is
an equation, and it preserves the shape \eqref{wit:eq:guard}.
\end{proof}

For an equation $E$ let the \emph{padding} be
$L(E)\defeq\D{r}{v}\gamma_v(E)$ with $r,v$ not free in $E$; it is canonical.
For equations $U\equiv(s=t)$ and $V\equiv(p=q)$ let
\[
 U\Rightarrow_{g}V\defeq\D{r}{w}L\bigl(\sel(w,s,p)=\sel(w,t,q)\bigr),
\]
a canonical formula with a compound matrix and an unused coordinate $r$.

\begin{lemma}[Padding, implication and cases without names]
\label{wit:lem:guarded-tools}
The following hold for stable derivations of $\mathsf{G}_{\nu}$, in every
scope and context.
\begin{enumerate}[label=\textup{(\roman*)}]
\item $E$ and $L(E)$ are interderivable.
\item If $\Pi;\Gamma,\mathsf{a}(U)\proves V$, then
 $\Pi;\Gamma\proves U\Rightarrow_{g}V$; and $U\Rightarrow_{g}V$ together
 with $U$ gives $V$.
\item If $\Pi;\Gamma\proves F(0)$ and $\Pi,z;\Gamma\proves F(\Suc z)$,
 then $\Pi;\Gamma\proves F(t)$ for every term $t$; here $F(y)$ is a
 compound formula or an equation.
\item If $E$ is an equation and $t$ a term over $\Pi$, and
 $\Pi;\Gamma,\mathsf{a}(t=0)\proves E$ and
 $\Pi,z;\ \Gamma,\mathsf{a}(t=\Suc z)\proves E$, then
 $\Pi;\Gamma\proves E$.
\end{enumerate}
\end{lemma}

\begin{proof}
(i) From $E$: introduce $L(E)$.  The matrix $\gamma_v(E)$ is an equation,
so the premise asks for $\gamma_{\Suc z}(E)$ from $\gamma_z(E)$; it follows
from $E$ by \rulename{IS} and \rulename{Tr}, and the assumption is unused.
From $L(E)$: open it at a fresh $b$; the antecedent $\gamma_0(E)$ of the
clause holds by \rulename{I0}, \rulename{Use} at $1$ gives $\gamma_1(E)$,
and \rulename{IS} gives $E$, which does not contain $b$.

(ii) Write $\mathrm{row}_w$ for $\sel(w,s,p)=\sel(w,t,q)$.  The matrix is
compound, so the introduction is by cases and asks for
$L(\mathrm{row}_{\Suc z})$ from $L(\mathrm{row}_0)$.  By (i) and
\rulename{I0} the assumption gives $U$; the cut of the given derivation
gives $V$; \rulename{IS} gives $\mathrm{row}_{\Suc z}$, and (i) pads it.
For the elimination open the formula at a fresh $b$, derive
$L(\mathrm{row}_0)$ from $U$ by \rulename{I0} and (i), use the clause at
$1$, and unpad and compute to obtain $V$.

(iii) For a compound $F$, introduce $\D{r}{y}F(y)$ with $r$ unused; it is
by cases, and its premise is the given derivation of $F(\Suc z)$ weakened by
the assumption $F(0)$.  Open it at a fresh $b$ and use the clause
$[F(0);\,y.F(y)]$ at $t$, the antecedent being the given derivation moved to
the scope $\Pi,b$ by \cref{wit:lem:guarded-stable}.  For an equation $E(y)$
apply this to $L(E(y))$ and use (i).

(iv) Let $F(y)$ be $(t=y)\Rightarrow_{g}E$.  By (ii) the hypotheses give
$F(0)$ and $F(\Suc z)$, by (iii) $F(t)$, and by (ii) with the reflexivity
$t=t$ the equation $E$.
\end{proof}

\begin{lemma}[Equation induction in the stable fragment]
\label{wit:lem:guarded-ei}
Let $E(y)\equiv(s(y)=t(y))$.  If $\Pi;\Gamma\proves E(0)$ and
$\Pi,z;\ \Gamma,\mathsf{a}(E(z))\proves E(\Suc z)$ have stable derivations
in $\mathsf{G}_{\nu}$, then so has $\Pi;\Gamma\proves E(a)$ for every term
$a$ over $\Pi$.  Consequently every derivation of the equational fragment
$\mathcal E$ is translated into a stable derivation of $\mathsf{G}_{\nu}$,
and for all terms $r,s$
\[
 \Pi;\Gamma\proves\operatorname{pred}(\Suc r)=r,
 \qquad
 \Pi;\ \Gamma,\mathsf{a}(\Suc r=\Suc s)\proves r=s
\]
by stable derivations.
\end{lemma}

\begin{proof}
Let $G(y)\defeq\gamma_y(E(y))$, an equation of the shape
\eqref{wit:eq:guard}.  First, $E(z)$ is derivable from $G(z)$: by
\cref{wit:lem:guarded-tools}(iv) with the term $z$, in the branch $z=0$
transport the base $E(0)$ along $0=z$, and in the branch $z=\Suc w$
transport $G(z)$ to $G(\Suc w)$, compute the two gates to obtain
$E(\Suc w)$, and transport back.  The step then gives $E(\Suc z)$, by
\cref{wit:lem:guarded-stable}, and computation gives $G(\Suc z)$.  This is
the premise of the introduction of $\D{r}{y}G(y)$ with $r$ unused, which is
by progress.  Open it at a fresh $b$; the antecedent $G(0)$ of its clause
holds by computation, and \rulename{Use} at $a$ gives $G(a)$.  A second
application of \cref{wit:lem:guarded-tools}(iv), with the term $a$ and with
$E(0)$ in the branch $a=0$, gives $E(a)$.  All introductions are dummy.

For the translation of $\mathcal E$, an assumption $E$ is the clause
$\mathsf{a}(E)$, the equality rules are instances of \rulename{Tr} and
reflexivity, and induction is the first statement; the instances of earlier
lemmas that an equational derivation uses are available by
\cref{wit:lem:guarded-stable}.  The last statement follows from
\cref{base:thm:pred}.
\end{proof}

\subsection{Uniform introduction: names and completeness}

Hereditary derivability $\Pi;\Gamma\Vdash A$ is defined for
$\mathsf{G}_{\nu}$ as in \cref{wit:def:hereditary}, with clause contexts.
The computation axioms, \rulename{Use} and \rulename{Tr} hold with
$\Vdash$; \rulename{Open} holds with $\Vdash$ by allocating its parameter in
each receiving scope; a stable derivation is a member of a hereditary
family; and, as in \cref{wit:lem:uniform-cut},
\begin{equation}
 \Pi;\Gamma\Vdash A\ \text{ and }\ \Pi;\Gamma,\mathsf{a}(A)\proves B
 \quad\text{imply}\quad\Pi;\Gamma\proves B .
 \label{wit:eq:guarded-cut}
\end{equation}
In the proof of \eqref{wit:eq:guarded-cut} the donor is generated at each
\rulename{Use} of the removed clause in the scope and context of that leaf,
the other clause addresses are adjusted, and every other node keeps its
scope and therefore its selected witness.  The hypotheses of the donor must
remain available after the deletion.

The packet $X(y)$ of \cref{wit:def:packet} is canonical: both sides of its
matrix are iterators with the counter $v$ and the seed $0$.

\begin{lemma}[The packet in the guarded calculus]
\label{wit:lem:guarded-packet}
In $\mathsf{G}_{\nu}$, for every scope $\Pi$, clause context $\Gamma$ and
term $t$ over $\Pi$: $\Pi;\Gamma\proves\D{a}{r}X(r)$ with $a$ unused, and
$\Pi;\Gamma\proves X(\Suc t)$.
\end{lemma}

\begin{proof}
$X(0)$ is introduced at whatever witness the scope selects.  Its matrix is
an equation, the introduction is by progress, and the instance asked for
computes to $0=0$ whatever the witness and the assumption are.  The matrix
of $\D{a}{r}X(r)$ is compound, so its introduction is by cases and asks for
$X(\Suc z)$ with a new eigenvariable $z$.  In that branch the scope ends in
$z$, so the introduction of $X(\Suc z)$ selects $z$, and the instance asked
for computes to $\Suc z=\Suc z$.  The assumption $X(0)$ of the outer
introduction is unused.  Open $\D{a}{r}X(r)$ at a fresh parameter; its
clause is $[X(0);\,r.X(r)]$.  Derive $X(0)$ again in the extended scope and
use the clause at $\Suc t$.  The result does not contain the opened
parameter, and \rulename{Open} discharges it.
\end{proof}

The lemma is about the whole construction, which allocates $z$ anew in each
receiving scope.  A separate derivation of $X(\Suc z)$ is not claimed to
remain a derivation after $z$ ceases to be the newest parameter.

\begin{lemma}[The alias in the guarded calculus]
\label{wit:lem:guarded-alias}
If a clause context over a scope contains
$[M_{\Suc t}(b,0);\,v.M_{\Suc t}(b,v)]$, then $b=t$ is hereditarily
derivable from it.
\end{lemma}

\begin{proof}
The antecedent holds by computation.  \rulename{Use} at $1$ and the
successor computations give $\Suc b=\Suc t$, and
\cref{wit:lem:guarded-ei} gives $b=t$.  The derivation is stable.
\end{proof}

\begin{theorem}[Uniform introduction on raw terms]
\label{wit:thm:guarded-uniform}
For every scope $\Pi$, canonical clause context $\Gamma$ and canonical
formula $A$,
\[
 \Pi;\Gamma\proves_{\mathsf{G}_{\nu}}A
 \quad\Longleftrightarrow\quad
 \Pi;\Gamma\proves_{\mathsf{G}}A
 \quad\Longleftrightarrow\quad
 \Pi;\Gamma\proves_{\mathsf{G}_{\mathrm d}}A .
\]
The naming exception of $\mathsf{G}_{\mathrm d}$ is not needed.
\end{theorem}

\begin{proof}
The second equivalence is \cref{wit:thm:gd}, and
$\mathsf{G}_{\nu}\subseteq\mathsf{G}$ is immediate.  For
$\mathsf{G}\subseteq\mathsf{G}_{\nu}$ it is shown by induction on a
derivation of $\mathsf{G}$ that $\Pi;\Gamma\Vdash A$.  The computation
axioms, \rulename{Use} and \rulename{Tr} are copied with their terms
renamed.  For \rulename{Open}, allocate a parameter in the receiving scope
and apply the induction hypothesis to the branch with the extended map.

Consider an \rulename{Intro} at a term $t$ with the matrix $B$.  The
induction hypothesis is
$\Pi,z;\ \Gamma,\mathsf{a}(B(t,\eta_B(z)))\Vdash B(t,\Suc z)$.  Let
$(\Pi',\rho,\Delta)$ be given and write $t'$, $B'$ for the images.  Derive
$X(\Suc t')$ by \cref{wit:lem:guarded-packet} and open it at a fresh $b$.
Introduce $\D{x}{y}B'$ in the scope $\Pi',b$; the rule selects $b$ and then
allocates its eigenvariable $z'$.  Let $\Theta$ be the scope $\Pi',b,z'$
with the context consisting of $\Delta$, the clause of the packet and
$\mathsf{a}(B'(b,\eta_{B'}(z')))$.  By \cref{wit:lem:guarded-alias},
$\Theta\Vdash b=t'$, and \rulename{Tr} gives
$\Theta\Vdash B'(t',\eta_{B'}(z'))$.  The induction hypothesis at the map
$\rho[z\mapsto z']$, with \eqref{wit:eq:guarded-cut}, gives
$\Theta\Vdash B'(t',\Suc z')$, and \rulename{Tr} along $t'=b$ gives
$\Theta\Vdash B'(b,\Suc z')$.  Its member at the identity is the premise of
the introduction.  The conclusion does not contain $b$, and \rulename{Open}
discharges it.  The two modes agree because $\eta_B$ depends on the root of
the matrix, which neither $\rho$ nor the change of witness alters.  All uses
of the induction hypothesis are for proper subderivations, and the packet
and the alias were constructed before.
\end{proof}

\begin{corollary}[The literal naming theorem is derived]
\label{wit:cor:guarded-name}
For every term $t$, the formula $\D{x}{y}\mathrm{Nm}_t$ has a hereditary
derivation in $\mathsf{G}_{\nu}$, without a selection by pattern.
\end{corollary}

\begin{proof}
Derive $X(\Suc t)$, open it at a fresh $b$ and extract $b=t$.  Introduce
$\D{x}{y}\mathrm{Nm}_t$; the selected witness is $b$.  In the premise the
alias and the selector computations give $\mathrm{Nm}_t(b,\Suc z)$, the
progress assumption being unused.  Discharge the packet.
\end{proof}

By \cref{wit:thm:guarded-uniform} and the inherited comparison of
$\mathsf{G}$ with $\Nat$, the ordinary realizations of $\mathsf{G}_{\nu}$
and of $\mathsf{G}_{\mathrm d}$ are bi-interpretable and equiconsistent with
$\iPRA$, with whole domains, actual equality and the identity graph in both
numerical composites; weakening, renaming, cut and substitution are
admissible in $\mathsf{G}_{\nu}$ as consequences of the theorem.

\begin{remark}[Versions]
\label{wit:rem:guarded-versions}
The stable fragment and the route to equation induction through
\cref{wit:lem:guarded-tools}(iv) are those of the guarded companion of
\pkg{newest\_parameter}, which states the translation with persistent
builders and an arbitrary substitution, as in \cref{wit:rem:families}.
\pkg{uniform\_raw} proves instead, by induction on a derivation of
$\mathcal E$ and for all substitutions $\sigma$ at once, that
$\Gamma\sigma\Vdash E\sigma$.  At an induction with the invariant $E(y)$ it
regenerates the step at $y=0$ and at $y=\Suc w$ from the source
derivation, proves $\Phi(y)\defeq G(y)\Rightarrow_{g}E(\Suc y)$ at $0$ and
at $\Suc w$, obtains $\Phi(z)$ by cases, and concludes as above; no
derivation of the target is substituted into.  A technical note of
\pkg{uniform\_raw} also observes that, once the names are available, the
proof of \cref{wit:thm:weakening} applies to $\mathsf{G}_{\nu}$ with no
case for a naming matrix, which gives weakening and then cut directly.
The checkers of \pkg{uniform\_raw} and of \pkg{newest\_parameter} for
$\mathsf{G}_{\nu}$ are obtained from the checker of
$\mathsf{G}_{\mathrm d}$ by deleting the function that selects the witness
(\cref{wit:sec:sizes}).
\end{remark}

% ---------------------------------------------------------------------
\section{A recursor whose predecessor computes}
\label{wit:sec:recursor}

This section draws on \pkg{parameter\_intro}; the normalization lemma and the
extensionality lemma are also in \pkg{conditional\_copy}.  The term
language, its normal forms, the decoder and the ordinary theory with its
bridge to $\iPRA$ are those of \cref{ctr:def:languages,ctr:lem:nf,ctr:lem:decoder,ctr:thm:transfer};
the section states what the witness-free introduction adds to them.  The
calculi on the same terms with a supplied witness are those of
\cref{ctr:sec:three}.

On iterator terms no term $p$ satisfies $p(\Suc z)=z$ by computation, and the
names of \cref{wit:sec:packet} need the derived predecessor and an equality
rule.  With a recursor whose body receives the counter, and with terms taken
in normal form, the predecessor computes, and the name of every term is
obtained with clause specialization alone.  This makes the restricted
introduction available in a calculus of three rule families in which
equality transport is itself part of the discharge rule and is not available
beforehand.

\subsection{Terms in normal form}

The raw terms are those of the counter-aware seeded recursor of
\cref{ctr:def:languages}(2): variables, $0$ and $R(n,a;\,i,u.\,h)$, with
$i,u$ bound in $h$ only, the successor $\Suc a\defeq R(0,a;\,i,u.\,0)$ and
the computation equations \eqref{ctr:eq:r2zero} and \eqref{ctr:eq:r2succ}.
The \emph{public terms} are the normal terms of \cref{ctr:def:normal}, and
every substitution performed by a rule is followed by normalization $\nf$
(\cref{ctr:def:eval}).  Computation is thus part of the syntax, and the
normalizer is part of a checker and of its measured size; it is an external
operation on finite terms, with no primitive recursive bound uniform in the
term (\cref{ctr:rem:fuel}).  By \cref{ctr:lem:nf}, normal forms are unique,
normalization respects both computation equations in every context, and
\[
 \nf(t[\sigma])=\nf\bigl(\nf(t)[\nf\circ\sigma]\bigr)
\]
for terms, formulas and templates with one distinguished variable
\eqref{ctr:eq:SC}.  Normalization introduces no free variable: by
induction on $t$, the free variables of $\operatorname{ev}(t,\rho)$ are
among those of the values $\rho(x)$ for $x$ free in $t$.

\begin{remark}[The normalization lemma in the packages]
\label{wit:rem:nf}
\pkg{parameter\_intro} and \pkg{conditional\_copy} define the normal forms
by rewriting, taking $0$, $\Suc$ and $R$ as distinct raw constructors (to
avoid the circular rewrite of $R(0,a;\,i,u.\,0)$ to itself), orienting the
two equations from left to right, reducing in every subterm including
bodies, and printing the normal constructor $\Suc a$ as
$R(0,a;\,i,u.\,0)$.  They prove strong normalization and confluence by the
reducibility argument for the base-type recursor of G\"odel's system $T$
recorded in \cref{ctr:rem:nfvariants}; the resulting normal forms are those
of \cref{ctr:def:eval}.
\end{remark}

The abbreviations are those of \eqref{ctr:eq:decoder}: the decoder $P$, the
selector $C$, the predecessor $p$, and $J(n,a;\,i,u.\,h)$ for
$\operatorname{Rec}(n,a;\,i,u.\,h)$.  By \cref{ctr:lem:decoder}, on normal
terms
\begin{gather*}
 P(0)\equiv 1,\quad P(\Suc n)\equiv n,\qquad
 C(0,a,b)\equiv a,\quad C(\Suc n,a,b)\equiv b,\qquad
 p(0)\equiv 0,\quad p(\Suc n)\equiv n,\\
 J(0,a;\,i,u.\,h)\equiv a,\qquad
 J(\Suc n,a;\,i,u.\,h)\equiv h[n/i,\ J(n,a;\,i,u.\,h)/u],
\end{gather*}
where $\equiv$ is literal identity of normal forms.  The iterator of
\cref{base:sec:terms} is $J$ with the counter coordinate unused
(\cref{ctr:thm:compiler}).  In this section and the next, $C$ takes the
place of $\sel$ in the gate and in the naming matrix of
\cref{wit:def:naming}.

\subsection{Three rule families}

Formulas are equations and $\D{x}{y}B$.  A clause is $[\ ;\,y.B]$ or
$[A;\,y.B]$ with $y$ not free in $A$; the first reads $\forall y\,B$.  Every
derivation starts from the one clause
\[
 \mathsf{Rf}\defeq[\ ;\,y.\,y=y],
\]
and an assumption $A$ is the clause $[\ ;\,v.A]$.  For a matrix $B$ and a
term $r$ write $B_r(y)$ for $\nf(B(r,y))$.  The calculus
$\mathsf{P}_{\nu}$ has three rule families.
\begin{enumerate}[label=\textup{(\roman*)}]
\item \rulename{Use}: for a clause of the context and a term $t$ over
 $\Pi$, derive the antecedent of the clause, if it has one, and conclude
 the normalized body at $t$.
\item \rulename{Intro}: with $\ell=\nu(\Pi)$ and a fresh $z$,
\[
 \infrule{\Pi,z;\ \Gamma,[\ ;\,v.B_\ell(\eta)]\proves B_\ell(\Suc z)}
         {\Pi;\Gamma\proves\D{x}{y}B}
 \qquad
 \eta=\begin{cases}
  z,&\text{$B_\ell(0)$ is literally an equation $r=r$,}\\
  0,&\text{otherwise.}
 \end{cases}
\]
 The instruction supplies no witness, no mode and no proof of the instance
 at $0$.
\item \rulename{Discharge}: with a fresh $b$, the clause
 $\mathcal{C}_B(b)\defeq[\nf(B(b,0));\,y.\nf(B(b,y))]$ and a template
 $F(v)$,
\[
 \infrule{\Pi;\Gamma\proves\D{x}{y}B\qquad
          \Pi,b;\ \Gamma,\mathcal{C}_B(b)\proves s=t\qquad
          \Pi,b;\ \Gamma,\mathcal{C}_B(b)\proves\nf(F(s))}
         {\Pi;\Gamma\proves\nf(F(t))}
\]
 provided $\nf(F(t))$ does not contain $b$.  The terms $s$, $t$ and the
 template may contain $b$.
\end{enumerate}
There are no other logical, equality, arithmetic or structural rules; the
computations \eqref{ctr:eq:r2zero} and \eqref{ctr:eq:r2succ} belong to the syntax.  The designated
contradiction is $0=1$.

\begin{proposition}[Soundness]
\label{wit:prop:P-sound}
Every rule of $\mathsf{P}_{\nu}$ is valid under the ordinary reading in
intuitionistic logic with equality, the equations \eqref{ctr:eq:r2zero} and \eqref{ctr:eq:r2succ},
numerical coverage, $0=1\to\bot$ and induction for literal equations.
\end{proposition}

\begin{proof}
\rulename{Use} is the application of a clause.  \rulename{Discharge} is the
elimination of the existential quantifier followed by the substitution of
equals in the branch; the condition on $\nf(F(t))$ is the eigenvariable
condition, and the rule is sound also when $t$ contains $b$.  For
\rulename{Intro}: if the instance at $0$ is not reflexive, the mode is by
cases, and coverage combines the retained instance at $0$ with the derived
instance at successors.  If it is reflexive, the matrix is a literal
equation, because normalization does not remove a binder of a formula; the
instance at $0$ is provable and induction for that equation gives all
instances.  In both cases $\ell$ is a witness of the existential
quantifier.  No induction on a quantified formula is used.
\end{proof}

A premise by cases can be used at any witness.  If the rule selects progress
instead, the discarded assumption at $0$ is literally reflexive: derive it
by \rulename{Use} of $\mathsf{Rf}$ and ignore the progress assumption.  In
the constructions below this is an operation on \emph{builders}, that is, on
constructions that can be carried out again after fresh parameters and
clauses have been adjoined; it is not a weakening principle for derivation
trees.  A builder is called \emph{persistent}.

\subsection{Constructions that precede equality}

Write $q_yA$ for $\D{w}{y}A$ with $w$ not free in $A$.  The witness of its
introduction is unused, so its instances do not depend on which parameter is
the newest.  The elimination of $\D{x}{y}B$ without substitution is
\rulename{Discharge} with $s=t=0$, the equation $0=0$ by \rulename{Use} of
$\mathsf{Rf}$, and the constant template whose value is the conclusion.
Hence the introduction by cases and the elimination of $q$ are available
from the three rules, with the repair of the previous paragraph when the
instance at $0$ is reflexive; this gives numerical cases for all formulas,
with the branches constructed in their own scopes.  Neither equality
transport nor equation induction is needed for these constructions.

For a formula $A$ let $G_vA$ be $A$ with each side $s$ of its final equation
replaced by $C(v,0,s)$, the binders being kept.  Then
$G_{\Suc z}A\equiv A$, and $G_0A$ is a tower of binders over $0=0$, which is
derivable in every scope by \rulename{Intro} and reflexivity, its witnesses
being unused.  For persistent premises the padding $L(A)\defeq q_yG_yA$ is
interderivable with $A$: one direction ignores its assumption at $0$ and
uses $A$, the other uses the closure at $1$ with the provable instance at
$0$.  Padding two formulas to the same number of binders permits their
selector, sidewise by $C(v,\cdot,\cdot)$, and the implication
$q_y\operatorname{Sel}(y,A,B)$ with its introduction and elimination.
Every introduction in these constructions has an unused witness, except
those of the towers over $0=0$, whose result does not depend on the witness.

\begin{lemma}[A name without equality]
\label{wit:lem:name-recursor}
For every term $t$ over $\Pi$ there is a persistent derivation in
$\mathsf{P}_{\nu}$ of
\[
 \operatorname{Name}(t)\defeq\D{x}{u}\bigl(C(u,0,x)=C(u,0,t)\bigr),
\]
and the clause obtained by opening $\operatorname{Name}(t)$ at a fresh $b$
gives $b=t$ by one \rulename{Use}.
\end{lemma}

\begin{proof}
The opened clause is $[0=0;\,u.\,C(u,0,b)=C(u,0,t)]$; \rulename{Use} at $1$,
with the antecedent from $\mathsf{Rf}$, gives $b=t$.  For the derivation let
\[
 F(y)\defeq\D{x}{u}\Bigl(C\bigl(y,0,C(u,0,x)\bigr)
   =C\bigl(y,0,C(u,0,P(y))\bigr)\Bigr).
\]
$F(0)$ normalizes to $\D{x}{u}(0=0)$ and is introduced at whatever parameter
is the newest.  In the branch for $F(\Suc z)$ the parameter $z$ is the
newest, and $F(\Suc z)$ normalizes to
$\D{x}{u}(C(u,0,x)=C(u,0,P(\Suc z)))$, that is, to
$\operatorname{Name}(z)$, because $P(\Suc z)\equiv z$.  Its introduction
selects the witness $z$; the instance at $0$ is $0=0$, so the mode is by
progress, and the instance asked for is $z=z$, by \rulename{Use} of
$\mathsf{Rf}$.  Numerical cases give $q_yF(y)$, and its elimination at the
term $\Suc t$ gives $\nf(F(\Suc t))$, which is literally
$\operatorname{Name}(t)$, since $C(\Suc t,0,b)\equiv b$ and
$P(\Suc t)\equiv t$.  The derivation is constructed anew in each scope.
\end{proof}

The proof uses no arbitrary witness, no equality substitution and no
predecessor theorem.  The shift by one is necessary: the value $P(0)=1$ is
erased by the outer selector.

\subsection{Completeness}

The source of the translation is the diagonal calculus $\Nat_{\mathrm{diag}}$ of
\cref{guard:thm:diagonal} (see also \cref{wit:sec:diagonal}), over the terms of \cref{base:sec:terms} extended by
the fixed descriptions $R$ with the equations \eqref{ctr:eq:r2zero} and \eqref{ctr:eq:r2succ}; it is stated
with transport in formula templates, which is a derived rule of it.  The
iterator is translated to $J$ with an unused counter coordinate, equations
and binders are translated compositionally, $\Q{y}B$ is translated to
$q_yB$, and all terms are normalized.  A diagonal matrix has the normalized
instance $0=0$ at $y=0$, whatever its witness and its bodies are.

\begin{theorem}[Completeness of the three families]
\label{wit:thm:recursor}
Let $\Gamma\proves A$ be derivable in $\Nat_{\mathrm{diag}}$ over the extended
vocabulary.  In every scope that contains the parameters, and from the
clauses $[\ ;\,v.B^{*}]$ for $B\in\Gamma$, $\mathsf{P}_{\nu}$ derives the
normalized translation $A^{*}$.  In particular closed theorems and the
contradiction are preserved.
\end{theorem}

\begin{proof}
By induction on the source derivation the following is shown: for every
target scope and context, every assignment of target terms to the source
parameters, and persistent builders of the translations of the source
assumptions, there is a persistent builder of the translation of the
conclusion.  At the root the builder of an assumption is its use.  Every
replacement of a local source assumption made below uses only the use of a
clause, reflexivity, and transport through the dummy major premise
$\D{x}{y}(0=0)$, whose introduction does not depend on the selected
witness; the builders therefore remain persistent.

\emph{Computation axioms} translate to reflexivity, their two sides having
the same normal form.  \emph{Hypotheses} invoke their builders.

\emph{Elimination of $\D{}{}$.}  Translate the major premise and apply
\rulename{Discharge} with a constant template and $0=0$.  The source
assumption $\Q{y}B(b,y)$ of the branch is represented by the clause
$[B(b,0);\,y.B(b,y)]$.  Turn the clause into the formula $q_yB(b,y)$ by an
introduction with an unused witness: in its premise use the clause at
$\Suc z$, the antecedent being the local assumption, or reflexivity if the
rule selected progress.  This builder is persistent and realizes the
assumption that the induction hypothesis for the branch requires.

\emph{Elimination of $\Q{}$.}  Translate the major premise to $q_yB(y)$ and
open it first.  Then apply the induction hypothesis to the derivation of
$B(0)$ in the extended context, use the new clause at the required term, and
close the \rulename{Discharge} with a constant template and $0=0$.  The
premise $B(0)$ is not translated outside the opening and then moved into the
branch.

\emph{Transport.}  Open the dummy major premise $\D{x}{y}(0=0)$, apply the
induction hypothesis to the two source premises in the extended context,
and finish the \rulename{Discharge} with the template of the source
inference.

\emph{Introduction at a witness $t$.}  Derive $\operatorname{Name}(t)$ by
\cref{wit:lem:name-recursor} and open it at a fresh $b$.  The branch has the
persistent equation $b=t$, and $b$ is the newest parameter, so
\rulename{Intro} selects $b$.  If the source introduction is by progress,
its matrix is diagonal, its normalized instance at $0$ is $0=0$ also at $b$,
and the target rule selects progress: transport the assumption at $b$ to the
assumption at $t$, apply the induction hypothesis to the source premise, and
transport the result back to $b$.  If the source introduction is by cases
and the target rule selects cases, do the same with the assumption at $0$.
If the target rule selects progress, the instance at $0$ at the witness $b$
is reflexive: derive it by \rulename{Use} of $\mathsf{Rf}$, transport it to
the instance at $0$ at the witness $t$, apply the induction hypothesis, and
transport back; the progress assumption is unused.  For each transport back,
the dummy major premise is opened before the induction hypothesis is applied
to the source premise, so that the recursive call is made in the extended
context; the builders of its assumptions are persistent there.  The
introduction yields $\D{x}{y}B$, which does not contain $b$, and the opening
of the name is closed with a constant template and $0=0$.

Every recursive call is for a proper subderivation of the source, and no
completed derivation of the target is moved into a new scope.
\end{proof}

The translation covers formulas that contain $R$ essentially and not only
the image of the iterator terms: the logical comparison and the diagonal
argument treat the two sides of an equation as arbitrary term templates, and
the additional computation axioms translate to reflexivity.

\subsection{The ordinary theory and the bridge}

Let $T_R$ be the theory $\mathsf T_R^2$ of \cref{ctr:def:ordinary}:
intuitionistic first-order logic with equality over the fixed descriptions
$R$, with the equations \eqref{ctr:eq:r2zero} and \eqref{ctr:eq:r2succ},
numerical coverage, the separation axiom $0=1\to\bot$ and contextual
induction for literal equations.  By body extensionality
(\cref{ctr:lem:bodyext}) $T_R$ proves $t=\nf(t)$ for every term, also with
reductions inside bodies, and every formula is provably equivalent to its
normalization (\cref{ctr:lem:normsound}).  No truth or evaluation predicate
is used.

\begin{theorem}[Bridge for the recursor calculus]
\label{wit:thm:recursor-bi}
$\mathsf{P}_{\nu}$ is sound and complete for $T_R$, and $T_R$ is
bi-interpretable with $\iPRA$: both domains are the whole numerical sort,
equality is equality, both numerical comparison maps have the graph $x=y$,
and every formula is provably equivalent to its double translation.  In
particular $\mathsf{P}_{\nu}$ derives $0=1$ if and only if $\iPRA$ proves
$\bot$.
\end{theorem}

\begin{proof}
Soundness is \cref{wit:prop:P-sound} with \cref{ctr:lem:normsound}.
Completeness is \cref{wit:thm:recursor} together with the equivalence of
$\Nat_{\mathrm{diag}}$ with $\Nat$ (\cref{guard:thm:diagonal}) and the
presentation of the ordinary theory by $\Nat$ (\cref{base:thm:bridge}),
which are parametric in the additional fixed descriptions.  The formula
round trips between $\mathsf{P}_{\nu}$ and $T_R$ follow by induction on
formulas from the reading of $\D{}{}$ and the logical macros
(\cref{wit:sec:liberal}).

The bi-interpretation of $T_R$ with $\iPRA$, with the stated domains,
equality, composites and double translations, is \cref{ctr:thm:transfer};
quantifier-free induction in $T_R$ is \cref{ctr:prop:qfi}
(\cref{wit:rem:reflection}).  The packages state the interpretations
directly, as recorded in \cref{ctr:rem:direct}: a description is
interpreted by the primitive recursion with the base $\Suc a$ and the step
$h$, and primitive recursion is compiled by $J$ without pairing
(\cref{ctr:thm:compiler}); $P$ applied to $\Suc x=\Suc y$ gives the
injectivity of the successor, and $C(\cdot,0,1)$ applied to $\Suc x=0$
gives $1=0$.  Falsity corresponds to $0=1$, which $T_R$ refutes by
separation, and derivations of a contradiction translate in both
directions.
\end{proof}

\subsection{Checking a template through its two instances}

In \rulename{Discharge} the checker receives a template $F(v)$, terms $s,t$
and derivations of $s=t$ and $\nf(F(s))$, and its result $\nf(F(t))$ must
equal the conclusion, which is checked.  A separate check that $F$ is
well-scoped and normal is redundant.

\begin{lemma}[Completion of an accepted template]
\label{wit:lem:template}
Let $F$ be a finite expression such that the substitution of the checked
terms $s$ and $t$ for its distinguished variable, followed by
normalization, succeeds, and such that the two results are well-scoped
formulas in the local scope.  Then there is a well-scoped template $F'(v)$
in normal form with
\[
 \nf(F'(s))=\nf(F(s)),\qquad\nf(F'(t))=\nf(F(t))
\]
literally.
\end{lemma}

\begin{proof}
\emph{Sorts.}  The traversal that substitutes and normalizes checks every
proper child at its required sort, term or formula, before any reduction
that erases it.  At the root it may accept an unclassified expression.  A
root that is neither an equation, a binder nor a recursor does not reduce
and yields no formula; a numerical root stays numerical under reduction.
Since an instance is a formula, the root of $F$ is an equation or a binder,
and $F$ is a sorted formula template that may contain additional numerical
parameters and reducible subterms.

\emph{Additional parameters.}  Read every parameter of $F$ outside the local
scope, and every bound-variable index that exceeds the binders above it by
more than the one position of the distinguished variable, as an additional
free parameter; the reading is consistent under a binder because the excess
is measured relative to the depth.  Let $\rho$ replace all additional
parameters by $0$ and fix $v$, the local parameters and the bound variables.
Neither $s$ nor $t$ contains an additional parameter.  For $r\in\{s,t\}$,
by \cref{ctr:lem:nf},
\[
 \nf\bigl((\rho F)[r/v]\bigr)=\nf\bigl(\rho(F[r/v])\bigr)
 =\nf\bigl(\rho\,\nf(F[r/v])\bigr)=\nf(F[r/v]),
\]
the last step because the checked instance is normal and contains no
additional parameter.  A reduction that becomes possible after an
additional parameter is closed does not affect the argument, which uses the
law of \cref{ctr:lem:nf} and not an identity of reduction sequences.

\emph{The template.}  Put $F'\defeq\nf(\rho F)$.  It is sorted, well-scoped
and normal, and a second application of \cref{ctr:lem:nf} gives the two
identities.  The distinguished variable may occur any number of times.
\end{proof}

Replacing $F$ by $F'$ leaves the two premises and the conclusion unchanged,
so every accepted use of the relaxed interface is an instance of
\rulename{Discharge}, and every canonical template is still accepted.  The
equation $s=t$ must be checked before the template is instantiated, so that
$s$ and $t$ are known to be terms, and the instance in the branch and the
conclusion keep their checks; the lemma does not justify omitting either.

% ---------------------------------------------------------------------
\section{Conditional copy}
\label{wit:sec:copy}

This section draws on \pkg{conditional\_copy}; the example of an unsound
probe is from \pkg{deterministic\_introduction}.

The calculus $\mathsf{P}_{\nu}$ takes no witness at its introduction and
still takes an equation, with its derivation, and a template at its
discharge.  Other calculi infer the substitution of a discharge from two
instances of the opened matrix and take an arbitrary witness at their
introduction (\cref{disch:def:selector,disch:thm:inferred}).  The two
restrictions are combined here.  The composition is
not immediate: with an inferred discharge, equality transport is recovered
from packets that are introduced at arbitrary witnesses, and the names of
\cref{wit:lem:name-recursor} give an equation $b=t$, whose use would need
the transport that is being recovered.  The combination rests on a formula
that copies an instance of a formula to a fresh parameter without any
equation.

\subsection{Use, Close and Export}

Terms, normal forms, the abbreviations $P$, $C$, $J$, formulas, clauses and
the initial clause $\mathsf{Rf}$ are those of \cref{wit:sec:recursor}.
Write $\top_{*}\defeq(0=0)$ and $\bot_{*}\defeq(0=1)$.  The calculus
$\mathsf{K}$ has three rule families; all displayed instances are
normalized.
\begin{enumerate}[label=\textup{(\roman*)}]
\item \rulename{Use}: as in $\mathsf{P}_{\nu}$.
\item \rulename{Close}: with $\ell=\nu(\Pi)$ and a fresh $z$,
\[
 \infrule{\Pi,z;\ \Gamma,[\ ;\,v.\nf(B(\ell,\eta))]\proves
          \nf(B(\ell,\Suc z))}
         {\Pi;\Gamma\proves\D{x}{y}B}
 \qquad
 \eta=\begin{cases}
  z,&\text{$\nf(B(\ell,0))$ is literally $\top_{*}$,}\\
  0,&\text{otherwise.}
 \end{cases}
\]
 The test for $\top_{*}$ is narrower than the test for a reflexive equation
 in $\mathsf{P}_{\nu}$ and suffices.
\item \rulename{Export}: with a fresh $b$, put
\[
 B_0=\nf(B(b,0)),\qquad B_1=\nf(B(b,1)),\qquad
 r_B=\begin{cases}
  t,&\text{$B_0\equiv\top_{*}$ and $B_1\equiv(b=t)$,}\\
  b,&\text{otherwise,}
 \end{cases}
\]
 where $\equiv$ is literal identity, including the orientation of the
 equation, and let $\mathcal{C}_B(b)=[B_0;\,y.\nf(B(b,y))]$.  Then
\[
 \infrule{\Pi;\Gamma\proves\D{x}{y}B\qquad
          \Pi,b;\ \Gamma,\mathcal{C}_B(b)\proves F}
         {\Pi;\Gamma\proves\nf(F[r_B/b])}
 \qquad b\notin\FV\bigl(\nf(F[r_B/b])\bigr).
\]
 $F$ may contain $b$, and so may $r_B$; the substitution is one simultaneous
 pass.  No equation, replacement term or template is supplied.
\end{enumerate}
Only \rulename{Use} takes a term.  The term $r_B$ is the selected witness
$\rho_B(b)$ of \cref{disch:def:selector}, and \rulename{Export} is the rule
\eqref{disch:eq:export} of the inferred discharge, here on the normal terms
of the counter-aware recursor (\cref{disch:thm:inferred} treats it with a
supplied witness at the introduction, on iterator terms).

\begin{lemma}[Soundness]
\label{wit:lem:k-sound}
Every rule of $\mathsf{K}$ is valid in $T_R$ under the ordinary reading.
\end{lemma}

\begin{proof}
\rulename{Use} and \rulename{Close} are as in \cref{wit:prop:P-sound}; if
the instance at $0$ is $\top_{*}$, the matrix is a literal equation, since
normalization does not remove a binder.  For \rulename{Export} the argument
is that of \cref{disch:prop:inferred-sound}: open the existential
quantifier; in the selected case reflexivity proves $B_0$, the clause at $1$
gives $b=r_B$, and the substitution of equals turns $F$ into $F[r_B/b]$; in
the default case the substitution is the identity.  The elimination of the
quantifier is legitimate because the normalized result does not contain
$b$.  Normalization is justified in $T_R$ by \cref{ctr:lem:normsound}.
\end{proof}

\begin{example}[The instance at $0$ must be tested]
\label{wit:ex:first-row}
An \rulename{Export} that inspected only the instance at $1$ would be
unsound.  Let $f$ be a term with $f(0)=f(1)=0$ and $f(2)=1$, and let
$B(x,y)$ be $x=f(y)$.  The reading of $\D{x}{y}B$,
$\exists x\,(x=f(0)\to\forall y\,x=f(y))$, is provable with the witness $1$,
whose antecedent is false.  The instance at $1$ at a fresh $b$ is $b=0$.  A
rule that replaced $b$ by $0$ on this evidence would produce the clause
$0=0\to\forall y\,0=f(y)$, which gives $0=1$ at $y=2$.  The test
$B_0\equiv\top_{*}$ excludes this matrix; \cref{disch:rem:conditions}
records a certificate for $1=0$ accepted without it.  For the same reason
the name
$\D{x}{y}\mathrm{Nm}_t$ of \cref{wit:sec:guarded} is not silently combined
with an inferred opening: opening a name directly at $t$ removes the fresh
alias $b=t$ on which the simulation of an arbitrary witness relies.
\end{example}

\subsection{Preparations without equality}
\label{wit:sec:copy-prep}

If $F$ does not contain the opened parameter $b$, the substitution of
\rulename{Export} does nothing to it; so the ordinary elimination of
$\D{x}{y}B$ is available whatever $r_B$ is.  As in
\cref{wit:sec:recursor}, $q_yA(y)$ is $\D{d}{y}A(y)$ with $d$ unused;
opening it and using its clause gives $A(t)$ from $q_yA$ and $A(0)$, and the
witness selected by \rulename{Close} for it is irrelevant.  Given persistent
constructions of $A(0)$ and of $A(\Suc z)$, \rulename{Close} introduces
$q_yA(y)$: if the instance at $0$ is $\top_{*}$, supply reflexivity in place
of the assumption at $0$ and ignore the progress assumption; otherwise the
rule is by cases.  The same holds for a premise that assumes $A(0)$.

Every formula is a chain of binders ending in an equation.  For a formula
$A$ with $d$ binders let $G_vA$ be as in \cref{wit:sec:recursor}; then
\[
 G_0A\equiv T_d,\qquad G_{\Suc z}A\equiv A,\qquad
 T_0=\top_{*},\quad T_{d+1}=\D{x}{y}T_d .
\]
Every tower $T_d$ is derivable in every scope, by reflexivity at depth $0$
and \rulename{Close} above it; none of its witnesses occurs in its body.
For formulas $A$ and $B$ with the same number of binders, align the binders
and replace corresponding sides $s,t$ of the final equations by $C(v,s,t)$.
This defines $\operatorname{Sel}(v,A,B)$, with
\[
 \operatorname{Sel}(0,A,B)\equiv A,\qquad
 \operatorname{Sel}(\Suc z,A,B)\equiv B .
\]
A substitution in terms does not change the number of binders, so $H(s)$
and $H(x)$ always have the same depth.  For formulas of different depth the
padding $L(A)=q_vG_vA$, interderivable with $A$, equalizes the depths.  All
of these are constructions carried out in the scope where they are needed.

\subsection{The conditional-copy theorem}

\begin{theorem}[Conditional copy]
\label{wit:thm:copy}
For every formula $H(x)$ with a distinguished variable and every term $t$
over $\Pi$, the formula
\[
 \Delta_t(H)\defeq\D{x}{y}\operatorname{Sel}\bigl(y,H(t),H(x)\bigr)
\]
is derivable in $\mathsf{K}$ in every scope and context, by a construction
that uses no equality transport and no equation other than reflexivity.
\end{theorem}

\begin{proof}
Construct the closure
\begin{equation}
 q_vG_v\Delta_{P(v)}(H).
 \label{wit:eq:copy-closure}
\end{equation}
Its instance at $v=0$ is a tower.  In the branch for $v=\Suc z$, with $z$
fresh, the conclusion normalizes to $\Delta_z(H)$, because
$P(\Suc z)\equiv z$ and the gate disappears.  The parameter $z$ is the
newest, so \rulename{Close} for $\Delta_z(H)$ selects $x=z$.  Its instance
at $0$ and its instance at a successor are both literally $H(z)$:
\[
 \operatorname{Sel}(0,H(z),H(z))\equiv H(z),\qquad
 \operatorname{Sel}(\Suc w,H(z),H(z))\equiv H(z).
\]
If $H(z)\not\equiv\top_{*}$, the rule is by cases and its assumption is its
conclusion.  If $H(z)\equiv\top_{*}$, the rule is by progress and
reflexivity derives the instance asked for.  The outer introduction is by
cases, since its instance at $0$ has at least one binder; its assumption is
unused.  Open \eqref{wit:eq:copy-closure} by the ordinary elimination,
supply the tower as the instance at $0$, and use the clause at $\Suc t$.
The gate disappears and $P(\Suc t)\equiv t$, so the result is literally
$\Delta_t(H)$; it does not contain the opened parameter.  Every step can be
carried out in every scope.
\end{proof}

The ordinary reading of $\Delta_t(H)$ is $\exists x\,(H(t)\to H(x))$.  The
theorem neither proves $H(t)$ nor identifies $x$ with $t$.  Opening
$\Delta_t(H)$ gives a fresh parameter $b$ and the clause
$[H(t);\,v.\operatorname{Sel}(v,H(t),H(b))]$, which yields $H(b)$ at $v=1$
once $H(t)$ is supplied.

\subsection{Removal of the witness}

Let $\mathsf{L}$, the \emph{liberal} calculus, have the syntax,
\rulename{Use} and \rulename{Export} of $\mathsf{K}$ and a \rulename{Close}
that takes an arbitrary term $t$ over $\Pi$ in place of $\ell$, with the
same test: progress exactly when $\nf(B(t,0))\equiv\top_{*}$.

\begin{theorem}[Witness removal]
\label{wit:thm:witness-removal}
For every scope $\Pi$, clause context $\Gamma$ and formula $A$, the calculi
$\mathsf{K}$ and $\mathsf{L}$ derive the same judgments
$\Pi;\Gamma\proves A$.
\end{theorem}

\begin{proof}
A \rulename{Close} of $\mathsf{K}$ is a \rulename{Close} of $\mathsf{L}$.
For the converse a derivation of $\mathsf{L}$ is translated by induction on
its tree.  The induction statement admits an injective map $\theta$ from the
source parameters into an arbitrary target scope and a map from the
positions of the source clauses to positions of target clauses, such that
the clause at each image position is literally the normalized
$\theta$-image of the source clause; further target parameters and clauses
are allowed.  This is an invariant of membership; it is not an assumption
that completed target derivations can be weakened.

\rulename{Use} follows the map of positions and translates the derivation
of the antecedent.  \rulename{Export} translates its major premise, sends
its eigenvariable to the new target parameter and extends the map of
positions by the opened clause; its two tests are preserved and reflected by
injective renaming and \cref{ctr:lem:nf}.

Consider a \rulename{Close} at a term $t$ with the matrix $B$ and the
premise $\Gamma,[\ ;\,v.\nf(B(t,\eta))]\proves\nf(B(t,\Suc z))$.  Write
$t$, $B$ also for their images, and put $H(x)\defeq q_yB(x,y)$.
\begin{enumerate}[label=\textup{(\alph*)}]
\item Derive $\Delta_t(H)$ by \cref{wit:thm:copy} and open it at a fresh
 $b$; its clause is $[H(t);\,v.\operatorname{Sel}(v,H(t),H(b))]$.
\item Apply \rulename{Close} to $\D{x}{y}B$.  The newest parameter is $b$,
 so the rule selects $b$; let $z'$ be its eigenvariable.
\item In the premise derive $H(t)$ by a \rulename{Close} with an unused
 witness.  Its test is the source test at $t$, and its premise, with a
 further eigenvariable, is the image of the source premise: the induction
 hypothesis applies with the source eigenvariable sent to the new one and
 the source assumption sent to the local clause of this \rulename{Close}.
\item \rulename{Use} of the clause of (a) at $1$ gives $H(b)$.
\item Open $H(b)$; its clause is $[B(b,0);\,y.B(b,y)]$.  If the
 \rulename{Close} of (b) is by cases, its assumption is $B(b,0)$.  If it is
 by progress, $B(b,0)$ is literally $\top_{*}$ and reflexivity supplies it.
 \rulename{Use} at $\Suc z'$ gives $B(b,\Suc z')$, and the unused witness
 of $H(b)$ is discharged.
\end{enumerate}
The \rulename{Close} of (b) yields $\D{x}{y}B$, which does not contain $b$,
and the ordinary elimination discharges $b$.  Every recursive call is for
the premise of the source \rulename{Close}, a proper subderivation, although
it may be made several times; the packet and the towers are explicit
constructions.  At each call the tests agree, because they are made on the
same instantiated formula.
\end{proof}

\begin{remark}[The modes need not agree]
\label{wit:rem:mode}
The test at $t$ and the test at $b$ may differ.  The proof transfers the
whole conditional closure $H$ and not a progressive premise; at $b$ a
progressive mode has a reflexive instance at $0$, so the closure still
yields the instance at a successor.  No equality reasoning repairs a failed
test.  The same proof applies with the rule \rulename{Name} of
\cref{wit:sec:name-variant} in place of \rulename{Export}.
\end{remark}

\subsection{The liberal calculus presents the ordinary theory}
\label{wit:sec:liberal}

\begin{lemma}[Structural facts of the liberal calculus]
\label{wit:lem:liberal-structural}
In $\mathsf{L}$, weakening by clauses, injective renaming of parameters and
the cut of an assumption $[\ ;\,v.A]$ are admissible.
\end{lemma}

\begin{proof}
By induction on derivations.  A supplied witness remains available after a
renaming, an injective renaming preserves and reflects the literal tests of
\rulename{Close} and \rulename{Export}, and normalization commutes with
renaming.  For the cut, a \rulename{Use} of the clause is replaced by the
given derivation, weakened and renamed into the scope of the leaf.  No
equality transport and no substitution of terms into derivations is used.
The lemma is about $\mathsf{L}$ and is not claimed for $\mathsf{K}$ at this
stage.
\end{proof}

\begin{lemma}[Transport from an inferred discharge]
\label{wit:lem:liberal-transport}
In $\mathsf{L}$, from $s=t$ and $h(s)=u$ one derives $h(t)=u$.
Consequently symmetry, transitivity, congruence and transport in arbitrary
formulas are derived.
\end{lemma}

\begin{proof}
Let
\[
 E(x,y)\defeq\Bigl(C\bigl(y,0,C(P(y),x,h(x))\bigr)
   =C\bigl(y,0,C(P(y),t,u)\bigr)\Bigr).
\]
Its normalized instances are $E(x,0)\equiv\top_{*}$, $E(x,1)\equiv(x=t)$
and $E(x,\Suc\Suc v)\equiv(h(x)=u)$.  Introduce $\D{x}{y}E$ at the witness
$s$; the mode is by progress, and the instance asked for,
$C(z,s,h(s))=C(z,t,u)$, is derived by numerical cases on $z$ from the two
given equations, the progress assumption being unused.  Open the formula at
$b$: the first two instances make \rulename{Export} select $t$.
\rulename{Use} at $2$, with the reflexive antecedent, gives $h(b)=u$, and
\rulename{Export} concludes $h(t)=u$.  This is the argument of
\cref{disch:thm:export-transport}, with the decoder packet
\eqref{ctr:eq:packet} of \cref{ctr:lem:openrules} in place of the packet on
iterator terms.

With reflexivity this rule gives symmetry, transitivity and congruence, and
transport in an equation template is two such replacements.  Transport in a
formula $\D{x}{y}B(v,x,y)$ follows by induction on the number of binders:
open the given formula at its witness, introduce the target at that
witness, transport the assumption at $0$ backwards through the matrix, use
the opened clause at the successor and transport the result forwards; if
the new instance at $0$ is $\top_{*}$, supply reflexivity instead.  The
result does not contain the opened witness.  The premises are moved into the
local scopes by \cref{wit:lem:liberal-structural}.
\end{proof}

\begin{lemma}[Cases and equation induction in the liberal calculus]
\label{wit:lem:liberal-induction}
In $\mathsf{L}$ the introduction by cases is available for every matrix and
witness, and for a literal equation $E(y)\equiv(s(y)=t(y))$, from $E(0)$
and a derivation of $E(\Suc z)$ from $E(z)$ one derives $E(a)$ for every
term $a$.
\end{lemma}

\begin{proof}
If an introduction by cases selects progress because its instance at $0$ is
$\top_{*}$, replace the assumption at $0$ by reflexivity and ignore the
progress assumption.  For the induction let
\[
 K(y)\defeq\bigl(C(y,0,s(y))=C(y,0,t(y))\bigr),
\]
whose instance at $0$ is literally $\top_{*}$.  Numerical cases, without
induction, give the two decoder equations
\[
 C\bigl(y,s(0),C(y,0,s(y))\bigr)=s(y),\qquad
 C\bigl(y,s(0),C(y,0,t(y))\bigr)=t(y);
\]
the first has reflexive branches, and the second uses $E(0)$ at $0$ and is
reflexive at successors.  From $K(z)$, congruence with
$v\mapsto C(z,s(0),v)$ and the decoders give $E(z)$.  The cut of this into
the step gives $E(\Suc z)$, which is literally $K(\Suc z)$.  The
introduction by progress gives $q_yK(y)$, \rulename{Use} with the reflexive
antecedent gives $K(a)$, and the decoders give $E(a)$.  No compound term is
substituted into a derivation.  A conditional closure $q_yE(y)$ is
introduced by cases, with $E(0)$ as its assumption and the induction just
derived for its successor instance.
\end{proof}

The ordinary logical vocabulary is given by the macros
\eqref{ctr:eq:macros} of \cref{ctr:def:macros}, with
$\operatorname{Sel}(v,A,B)$ in the role of $M_v(A,B)$, the depths
equalized by $L$ before each binary selector and all bound variables fresh;
in particular $(A\to B)^{*}=q_y\operatorname{Sel}(y,A^{*},B^{*})$,
$(A\wedge B)^{*}=q_yG_y\operatorname{Sel}(P(y),A^{*},B^{*})$,
$(\forall x\,A)^{*}=q_yG_yA^{*}(P(y))$ and $\bot^{*}=\bot_{*}$.  The proofs
of \cref{ctr:lem:gates,ctr:lem:macros} use \rulename{Use}, the ordinary
elimination, numerical cases, the introduction by cases at an arbitrary
witness, the introduction by progress, transport in formulas and binder
congruence (proved as \cref{ctr:lem:libderived}(5)).  All of these are
available in $\mathsf{L}$ by
\cref{wit:lem:liberal-structural,wit:lem:liberal-transport,wit:lem:liberal-induction}
and, for the ordinary elimination, \cref{wit:sec:copy-prep}, so the proofs
apply unchanged: the macros admit the introduction and
elimination rules of intuitionistic logic under arbitrary assumptions, with
explosion for $\bot_{*}$, and by induction on a formula the ordinary reading
of each macro is equivalent to the intended connective
(\cref{ctr:thm:adequacy}(3)).
\pkg{parameter\_intro} uses the same macros with two differences: the
conjunction selects by a parity term,
$(A\wedge B)^{*}=q_yG_y\operatorname{Sel}(\operatorname{par}(y),A^{*},B^{*})$,
and the universal quantifier is
$(\forall x\,A)^{*}=A^{*}(0)\wedge q_xA^{*}(x)$.

\begin{proposition}[The liberal calculus]
\label{wit:prop:liberal}
$\mathsf{L}$ is sound and complete for $T_R$.  Hence so is $\mathsf{K}$.
\end{proposition}

\begin{proof}
Soundness is \cref{wit:lem:k-sound} with a supplied witness.  For
completeness translate the diagonal calculus $\Nat_{\mathrm{diag}}$ over the
vocabulary with the fixed descriptions $R$, as in \cref{wit:thm:recursor}:
computation axioms become reflexivity after normalization; hypotheses are
used; equality is \cref{wit:lem:liberal-transport}; the eliminations of
$\Q{}$ and $\D{}{}$ are the derived ones; a diagonal matrix has the instance
$0=0$ at $0$, so its introduction by progress is the progress mode of
\rulename{Close}; an introduction by cases is
\cref{wit:lem:liberal-induction}.  The structural facts of
\cref{wit:lem:liberal-structural} move derivations between scopes.  This
gives the full consequence relation and not only closed equations.  The
statement for $\mathsf{K}$ follows by \cref{wit:thm:witness-removal}, which
regenerates the branches in the target scopes.
\end{proof}

\subsection{Quantifier-free induction and the bridge}

\begin{remark}[Reflection of equality in $T_R$]
\label{wit:rem:reflection}
$T_R$ proves the successor injective and equality decidable, and it proves
every instance of contextual quantifier-free induction
(\cref{ctr:prop:qfi}); the reflection term of \cref{base:thm:reflect}
transfers along \cref{ctr:thm:transfer}.  \pkg{parameter\_intro} and
\pkg{conditional\_copy} prove this directly in $T_R$, with the terms of
\cref{ctr:sec:counter}, by the argument recorded in
\cref{ctr:rem:reflection}: addition and truncated subtraction $d$ compiled
by $J$, four inductions on literal equations, one induction on the residual
equation $C(\operatorname{sg}(d(x,y)),x,y+d(x,y))=x$ whose case splits are
all coverage, the reflection term $e(a,b)=\operatorname{sg}(K(d(\Suc a,b)))$
with $K(u)=C(u,1,p(u))$, and characteristic terms $b_A$ with $0$ for truth.
\end{remark}

\begin{theorem}[Introduction and discharge without supplied terms]
\label{wit:thm:k-bi}
$\mathsf{K}$ is a sound and complete presentation of $T_R$, and $T_R$ is
bi-interpretable with $\iPRA$, with the domains $x=x$, equality as
equality, the graph $x=y$ in both numerical composites and a provable
equivalence of every formula with its double translation.  Hence
$\mathsf{K}$ and $\iPRA$ are equiconsistent.
\end{theorem}

\begin{proof}
The first statement is \cref{wit:prop:liberal}.  The bi-interpretation of
$T_R$ with $\iPRA$, with the stated domains, equality, composites and
double translations, is \cref{ctr:thm:transfer}, as in the proof of
\cref{wit:thm:recursor-bi}; quantifier-free induction in $T_R$ is
\cref{ctr:prop:qfi}; the formula round trips between $\mathsf{K}$ and $T_R$
follow from the macros of \cref{wit:sec:liberal} by induction on formulas.
Falsity is translated to $0=1$, whose reading implies $\bot$ by the
separation axiom of $T_R$.
\end{proof}

Here $\iPRA$ has quantifier-free induction, and $T_R$ has induction for
literal equations only; the words ``first-order'' refer to the logical
language.

\subsection{Discharge by a supplied name}
\label{wit:sec:name-variant}

Let $\mathsf{K}_{\mathrm N}$ have the syntax, the clause $\mathsf{Rf}$,
\rulename{Use} and \rulename{Close} of $\mathsf{K}$, and in place of
\rulename{Export} the rule
\[
 \infrule{\Pi;\Gamma\proves\D{x}{y}B\qquad
          \Pi,b;\ \Gamma,\mathcal{C}_B(b)\proves b=t\qquad
          \Pi,b;\ \Gamma,\mathcal{C}_B(b)\proves F}
         {\Pi;\Gamma\proves\nf(F[t/b])}\ \rulename{Name}
 \qquad b\notin\FV\bigl(\nf(F[t/b])\bigr).
\]
The term $t$ is over $\Pi,b$ and may contain $b$; it is inserted once.  The
left side of the supplied equation is fixed to the opened parameter, whereas
\rulename{Discharge} of $\mathsf{P}_{\nu}$ takes both sides of an equation
and a template.

\begin{proposition}[The name variant]
\label{wit:prop:name-variant}
$\mathsf{K}_{\mathrm N}$, $\mathsf{K}$ and $T_R$ have the same open
consequences under their readings.
\end{proposition}

\begin{proof}
\rulename{Name} is sound in $T_R$: equality gives the substituted formula
and the condition on the scope permits the elimination of the quantifier.
Every instance of \rulename{Export} is an instance of \rulename{Name} in the
same scope: in the default case with the reflexivity $b=b$, in the selected
case with the equation $b=t$ obtained from the opened clause by reflexivity
and \rulename{Use} at $1$.  The replacement adds no parameter, so no
\rulename{Close} changes its witness or its mode, and induction on the
derivation gives $\mathsf{K}\subseteq\mathsf{K}_{\mathrm N}$.  The
completeness of $\mathsf{K}$ and the soundness of $\mathsf{K}_{\mathrm N}$
give the rest.
\end{proof}

A backward checker for $\mathsf{K}_{\mathrm N}$ receives a template $T(v)$
and a term $t$, checks the conclusions $b=t$ and $F=\nf(T[b/v])$ in the
branch, and requires the goal to be $A=\nf(T[t/v])$.  It does not forbid
occurrences of $b$ in $T$.

\begin{lemma}[Template representation]
\label{wit:lem:name-template}
Let $T(v)$ be a well-scoped template and $t$ a term of the local scope, and
put $F=\nf(T[b/v])$ and $A=\nf(T[t/v])$.  If $A$ does not contain $b$, then
$\nf(F[t/b])=A$ literally.
\end{lemma}

\begin{proof}
A renaming of parameters, injective or not, preserves normal forms: the
reductions test only for the constructors $0$ and $\Suc$ at a counter, and
an identification of parameters creates neither; no reduction turns an
equation with equal sides into $0=0$.  Choose a parameter $c$ not in $T$ or
$t$ and let $U=\nf(T[v:=t,\,b:=c])$, with simultaneous substitution.
Identifying $c$ with $b$ gives
\[
 U[c\mapsto b]=\nf(U[c\mapsto b])=\nf(T[t/v])=A .
\]
Since $A$ does not contain $b$, $U$ contains neither $b$ nor $c$, and
$U=A$.  Substituting $t$ for $c$ and \eqref{ctr:eq:SC} give
$\nf(T[v:=t,\,b:=t])=A$.  Finally
\[
 \nf(F[t/b])=\nf\bigl((T[b/v])[t/b]\bigr)=\nf(T[v:=t,\,b:=t])=A .
\]
No substitution is made inside an inserted copy of $t$.
\end{proof}

The checker validates a template through its two normalized instances, as
in \cref{wit:lem:template}, and the same completion applies: the excess
parameters of an accepted raw template are replaced by $0$, which preserves
both checked instances literally because neither they nor the checked
arguments $b$, $t$ contain an excess parameter; the normalized completed
template is well-scoped, and \cref{wit:lem:name-template} applies to it.
The equation premise validates $t$, and each instance is validated in its
own scope; these checks are needed.  Conversely, abstracting all free
occurrences of $b$ in $F$ into $v$ represents every instance of
\rulename{Name}.  The backward representation is therefore exact.

% ---------------------------------------------------------------------
\section{Specialization at the newest parameter}
\label{wit:sec:spec}

This section draws on \pkg{export\_specialization} for
\cref{wit:sec:spec-x} and on \pkg{vacuous\_specialization} for
\cref{wit:sec:vac,wit:sec:vac-copy}.

In every calculus considered so far the use of a clause takes an arbitrary
term.  Here it is restricted to the newest parameter, with $0$ in an empty
scope, so that it takes no term.  The term is delivered instead by an
inferred discharge, which substitutes a term read off the opened matrix.  In
\cref{wit:sec:spec-x} the introduction keeps an arbitrary witness, and the
names introduced at that witness carry the term.  In
\cref{wit:sec:vac,wit:sec:vac-copy} the introduction is that of
$\mathsf K$, at the newest parameter, so that no name can be introduced at
an arbitrary term; a closure whose matrix does not depend on its second
variable carries the term, and the selector of \rulename{Export} receives
one additional case.  Neither result needs a primitive alias rule
(\cref{disch:sec:alias-newest}) or the supplied witness of the
deterministic specialization (\cref{disch:def:spec}).

\subsection{Newest specialization with the inferred discharge}
\label{wit:sec:spec-x}

The calculus $\mathcal X$ of \cref{disch:def:inferred} has the rules
\rulename{Refl}, \rulename{Use}, \rulename{Intro} and \rulename{Export} of
\cref{disch:def:rules} and \eqref{disch:eq:export}, on full normal forms of
iterator terms with normalizing substitution, in the convention with
\rulename{Refl}.  Its scopes are written explicitly as in
\cref{wit:def:scope}: \rulename{Intro} appends its eigenvariable $z$ and
\rulename{Export} its parameter $b$.

\begin{definition}[Newest specialization]
\label{wit:def:xnu}
$\mathcal X_\nu$ is $\mathcal X$ with \rulename{Use} replaced by
\[
 \infrule{\Pi;\Gamma\proves A}{\Pi;\Gamma\proves\nf\bigl(B(\nu(\Pi))\bigr)}
 \ \rulename{Use}_\nu,
 \qquad [A;\,y.\,B(y)]\text{ an occurrence in }\Gamma .
\]
\rulename{Refl} is kept at every term over $\Pi$, and \rulename{Intro} at
every witness over $\Pi$, with the mode of \eqref{disch:eq:intro}.
\end{definition}

A certificate of $\rulename{Use}_\nu$ consists of the occurrence and the
derivation of its antecedent.  For a term $t$ over $\Pi$ let
\[
 \mathrm{Nm}_t\defeq\D{x}{y}\bigl(\sel(y,0,x)=\sel(y,0,t)\bigr),
\]
the name of \cref{disch:sec:spec}.

\begin{lemma}[A name without \rulename{Use}]
\label{wit:lem:x-name}
In every scope $\Pi$ over which $t$ is a term, and in every context,
$\mathcal X_\nu$ derives $\mathrm{Nm}_t$ by one \rulename{Intro} and one
\rulename{Refl}.  \rulename{Export} of $\mathrm{Nm}_t$ at a fresh $b$
selects $t$.
\end{lemma}

\begin{proof}
Introduce $\mathrm{Nm}_t$ at the witness $t$.  The instance asked for is
$\sel(\Suc z,0,t)=\sel(\Suc z,0,t)$, whose normal form is $t=t$; it is
proved by \rulename{Refl}, and the assumption of the premise is unused.
This holds in both modes, in particular for $t\equiv 0$, when the rule is by
progress.  The normalized rows at a fresh $b$ are $\sel(0,0,b)=\sel(0,0,t)$,
that is $0=0$, and $\sel(1,0,b)=\sel(1,0,t)$, that is $b=t$; so the selector
of \cref{disch:def:selector} returns $t$.
\end{proof}

\begin{theorem}[Newest specialization]
\label{wit:thm:x-newest}
For every scope $\Pi$, clause context $\Gamma$ and formula $A$,
\[
 \Pi;\Gamma\proves_{\mathcal X_\nu}A
 \quad\Longleftrightarrow\quad
 \Pi;\Gamma\proves_{\mathcal X}A ,
\]
by effective transformations that leave the end sequent unchanged.  In the
convention with \rulename{Refl}, a derivation of $\mathcal X$ with $n$
inference nodes, $u$ of them uses of a clause, is translated into a
derivation of $\mathcal X_\nu$ with exactly $n+3u$ nodes.
\end{theorem}

\begin{proof}
Every instance of $\rulename{Use}_\nu$ is an instance of \rulename{Use}.
For the converse the following is shown by induction on a derivation $\pi$
of $\Pi;\Gamma\proves_{\mathcal X}A$: for every injective map $\theta$ of
$\Pi$ into a scope $\Pi'$, and every context $\Delta$ over $\Pi'$ with an
injective map of the occurrences of $\Gamma$ to occurrences of $\Delta$ that
sends each clause to its normalized $\theta$-image, $\mathcal X_\nu$ derives
$\Pi';\Delta\proves\theta A$.  Further parameters and clauses of $\Delta$
are arbitrary, and $\theta$ need not preserve which parameter is the newest.
This is a construction carried out again in each target situation; it is not
a claim that a derivation of $\mathcal X_\nu$ can be weakened, which is false
node by node, because an extension of the scope changes the newest
parameter.

\rulename{Refl} is translated literally.  At \rulename{Intro} extend
$\theta$ by the eigenvariable sent to a fresh target parameter, extend the
map of occurrences by the assumption, and translate the premise; an
injective renaming preserves the root of the matrix and the literal $0$, so
the mode is unchanged.  At \rulename{Export} translate the major premise,
open at a fresh target $b'$, extend $\theta$ by $b\mapsto b'$ and the map of
occurrences by the opened clause, and translate the branch; by
\cref{disch:lem:selector}(a) the selected term is the image of the source
selection, and by \eqref{disch:eq:coherence} the exported formula is the
image of the source conclusion, which does not contain $b'$.

At a \rulename{Use} of the occurrence $[A_0;\,y.\,B(y)]$ at a term $t$, put
$t'\equiv\theta t$.  Derive $\mathrm{Nm}_{t'}$ by \cref{wit:lem:x-name} and
apply \rulename{Export} at a fresh $b$, which is the newest parameter of the
branch.  In the branch apply the induction hypothesis to the source
derivation of $A_0$, with the same $\theta$ and the same map of occurrences,
the opened clause and $b$ being extras; this gives $\theta A_0$.
$\rulename{Use}_\nu$ of the image of the occurrence gives
$\nf((\theta B)(b))$, and \rulename{Export} concludes
\[
 \nf\bigl((\theta B)(b)\subst{t'}{b}\bigr)\equiv\theta\,\nf\bigl(B(t)\bigr)
\]
by \eqref{disch:eq:coherence}.  Neither the context nor $t'$ contains $b$,
so the side condition holds.  For a clause without antecedent the inductive
call is omitted.  The source derivation of $A_0$ is a proper subderivation
of $\pi$, and it is translated in the situation where it is used.  The
identity map gives the equivalence.

For the count, every node other than a \rulename{Use} is copied once with
its translated premises, and a \rulename{Use} becomes four nodes: the
\rulename{Export}, the \rulename{Intro} of the name with its \rulename{Refl}
leaf, and the $\rulename{Use}_\nu$, above the translated antecedent, which
is translated once.  Induction on the derivation gives $n+3u$.
\end{proof}

\begin{remark}[The Lean module]
\label{wit:rem:x-lean}
\Lean{NativeNewestUse} formalizes $\mathcal X_\nu$ (\Lean{XNPrf}) over the
calculus \Lean{XCPrf} of $\mathcal X$: \rulename{Use} at \Lean{Tm.newest},
the variable $0$ of the de Bruijn scope or the literal $0$ in an empty scope,
and \rulename{Refl} at every term.  \Lean{xnprf\_iff\_xcprf} is
\cref{wit:thm:x-newest} with the reflexivity clause in the context, and
\Lean{xnprf\_iff\_nprf\_of\_qfree} the bridge on $Q$-free sequents.  The
translation is the one of the proof above: it carries a renaming and an
extension of the context (\Lean{XCPrf.toXN}); the name formula is introduced
at the normal form of the witness (\Lean{XNPrf.nm}, \Lean{NFm.nm\_step}),
its rows at the fresh variable select it (\Lean{NFm.nm\_otest}), and the
exported formula is the image of the source conclusion up to full normal
forms.  The count $n+3u$ is not formalized.
\end{remark}

The count is in inference nodes.  A name copies its term, and normalization
can enlarge formulas, so no bound on the size of certificates in bytes or
on the time of checking follows.

\begin{corollary}[The bridge for $\mathcal X_\nu$]
\label{wit:cor:x-bridge}
$\mathcal X_\nu$ derives the same sequents as the reference calculus
$\mathcal T$ of \cref{disch:def:reference}; weakening, injective renaming,
substitution of terms for parameters and cut are admissible in it on
derivable judgments; and its ordinary reading is bi-interpretable with
$\iPRA$, with whole domains, equality as equality and the graph $x=y$ in
both numerical composites, as for $\mathcal X^{\mathcal R}$ in
\cref{disch:thm:inferred}.
\end{corollary}

\begin{proof}
\Cref{wit:thm:x-newest}, \cref{disch:thm:inferred} and, for the two
conventions, \cref{disch:lem:conventions}.  The admissible
rules hold for the derivable judgments of $\mathcal X$ and transfer through
the equivalence; they are not preserved node by node.
\end{proof}

\begin{remark}[What the proof uses]
\label{wit:rem:x-uses}
The name is introduced at the arbitrary witness $t$, and its instance is the
reflexivity $t=t$ at a compound term.  The proof therefore restricts neither
the witness of \rulename{Intro} nor \rulename{Refl}.  In the convention with
the reflexivity clause $\mathcal R=[\ ;\,y.\,y=y]$ the restriction of all
uses of $\mathcal R$ to $\nu(\Pi)$ would leave only $\nu(\Pi)=\nu(\Pi)$, and
a derivation of $t=t$ from a name would presuppose the naming lemma; the
theorem is therefore stated in the convention with \rulename{Refl}, and its
form $\mathcal X^{\mathcal R}_\nu$ with the reflexivity clause exempts
$\mathcal R$ from the restriction.  The opened clause of the name is never
used: \rulename{Export} reads $t$ from its row at $1$.  This is the
difference from \cref{disch:lem:spec-name}, where $b=t$ is derived from the
naming clause by $\rulename{Use}_{*}$ at $1$, and from
\cref{disch:cor:alias-newest}, where a primitive \rulename{Alias} supplies
the equation.  \Cref{wit:sec:vac} restricts the witness of the introduction
as well.
\end{remark}

The same argument applies without a normalizer.  Let $\mathsf{C}$ be the
five-rule clause calculus of \cref{claus:def:c3}, with \rulename{Refl},
\rulename{Eval}, \rulename{Use}, \rulename{Elim} and \rulename{Intro} on raw
terms with ordinary substitution, and with scopes written explicitly.  The
\emph{raw selector} returns $t$ if the matrix $B(x,y)$ is literally
$\sel(y,0,x)=\sel(y,0,t)$ with $t$ free of $x$, $y$ and the variables bound
by the two selectors, and returns $b$ otherwise.  $\rulename{Export}_{\mathrm r}$
is \eqref{disch:eq:export} with this selector and ordinary substitution,
the exported formula being required to be over the old scope.

\begin{proposition}[The raw companion]
\label{wit:prop:x-raw}
The calculus with the five rule families \rulename{Refl}, \rulename{Eval},
$\rulename{Use}_\nu$, \rulename{Intro} and $\rulename{Export}_{\mathrm r}$
derives the same judgments $\Pi;\Gamma\proves A$ as $\mathsf{C}$.
\end{proposition}

\begin{proof}
Every instance of $\rulename{Export}_{\mathrm r}$ is derivable in
$\mathsf{C}$.  If the selector returns $b$, the exported formula is the
branch formula, which does not contain $b$, and the rule is \rulename{Elim}.
Otherwise open the major premise by \rulename{Elim} at $b$; the opened clause
is $[\sel(0,0,b)=\sel(0,0,t);\,y.\,\sel(y,0,b)=\sel(y,0,t)]$.  Both sides of
its antecedent are equal to $0$ by the computation equations of
\cref{claus:lem:eq}, so the antecedent is derived by symmetry and
transitivity; \rulename{Use} at $1$ gives $\sel(1,0,b)=\sel(1,0,t)$, and the
computation equations give $b=t$.  Transport in the formula $F\subst{v}{b}$
(\cref{claus:lem:ftransport}) turns the branch conclusion $F$ into
$F\subst{t}{b}$, and \rulename{Elim} discharges $b$, since $F\subst{t}{b}$
does not contain it.  $\rulename{Use}_\nu$ is an instance of \rulename{Use}.

Conversely translate a derivation of $\mathsf{C}$ with the induction
statement of the proof of \cref{wit:thm:x-newest}.  \rulename{Refl},
\rulename{Eval} and \rulename{Intro} are copied; an injective renaming
preserves the literal matches of \rulename{Eval}, the root of the matrix and
the literal $0$.  \rulename{Elim} becomes $\rulename{Export}_{\mathrm r}$:
its conclusion does not contain $b$, so the substitution of the selected
term changes nothing.  A \rulename{Use} at $t$ is simulated as before:
$\mathrm{Nm}_{t'}$ is introduced at $t'$, where the raw instance
$\sel(\Suc z,0,t')=\sel(\Suc z,0,t')$ has literally equal sides and is
proved by \rulename{Refl} in either mode; its matrix is the pattern of the
raw selector, which therefore returns $t'$; the antecedent is translated in
the branch, $\rulename{Use}_\nu$ gives $(\theta B)(b)$, and
$\rulename{Export}_{\mathrm r}$ gives $(\theta B)(b)\subst{t'}{b}\equiv
\theta(B(t))$, since $B$ does not contain $b$.  The pattern and its
condition on free variables are preserved and reflected by injective
renaming.
\end{proof}

No checker of the raw companion is measured.

\subsection{The vacuous selector}
\label{wit:sec:vac}

Terms, normal forms, the abbreviations $P$, $C$, $J$, formulas, clauses,
$\top_{*}$, $\bot_{*}$ and the initial clause $\mathsf{Rf}$ are those of
\cref{wit:sec:recursor,wit:sec:copy}.  For a matrix $B$ and a fresh $b$ write
\[
 \widehat B(y)\defeq\nf\bigl(B(b,y)\bigr),\qquad
 B_0\defeq\nf\bigl(\widehat B(0)\bigr),\qquad
 B_1\defeq\nf\bigl(\widehat B(1)\bigr).
\]
Since normalization introduces no free variable (\cref{wit:sec:recursor}),
$\widehat B\equiv B_0$ holds exactly when $y$ does not occur in
$\widehat B$; in that case $\widehat B\equiv B_0\equiv B_1$.

\begin{definition}[The calculus $\mathsf V$]
\label{wit:def:v}
$\mathsf V$ has three rule families.
\begin{enumerate}[label=\textup{(\roman*)}]
\item $\rulename{Use}_\nu$: for an occurrence $[A;\,y.\,B]$ in $\Gamma$, from
 $\Pi;\Gamma\proves A$ infer $\Pi;\Gamma\proves\nf(B(\nu(\Pi)))$; for a
 clause without antecedent the premise is omitted.
\item \rulename{Close}: the rule \rulename{Close} of $\mathsf K$
 (\cref{wit:sec:copy}), at $\ell=\nu(\Pi)$, by progress exactly when
 $\nf(B(\ell,0))\equiv\top_{*}$.
\item \rulename{Export}: the rule \rulename{Export} of $\mathsf K$ with the
 selector
\[
 r^{\mathsf V}_B\defeq
 \begin{cases}
  r,&\text{$B_1\equiv(b=r)$, and $B_0\equiv\top_{*}$ or
        $\widehat B\equiv B_0$,}\\
  b,&\text{otherwise,}
 \end{cases}
\]
 in place of $r_B$:
\[
 \infrule{\Pi;\Gamma\proves\D{x}{y}B\qquad
          \Pi,b;\ \Gamma,[B_0;\,y.\,\widehat B(y)]\proves F}
         {\Pi;\Gamma\proves\nf\bigl(F\subst{r^{\mathsf V}_B}{b}\bigr)}
 \qquad b\notin\FV\bigl(\nf(F\subst{r^{\mathsf V}_B}{b})\bigr).
\]
\end{enumerate}
There are no other rules; every derivation starts from the context
$\mathsf{Rf}$.
\end{definition}

No rule of $\mathsf V$ takes a term: \rulename{Use} takes an occurrence and
the derivation of its antecedent, \rulename{Close} a matrix and the
derivation of its premise, \rulename{Export} two derivations.  Terms still
occur in the matrices supplied to \rulename{Close}; what is removed is the
separate term argument, not the information it carries.  The first case of
$r^{\mathsf V}_B$ is $r_B$ of $\mathsf K$; the second is new, and $r$ may
contain $b$.  The two cases exclude each other, because in the second $B_0$
is the equation $b=r$, whose left side is a parameter.

\begin{lemma}[Soundness of the vacuous case]
\label{wit:lem:v-sound}
Every rule of $\mathsf V$ is valid in $T_R$ under the ordinary reading, also
when the selected term contains $b$.
\end{lemma}

\begin{proof}
$\rulename{Use}_\nu$, \rulename{Close} and the first case of
\rulename{Export} are covered by \cref{wit:lem:k-sound}.  In the second case
the opened clause reads $(b=r)\to\forall y\,(b=r)$, with $y$ not free, which
is a theorem of intuitionistic logic.  Suppose the reading of the branch is
derivable in $T_R$ from the readings of $\Gamma$ and of this clause; cutting
the clause gives a derivation of $F$ from
$\Gamma$ in which $b$ is a free variable not free in $\Gamma$.  Substituting
the term $r$ for $b$ throughout this derivation gives $F\subst{r}{b}$ from
$\Gamma$, also when $r$ contains $b$, and $T_R$ proves it equivalent to
$\nf(F\subst{r}{b})$ by \cref{ctr:lem:normsound}.  This formula does not
contain $b$.  The major premise is not needed in this case.
\end{proof}

\begin{lemma}[The vacuous packet]
\label{wit:lem:v-packet}
For every scope $\Pi$, context $\Gamma$ and term $t$ over $\Pi$, the formula
\[
 \mathsf V_t\defeq\D{x}{y}(x=t)
\]
is derived in $\mathsf V$ by one \rulename{Close} and one $\rulename{Use}_\nu$.
\end{lemma}

\begin{proof}
\rulename{Close} selects $\ell=\nu(\Pi)$; its matrix at $\ell$ is the
equation $\ell=t$, which does not contain $y$.  In either mode the
assumption of the premise and the instance asked for are both $\ell=t$, the
mode being by progress only when $\ell\equiv t\equiv 0$.
$\rulename{Use}_\nu$ of the assumption clause $[\ ;\,v.\,\ell=t]$ derives
the instance.
\end{proof}

The reading of $\mathsf V_t$ is $\exists x\,\bigl(x=t\to\forall y\,(x=t)\bigr)$,
a tautology.  It does not assert that an opened witness equals $t$, and its
opened clause $[b=t;\,y.\,b=t]$ gives no equation $b=t$; it is the selector
that substitutes $t$ for $b$ in the exported formula.

\begin{lemma}[Specialization at a term]
\label{wit:lem:v-spec}
Let $[A;\,y.\,M(y)]$ be an occurrence in $\Gamma$ and $t$ a term over $\Pi$,
and suppose a derivation of $A$ can be constructed over the scope $\Pi,b$
and the context $\Gamma,[b=t;\,y.\,b=t]$ for a fresh $b$.  Then $\mathsf V$
derives $\Pi;\Gamma\proves\nf(M(t))$ by four inferences above that
derivation: the two of \cref{wit:lem:v-packet}, one \rulename{Export} and one
$\rulename{Use}_\nu$.  For a clause without antecedent no derivation of $A$
is needed.
\end{lemma}

\begin{proof}
Derive $\mathsf V_t$ by \cref{wit:lem:v-packet} and apply \rulename{Export} at
the fresh $b$.  The opened matrix is $b=t$, so $\widehat B\equiv B_0\equiv
B_1\equiv(b=t)$ and $r^{\mathsf V}_B\equiv t$.  In the branch $b$ is the
newest parameter; $\rulename{Use}_\nu$ of the occurrence, with the given
derivation of $A$, gives $\nf(M(b))$.  The exported formula is
$\nf(\nf(M(b))\subst{t}{b})\equiv\nf(M(t))$ by \cref{ctr:lem:nf}, and it does
not contain $b$ because neither $M$ nor $t$ does.
\end{proof}

The derivation of $A$ must be constructed in the extended scope: a
derivation of $A$ over $\Pi$ is in general not a derivation over $\Pi,b$,
since its uses of clauses and its introductions select $\nu(\Pi)$, which is
no longer the newest parameter (\cref{wit:ex:weakening}).  The lemma is
therefore applied to persistent constructions.  Applied to $\mathsf{Rf}$, it
gives $b=b$ in the branch and $t=t$ as the exported formula: reflexivity at
every term, in every scope, by four inferences and without a term argument.
This is the \emph{persistent reflexivity} used below.

\begin{remark}[The obstruction does not apply]
\label{wit:rem:v-obstruction}
At the closed root both $\rulename{Use}_\nu$ and \rulename{Close} use $0$,
and in every scope both use the same parameter.  This does not contradict
\cref{wit:thm:finite-range}, which concerns the rules of $\Nat$, whose
elimination \rulename{DE} substitutes nothing.  The validity invariant of
that proof, which sends every parameter to $0$, is not preserved by the
vacuous case of \rulename{Export}: the branch is valid at $b\mapsto 0$,
while the conclusion is the branch formula at the term $r$, whose value can
be any number.  The restriction of the targets is compensated by a discharge
that substitutes a term read off the major premise.
\end{remark}

\subsection{Conditional copy and the removal of both terms}
\label{wit:sec:vac-copy}

The constructions of \cref{wit:sec:copy-prep} transfer to $\mathsf V$, with
every use of a clause at a term other than the newest parameter replaced by
\cref{wit:lem:v-spec} and every reflexivity by the persistent reflexivity.
The ordinary elimination of $\D{x}{y}B$ is \rulename{Export} with a branch
conclusion $F$ that does not contain $b$: then $\nf(F\subst{r}{b})\equiv F$
for every selected $r$, so the extended selector does no harm.  The tower
$T_0=\top_{*}$ is the persistent reflexivity at $0$, and $T_{d+1}$ is
derived by \rulename{Close} with an unused witness, whose assumption in
either mode is $T_d$, the instance asked for.  The closure $q_yA$ is
introduced as in \cref{wit:sec:copy-prep}, with the persistent reflexivity
for $A(0)$ when $A(0)\equiv\top_{*}$, and $A(t)$ is obtained from $q_yA$ and
$A(0)$ by opening $q_yA$ and applying \cref{wit:lem:v-spec} to its clause at
$t$, with $A(0)$ as the antecedent.

\begin{lemma}[Conditional copy in $\mathsf V$]
\label{wit:lem:v-copy}
For every formula $H(x)$ with a distinguished variable and every term $t$
over $\Pi$, $\mathsf V$ derives $\Delta_t(H)$ of \cref{wit:thm:copy} in every
scope and context, by a construction that is carried out anew in each
scope.
\end{lemma}

\begin{proof}
The proof of \cref{wit:thm:copy} applies with two changes.  The outer
introduction of the closure \eqref{wit:eq:copy-closure} is by cases, since
its instance at $0$ is the tower $T_{d+1}$, where $d$ is the depth of $H$.
In the branch for $\Suc z$ the conclusion normalizes to $\Delta_z(H)$, and
$z$ is the newest parameter, so \rulename{Close} for $\Delta_z(H)$ selects
$z$ and both of its instances are $H(z)$.  If $H(z)\not\equiv\top_{*}$ the
rule is by cases and $\rulename{Use}_\nu$ of the assumption gives the
instance; if $H(z)\equiv\top_{*}$ the instance is the persistent reflexivity
at $0$, the progress assumption being unused; this is the first change.
Then open the closure at a fresh $d'$; its clause is
$[T_{d+1};\,v.\,G_v\Delta_{P(v)}(H)]$.  The second change: in place of the
use of this clause at $\Suc t$, apply \cref{wit:lem:v-spec} at $\Suc t$,
with the persistent tower $T_{d+1}$ as the antecedent, constructed in the
scope of its use.  The result is
$\nf(G_{\Suc t}\Delta_{P(\Suc t)}(H))\equiv\Delta_t(H)$, which contains
neither $d'$ nor the parameter of the packet, and the ordinary elimination
discharges $d'$.  The construction is explicit and makes no recursive call
to a translation of derivations.
\end{proof}

Let $\mathsf L$ be the liberal calculus of \cref{wit:thm:witness-removal}:
\rulename{Use} at every term, \rulename{Close} at every witness, and
\rulename{Export} with $r_B$.

\begin{theorem}[Both terms removed]
\label{wit:thm:v-equiv}
For every scope $\Pi$, clause context $\Gamma$ and formula $A$ the calculi
$\mathsf V$, $\mathsf K$ and $\mathsf L$ derive the same judgments
$\Pi;\Gamma\proves A$.  The transformations are effective.
\end{theorem}

\begin{proof}
$\mathsf K$ and $\mathsf L$ derive the same judgments by
\cref{wit:thm:witness-removal}.

\emph{From $\mathsf V$ to $\mathsf L$}, by induction on the derivation.
$\rulename{Use}_\nu$ is \rulename{Use} at $\nu(\Pi)$, \rulename{Close} is
\rulename{Close} at the witness $\nu(\Pi)$, and an \rulename{Export} whose
selector is not in the vacuous case is an \rulename{Export} of $\mathsf L$
with the same conclusion, because then $r^{\mathsf V}_B\equiv r_B$.  In the
vacuous case the induction hypothesis gives derivations in $\mathsf L$ of
$\Pi;\Gamma\proves\D{x}{y}B$ and of
\[
 \Pi,b;\ \Gamma,[b=r;\,y.\,b=r]\proves F,
\]
and $G\defeq\nf(F\subst{r}{b})$ does not contain $b$.  (1) Every use of the
clause $[b=r;\,y.\,b=r]$ concludes its antecedent $b=r$; replacing each such
node by the derivation of its premise, and deleting the occurrence from all
contexts, gives $\Pi,b;\Gamma\proves F$ in $\mathsf L$; no rule of
$\mathsf L$ depends on the length of the context.  (2) For a fresh $c$, the
injective renaming $b\mapsto c$ into the scope $\Pi,b,c$ and weakening by a
clause (\cref{wit:lem:liberal-structural}) give
$\Pi,b,c;\Gamma,\mathcal{C}\proves F\subst{c}{b}$ for every clause
$\mathcal C$ over $\Pi,b,c$.  (3) Over $\Pi,b$ the name
$\D{x}{y}\bigl(C(y,0,x)=C(y,0,r)\bigr)$ is derived by \rulename{Close} at the
witness $r$: the instance at $0$ is $\top_{*}$, the mode is progress, and
the instance asked for normalizes to $r=r$, a \rulename{Use} of
$\mathsf{Rf}$ at $r$.  \rulename{Export} of the name at $c$ has the rows
$\top_{*}$ and $c=r$, selects $r$, and with the branch of (2), $\mathcal C$
being the opened clause, concludes
$\nf(F\subst{c}{b}\subst{r}{c})\equiv G$, which contains neither $c$ nor
$b$.  (4) Weakening by the clause $[b=r;\,y.\,b=r]$ and \rulename{Export} of
the given major premise at $b$ with this branch conclude $G$ over $\Pi$,
since $G$ does not contain $b$.
% ed.: the package obtains V within L from the soundness of V and the
% completeness of L for T_R (wit:prop:liberal); the vacuous case is derived
% here directly in L, which needs neither.

\emph{From $\mathsf L$ to $\mathsf V$}, by induction on a derivation $\pi$
of $\mathsf L$, with the statement of the proof of
\cref{wit:thm:witness-removal}: for every injective map $\theta$ of the
source parameters into a target scope and every map of the source clause
occurrences to target occurrences whose clauses are literally the
normalized $\theta$-images, $\mathsf V$ derives the normalized
$\theta$-image of the conclusion; further parameters and clauses are
allowed, and $\mathsf{Rf}$ is sent to $\mathsf{Rf}$.  The images of $t$, $B$
and so on are written with the same letters.

\rulename{Use} at a term $t$: apply \cref{wit:lem:v-spec} at $t$ to the image
occurrence; the derivation of the antecedent in the extended situation is
the induction hypothesis for the source derivation of the antecedent, with
the same maps, the packet parameter and its clause being extras.

\rulename{Export}: translate the major premise and the branch, the source
eigenvariable sent to the fresh target parameter and the opened clause to
the opened clause.  An injective renaming preserves and reflects the
conditions $B_0\equiv\top_{*}$, $B_1\equiv(b=r)$ and $\widehat B\equiv B_0$,
by \cref{ctr:lem:nf}.  If the source selector succeeds, so does the first
case of $r^{\mathsf V}$, with the image of the source term.  If neither case
applies, both fall back.  If the vacuous case applies, the source selector
falls back, because $B_0$ is the equation $b=r$ and not $\top_{*}$; the
source branch conclusion therefore does not contain the source
eigenvariable, its image $F$ does not contain $b$, and
$\nf(F\subst{r}{b})\equiv F$.  In each case the target conclusion is the
image of the source conclusion.

\rulename{Close} at a witness $t$ with the matrix $B$ and the premise
$\Gamma,[\ ;\,v.\,\nf(B(t,\eta))]\proves\nf(B(t,\Suc z_s))$: put
$H(x)\defeq q_yB(x,y)$ and follow the steps (a) to (e) of the proof of
\cref{wit:thm:witness-removal}.
\begin{enumerate}[label=\textup{(\alph*)}]
\item Derive $\Delta_t(H)$ by \cref{wit:lem:v-copy} and open it at a fresh
 $b$, with the clause $[H(t);\,v.\operatorname{Sel}(v,H(t),H(b))]$.
\item Apply \rulename{Close} to $\D{x}{y}B$; the newest parameter is $b$, so
 the rule selects $b$; let $z$ be its eigenvariable.
\item In the premise derive $H(t)$ by a \rulename{Close} with an unused
 witness and a further eigenvariable $w$.  Its test is the source test at
 $t$, so its assumption clause is literally the image of the source
 assumption with $z_s\mapsto w$; the induction hypothesis for the source
 premise, with these maps, gives $\nf(B(t,\Suc w))$.  This construction is
 carried out anew wherever $H(t)$ is needed.
\item Apply \cref{wit:lem:v-spec} at $1$ to the clause of (a), with the
 construction (c) of $H(t)$, made in the scope of the packet, as the
 antecedent: this gives $\operatorname{Sel}(1,H(t),H(b))\equiv H(b)$.
\item Open $H(b)$ at a fresh $d$, with the clause $[B(b,0);\,y.\,B(b,y)]$,
 and apply \cref{wit:lem:v-spec} at $\Suc z$ to it.  Its antecedent
 $B(b,0)$ is $\rulename{Use}_\nu$ of the assumption of the \rulename{Close}
 of (b) if that rule is by cases, and the persistent reflexivity at $0$ if it
 is by progress, since then $B(b,0)\equiv\top_{*}$.  The result
 $\nf(B(b,\Suc z))$ does not contain $d$, and the ordinary elimination
 discharges $d$.
\end{enumerate}
The \rulename{Close} of (b) yields $\D{x}{y}B$, which does not contain $b$,
and the ordinary elimination discharges the opening of (a).  The recursive
calls are made for the premise of the source \rulename{Close} and for the
derivation of a source antecedent, both proper subderivations, possibly
several times and each time in the situation where the result is used; the
packets, the towers and the conditional copy are explicit constructions,
which are larger than the source premise but are not sent through the
translation.  The tests at $t$ and at $b$ may differ, as in
\cref{wit:rem:mode}.
\end{proof}

\begin{corollary}[The bridge for $\mathsf V$]
\label{wit:cor:v-bi}
$\mathsf V$ is a sound and complete presentation of $T_R$, and its ordinary
reading is bi-interpretable with $\iPRA$, with whole domains, equality as
equality, the graph $x=y$ in both numerical composites and a provable
equivalence of every formula with its double translation.  Hence
$\mathsf V$ and $\iPRA$ are equiconsistent.
\end{corollary}

\begin{proof}
\Cref{wit:thm:v-equiv,wit:thm:k-bi}.  $\mathsf V$ has the formulas, terms
and readings of $\mathsf K$; only derivations change, so the formula round
trips of \cref{wit:thm:k-bi} hold unchanged.
\end{proof}

\begin{remark}[The two packages]
\label{wit:rem:spec-packages}
\pkg{export\_specialization} proves \cref{wit:thm:x-newest} and
\cref{wit:prop:x-raw} and the count of nodes; it keeps the witness of the
introduction and \rulename{Refl} unrestricted and asks whether both
restrictions can be combined, and whether arbitrary reflexivity survives a
restriction of the reflexivity clause.  On the recursor syntax with
normalization, \pkg{vacuous\_specialization} answers both:
\cref{wit:thm:v-equiv} restricts the witness, the target of every use,
including the uses of $\mathsf{Rf}$, and reflexivity at every term is the
persistent reflexivity.  The conditional copy of \cref{wit:thm:copy} with
the restricted introduction is from \pkg{conditional\_copy}; the new
ingredient is the vacuous case of the selector.  The package proves the
inclusion of $\mathsf V$ in $\mathsf L$ semantically; the proof above
derives the vacuous case in $\mathsf L$.  On iterator terms no term
satisfies $P(\Suc z)\equiv z$ by computation, and the combination there is
open (\cref{wit:q:raw-three}).
\end{remark}

% ---------------------------------------------------------------------
\section{Introduction with a reconstructed matrix}
\label{wit:sec:recon}

This section draws on \pkg{reconstructed\_introduction}.

The calculi of the preceding sections remove the witness of the
introduction and keep its matrix.  The calculus below keeps the witness and
removes the matrix.  The rule $\rulename{DI}_{\mathrm{w}}$ of
\cref{kern:def:W} is given a matrix, a witness and a derivation of its
premise.  Its reconstructed form is given the witness, the eigenvariable, the
assumption of the premise and the derivation; the matrix is computed from the
conclusion of that derivation by abstraction, after the derivation has been
checked.  This is the counterpart for the introduction of the rule
\rulename{AQE} of \cref{kern:sec:aqe}, which computes the matrix of the
elimination of a closure from its conclusion.  The calculus is a different
presentation of the same consequence relation, not the deletion of an
arithmetic principle: its equivalence with $\mathsf{W}$ and $\Nat$
(\cref{wit:thm:recon}) refines \cref{kern:thm:witness}.

\subsection{The rule and the calculus}
\label{wit:sec:recon-rule}

$\operatorname{Rep}_{t}^{s}$ is the replacement in terms and native formulas
of \cref{kern:sec:abstr}: one traversal from the root in which every term node
identical with $t$ is replaced by $s$ and not traversed further, inserted
copies of $s$ are not traversed, formula nodes are never compared with $t$,
and bound variables are renamed away from $\FV(s)\cup\FV(t)$.

\begin{definition}[Reconstructed introduction; the calculus
$\mathsf{W}_{\mathrm r}$]
\label{wit:def:recon}
Let $\Gamma$ be a context, $w$ a term, $z$ a variable not free in $\Gamma$ or
in $w$, and $A$ and $C$ formulas.  Let $x$ and $y$ be distinct variables,
different from $z$ and not free in $\Gamma$, $w$, $A$ or $C$, and let the
bound variables of $C$ be renamed away from $\FV(w)\cup\{x,y,z\}$.  The
\emph{reconstructed body} is
\[
 \rho_{w,z}(C)\defeq
 \begin{cases}
  \operatorname{Rep}_{\Suc z}^{y}(C),
   & w\equiv 0\text{ and }C\text{ is a literal equation}
     \quad(\text{mode of progress}),\\[0.5ex]
  \operatorname{Rep}_{\Suc z}^{y}\bigl(\operatorname{Rep}_{w}^{x}(C)\bigr),
   & \text{otherwise}\quad(\text{mode of cases}).
 \end{cases}
\]
The rule of \emph{reconstructed introduction} is
\[
 \infrule{\Gamma,\ A\proves C}
         {\Gamma\proves\D{x}{y}\rho_{w,z}(C)}\ \rulename{DI}_{\mathrm{r}},
\]
with the side conditions that $z$ is not free in $\rho_{w,z}(C)$ and that
\[
 A\equiv\rho_{w,z}(C)\subst{z}{y}\ \text{ in the mode of progress},\qquad
 A\equiv\rho_{w,z}(C)\subst{w}{x}\subst{0}{y}\ \text{ in the mode of cases}.
\]
The data of an instance are $w$, $z$, $A$ and the derivation of the premise;
an instance carries no matrix.  $\mathsf{W}_{\mathrm r}$ has the seven rules
\rulename{Hyp}, \rulename{I0}, \rulename{IS}, \rulename{Tr}, \rulename{QE},
\rulename{DE} and $\rulename{DI}_{\mathrm{r}}$.
% ed.: the package states the calculus with transport in equation templates in place of Tr; the two forms are interderivable over I0 (wit:rem:recon-package).
\end{definition}

The replacement maps an equation to an equation and keeps the root $\Q{}$ or
$\D{}{}$ of a compound formula, so $\rho_{w,z}(C)$ is a literal equation
exactly when $C$ is one.  The mode of progress is therefore selected exactly
when the witness is the literal $0$ and the reconstructed matrix is a literal
equation, as in $\rulename{DI}_{\mathrm{w}}$.  In that mode $x$ is not free in
$\rho_{w,z}(C)$, since it is not free in $C$, and the coordinate $x$ of the
conclusion is unused.  Since $w$ contains no $y$, the two substitutions in the
condition of the mode of cases may be performed in either order.

The operation $\operatorname{Rep}$ is total and deterministic, so
$\rho_{w,z}(C)$ is a function of $w$, $z$ and $C$; another choice of the
placeholders $x$, $y$ changes only the names of the bound variables of the
conclusion.  The seed $A$ is part of the instance and is fixed before the
premise is derived; the derivation of the premise may use $A$ through
\rulename{Hyp}.  The body is computed only after $C$ is known and is then
compared with the seed.  A checker that infers the conclusion of every node
therefore needs no matrix: it checks the derivation of the premise in the
context $\Gamma,A$, computes $\rho_{w,z}(C)$, tests that $z$ is not free in
it and compares $A$ with the required instance.  Neither the context of the
premise nor the seed is chosen retrospectively.

The test that $z$ is not free in $\rho_{w,z}(C)$ is not implied by the
reconstruction.  With $w\defeq 1$ and the premise $\Gamma,A\proves z=z$, the
two abstractions leave $z=z$ unchanged, and without the test the rule would
conclude $\D{x}{y}(z=z)$, which still contains the discharged eigenvariable.

\begin{lemma}[Syntax of the reconstruction]
\label{wit:lem:recon-syntax}
\leavevmode
\begin{enumerate}[label=\textup{(\roman*)}]
\item Let $H(y)$ be a formula with the designated variable $y$, let $z$ be a
 variable different from $y$ and not free in $H(y)$, and let the bound
 variables of $H$ be renamed away from $\{y,z\}$.  Then
 $\operatorname{Rep}_{\Suc z}^{y}\bigl(H(\Suc z)\bigr)\equiv H(y)$.
\item In the situation of \cref{wit:def:recon}, in the mode of progress
 $\rho_{w,z}(C)\subst{\Suc z}{y}\equiv C$, and in the mode of cases
 $\rho_{w,z}(C)\subst{w}{x}\subst{\Suc z}{y}\equiv C$.
\item Let $B(x,y)$ be a formula, $t$ a term in which $x$ and $y$ are not
 free, and $z$ a variable not free in $B(x,y)$ or in $t$, with the bound
 variables of $B$ renamed away from $\FV(t)\cup\{x,y,z\}$.  Put
 \[
  n\defeq 1+\bigl|B(x,\Suc z)\bigr|,\qquad
  K\defeq\It{0}{t}{u}{\Suc^{n}0},
 \]
 with $|\cdot|$ the number of nodes of \cref{kern:lem:fmarker}.  Then
 $K\not\equiv 0$, $K=t$ is an instance of \rulename{I0},
 $\FV(K)=\FV(t)$, so that $z$ is not free in $K$, and
 \[
  \rho_{K,z}\bigl(B(K,\Suc z)\bigr)\equiv B(x,y),
 \]
 computed in the mode of cases.
\end{enumerate}
\end{lemma}

\begin{proof}
(i) The variable $z$ is not free in $H(y)$, so the free occurrences of $z$ in
$H(\Suc z)$ are exactly the arguments of the inserted copies of $\Suc z$.  A
term node identical with $\Suc z$ is the successor of a free occurrence of
$z$, hence the successor node of an inserted copy.  Conversely, a node that
properly contains an inserted copy has more than two nodes and is not
identical with $\Suc z$, so the traversal reaches every inserted copy,
replaces it by $y$ and does not descend further.  Nothing else is replaced,
and the result is $H$ with $y$ at the positions of its designated variable,
that is $H(y)$.

(ii) By \cref{kern:lem:abstr}(i), applied to $C$ and the source $\Suc z$,
the replaced occurrences give a template $H$ with $C\equiv H(\Suc z)$ and
$\operatorname{Rep}_{\Suc z}^{y}(C)\equiv H(y)$.  Since $y$ is not free in $C$,
the free occurrences of $y$ in $H(y)$ are exactly the replaced positions, and
$H(y)\subst{\Suc z}{y}\equiv H(\Suc z)\equiv C$; no variable of $\Suc z$ is
captured, because the bound variables of $C$ avoid $z$.  This is the mode of
progress.  In the mode of cases put $C_1\defeq\operatorname{Rep}_{w}^{x}(C)$.
By \cref{kern:lem:abstr}(i) there is a template $H_1$ with $C\equiv H_1(w)$ and
$C_1\equiv H_1(x)$, and since $x$ is not free in $C$, the free occurrences of
$x$ in $C_1$ are exactly the replaced positions.  By the argument of the mode
of progress applied to $C_1$, in which $y$ is not free,
$\rho_{w,z}(C)\subst{\Suc z}{y}\equiv C_1$.  The second abstraction does not
remove an occurrence of $x$: a replaced node is identical with $\Suc z$, and
its only variable is $z\neq x$.  Hence the free occurrences of $x$ in
$\rho_{w,z}(C)$ are those of $C_1$.  Since $w$ contains no $y$ and $\Suc z$
contains no $x$, the two substitutions commute, and
\[
 \rho_{w,z}(C)\subst{w}{x}\subst{\Suc z}{y}
 \equiv\rho_{w,z}(C)\subst{\Suc z}{y}\subst{w}{x}
 \equiv C_1\subst{w}{x}\equiv H_1(w)\equiv C .
\]

(iii) $K$ is an iterator, so $K\not\equiv 0$ and $\rho_{K,z}$ is computed in
the mode of cases; $K=t$ is \rulename{I0}, and $\FV(K)=\FV(t)$ because the
body $\Suc^{n}0$ is closed.  Apply \cref{kern:lem:fmarker} to the template
$C(v)\defeq B(v,\Suc z)$, the term $p\defeq t$ and $a\defeq x$; its bound
$1+|C|$ is $n$.  It gives
$\operatorname{Rep}_{K}^{x}\bigl(B(K,\Suc z)\bigr)\equiv B(x,\Suc z)$.  The
variable $z$ is different from $x$ and $y$ and not free in $B(x,y)$, so (i)
with $H(y)\defeq B(x,y)$ gives
$\operatorname{Rep}_{\Suc z}^{y}\bigl(B(x,\Suc z)\bigr)\equiv B(x,y)$.  The
placeholders $x$, $y$ are admissible: they are not free in $t$, $z$ or $K$,
hence not in $B(K,\Suc z)$.
\end{proof}

Without the marker the matrix need not be recovered.  If $B(x,y)$ is $x=1$ and
the witness is $1$, the conclusion of the premise is $1=1$; the abstraction of
every occurrence of $1$ gives $x=x$, and the rule concludes $\D{x}{y}(x=x)$,
not $\D{x}{y}(x=1)$.  The proof below therefore introduces $\D{x}{y}B$ at the
marker $K$ of \cref{wit:lem:recon-syntax}(iii), whose occurrences in
$B(K,\Suc z)$ are exactly the copies at the positions of $x$, and moves the
derivation of the premise from $t$ to $K$ by transport.

\begin{lemma}[Instances of $\rulename{DI}_{\mathrm{r}}$]
\label{wit:lem:recon-instance}
Every instance of $\rulename{DI}_{\mathrm{r}}$ is an instance of
$\rulename{DI}_{\mathrm{w}}$ of \cref{kern:def:W}, with the matrix
$\rho_{w,z}(C)$, the witness $w$, and the same premise and conclusion.
\end{lemma}

\begin{proof}
Let $B(x,y)\defeq\rho_{w,z}(C)$.  The two rules select the same mode, by the
remark after \cref{wit:def:recon}.  In the mode of progress $w\equiv 0$, $B$ is
a literal equation in which $x$ is not free, and the premise of
$\rulename{DI}_{\mathrm{w}}$ is $\Gamma,B(0,z)\proves B(0,\Suc z)$, that is
$\Gamma,\rho_{w,z}(C)\subst{z}{y}\proves\rho_{w,z}(C)\subst{\Suc z}{y}$; by the
side condition and \cref{wit:lem:recon-syntax}(ii) this is
$\Gamma,A\proves C$.  In the mode of cases the premise of
$\rulename{DI}_{\mathrm{w}}$ is $\Gamma,B(w,0)\proves B(w,\Suc z)$, which is
$\Gamma,A\proves C$ by the side condition and
\cref{wit:lem:recon-syntax}(ii).  The variable $z$ is not free in $\Gamma$, in
$w$ or in $\D{x}{y}B$, which is the condition on the eigenvariable of
$\rulename{DI}_{\mathrm{w}}$.
\end{proof}

\subsection{Equivalence with the native calculus}
\label{wit:sec:recon-equiv}

Let $r(A)$ be the binder rank of \cref{base:rem:single-binder-size}.  For
$r\ge 0$ consider the following statements about $\mathsf{W}_{\mathrm r}$,
for every context $\Gamma$.
\begin{itemize}
\item[$(\mathrm T_r)$] If $F(v)$ is a formula template with $r(F)\le r$,
 $\Gamma\proves_{\mathsf{W}_{\mathrm r}}s=t$ and
 $\Gamma\proves_{\mathsf{W}_{\mathrm r}}F(s)$, then
 $\Gamma\proves_{\mathsf{W}_{\mathrm r}}F(t)$.
\item[$(\mathrm J_r)$] If $B(x,y)$ is a matrix with $r(B)\le r$, $t$ a term,
 $z$ a variable not free in $\Gamma$, in $t$ or in $\D{x}{y}B$, and
 $\Gamma,B(t,0)\proves_{\mathsf{W}_{\mathrm r}}B(t,\Suc z)$, then
 $\Gamma\proves_{\mathsf{W}_{\mathrm r}}\D{x}{y}B(x,y)$.
\end{itemize}
$(\mathrm J_r)$ is the admissibility of \rulename{DI} of $\Nat$ at matrices
of rank at most $r$, and $(\mathrm T_r)$ is the admissibility of transport in
templates of rank at most $r$.

\begin{lemma}[Structural rules and the introductions]
\label{wit:lem:recon-rules}
\leavevmode
\begin{enumerate}[label=\textup{(\roman*)}]
\item Weakening, injective renaming of free variables and cut are
 admissible in $\mathsf{W}_{\mathrm r}$ by induction on derivations.
\item $\mathsf{W}_{\mathrm r}$ derives reflexivity, symmetry and transport in
 an equation template; in particular $(\mathrm T_0)$ holds.
\item \rulename{QI} of $\Nat$ is derivable in $\mathsf{W}_{\mathrm r}$ by a
 derivation that uses \rulename{DI}$_{\mathrm r}$, \rulename{DE} and
 \rulename{Hyp} only.
\item $(\mathrm T_r)$ implies $(\mathrm J_r)$.
\item $(\mathrm T_r)$ and $(\mathrm J_r)$ imply $(\mathrm T_{r+1})$.
\end{enumerate}
\end{lemma}

\begin{proof}
(i) The proof of \cref{kern:lem:struct} applies.  The side conditions of
$\rulename{DI}_{\mathrm{r}}$ are the freshness of $z$, the test whether $w$ is
the literal $0$, the test whether $C$ is an equation, the matches of the
replacements, the test that $z$ is not free in the body and the comparison of
$A$ with an instance of the body.  An injective renaming $\theta$ of free
variables, with the placeholders chosen outside its range, preserves and
reflects identity of syntax, so it commutes with $\operatorname{Rep}$ and
$\rho_{\theta w,\theta z}(\theta C)\equiv\theta\,\rho_{w,z}(C)$; it maps a
variable to a variable, so it preserves the literal $0$ and the root of a
formula.  Hence it maps an instance to an instance with the renamed data.
Weakening, after the eigenvariables are renamed away from the added formulas,
and cut change neither $w$, $z$, $A$ nor the conclusion $C$ of the premise of
an instance: cut replaces the derivation of a retained assumption and leaves
every formula of the instance unchanged.  So both keep every side condition.
No substitution into a derivation is used.

(ii) $\mathsf{W}_{\mathrm r}$ contains \rulename{I0} and \rulename{Tr};
\cref{kern:lem:local}(iii).

(iii) Let $\Gamma,E(z)\proves_{\mathsf{W}_{\mathrm r}}E(\Suc z)$, with $E$ a
literal equation and $z$ not free in $\Gamma$ or in $\Q{y}E(y)$.  Apply
$\rulename{DI}_{\mathrm{r}}$ with the witness $0$, the eigenvariable $z$ and
the seed $E(z)$.  The conclusion $E(\Suc z)$ is a literal equation and the
witness is $0$, so the mode is progress, and by
\cref{wit:lem:recon-syntax}(i) the body is
$\operatorname{Rep}_{\Suc z}^{y}\bigl(E(\Suc z)\bigr)\equiv E(y)$.  The variable
$z$ is not free in $E(y)$, and the seed is $E(y)\subst{z}{y}\equiv E(z)$.  The
rule concludes $\Gamma\proves\D{x}{y}E(y)$, with $x$ unused.  \rulename{DE}
with a fresh eigenvariable $b$ and the second premise
$\Gamma,\Q{y}E(y)\proves\Q{y}E(y)$, an instance of \rulename{Hyp}, gives
$\Gamma\proves\Q{y}E(y)$.

(iv) Assume $(\mathrm T_r)$, and let $B$, $t$, $z$ be as in $(\mathrm J_r)$
with a derivation $\pi$ of $\Gamma,B(t,0)\proves B(t,\Suc z)$.  Rename the
bound variables of $B$ away from $\FV(t)\cup\{z\}$ and the bound variables
$x$, $y$ of the conclusion away from $\FV(\Gamma)\cup\FV(t)\cup\{z\}$, so that
they serve as the placeholders of \cref{wit:def:recon}, and let $K$ be the
marker of
\cref{wit:lem:recon-syntax}(iii) for $B$ and $t$.  Put
$\Gamma'\defeq\Gamma,B(K,0)$.  Over $\Gamma'$: \rulename{Hyp} gives $B(K,0)$;
\rulename{I0} gives $K=t$; $(\mathrm T_r)$ with the template $B(v,0)$, of rank
$r(B)\le r$, gives $B(t,0)$; $\pi$ weakened to $\Gamma',B(t,0)$ and cut with
$B(t,0)$ give $B(t,\Suc z)$, by (i); symmetry gives $t=K$, and $(\mathrm T_r)$
with the template $B(v,\Suc z)$ gives $B(K,\Suc z)$.  Thus
\[
 \Gamma,\ B(K,0)\proves_{\mathsf{W}_{\mathrm r}}B(K,\Suc z).
\]
Apply $\rulename{DI}_{\mathrm{r}}$ with the witness $K$, the eigenvariable $z$
and the seed $B(K,0)$.  The variable $z$ is not free in $\Gamma$ or in $K$.
By \cref{wit:lem:recon-syntax}(iii) the mode is cases and the body is $B(x,y)$,
in which $z$ is not free, and the seed is
$B(x,y)\subst{K}{x}\subst{0}{y}\equiv B(K,0)$.  The rule concludes
$\Gamma\proves\D{x}{y}B(x,y)$.  The mode is cases also when $B$ is a literal
equation and $t\equiv 0$, because $K\not\equiv 0$.

(v) Assume $(\mathrm T_r)$ and $(\mathrm J_r)$, and let $F(v)$ be a template
with $r(F)\le r+1$, $\Gamma\proves s=t$ and $\Gamma\proves F(s)$.  The bound
variables of $F$ are chosen fresh for the other data.  If $F$ is a literal
equation, (ii) applies.

\emph{$F(v)\equiv\Q{y}B(v,y)$}, with $r(B)\le r$.  Let $z$ be a variable not
free in $\Gamma$, $s$, $t$ or $B$.  Over $\Gamma,B(t,0)$, by weakening: $s=t$,
symmetry $t=s$, and $\Q{y}B(s,y)$; $(\mathrm T_r)$ with the template $B(v,0)$
and $t=s$ gives $B(s,0)$; \rulename{QE} with the matrix $B(s,y)$, the premise
$B(s,0)$ and the term $\Suc z$ gives $B(s,\Suc z)$; and $(\mathrm T_r)$ with
the template $B(v,\Suc z)$ and $s=t$ gives $B(t,\Suc z)$.  Hence
$\Gamma,B(t,0)\proves B(t,\Suc z)$.  Regard $B(t,y)$ as a matrix in $x$ and
$y$ in which $x$ is not used; its rank is at most $r$, and its instance at
the witness $0$ is the given premise.  $(\mathrm J_r)$ gives
$\Gamma\proves\D{x}{y}B(t,y)$, and \rulename{DE} with a fresh eigenvariable
and \rulename{Hyp} give $\Gamma\proves\Q{y}B(t,y)$, as in (iii).

\emph{$F(v)\equiv\D{x}{y}B(v,x,y)$}, with $r(B)\le r$.  Let $b$, $z$ be distinct
variables not free in $\Gamma$, $s$, $t$ or $B$, and
$\Delta\defeq\Gamma,\Q{y}B(s,b,y)$.  Over $\Delta,B(t,b,0)$ the steps of the
preceding case, with the templates $B(v,b,0)$ and $B(v,b,\Suc z)$ and the
matrix $B(s,b,y)$ of \rulename{QE}, give $B(t,b,\Suc z)$.  $(\mathrm J_r)$ with
the matrix $B(t,x,y)$ and the witness $b$ gives
$\Delta\proves\D{x}{y}B(t,x,y)$; $z$ is not free in $\Delta$, $b$ or the
conclusion.  \rulename{DE} with the major premise
$\Gamma\proves\D{x}{y}B(s,x,y)$ and the eigenvariable $b$, which is not free in
$\Gamma$, in the major premise or in $\D{x}{y}B(t,x,y)$, gives
$\Gamma\proves\D{x}{y}B(t,x,y)$.

Every formula of rank at most $r+1$ is an equation or has one of the two
forms, with a matrix of rank at most $r$.
\end{proof}

The order $(\mathrm T_0),(\mathrm J_0),(\mathrm T_1),(\mathrm J_1),\dots$ is
not circular.  $(\mathrm T_0)$ is primitive transport.  $(\mathrm J_r)$ uses
transport at rank $r$ only in the templates $B(v,0)$ and $B(v,\Suc z)$, and
one instance of $\rulename{DI}_{\mathrm{r}}$, at the marker.
$(\mathrm T_{r+1})$ uses transport and introduction at rank $r$, the primitive
\rulename{QE} and \rulename{DE}, and \rulename{Hyp}; the extraction of
$\Q{y}B$ from $\D{x}{y}B$ with $x$ unused uses no transport at all.  The
derivation of \rulename{QI} in (iii) uses no transport.  Compared with
\cref{kern:lem:aqerank} the roles are exchanged: in $\mathsf{A}_7$ the
introductions are primitive and the elimination of a closure is recovered
jointly with transport, whereas in $\mathsf{W}_{\mathrm r}$ the elimination is
primitive and the introduction by cases is recovered jointly with transport.

\begin{theorem}[Reconstructed introduction]
\label{wit:thm:recon}
$\mathsf{W}_{\mathrm r}$, $\mathsf{W}$ and $\Nat$ are equivalent: for every
context $\Gamma$ and every native formula $A$,
\[
 \Gamma\proves_{\mathsf{W}_{\mathrm r}}A
 \iff\Gamma\proves_{\mathsf{W}}A
 \iff\Gamma\proves_{\Nat}A .
\]
The translations leave terms, formulas, contexts and conclusions unchanged.
\end{theorem}

\begin{proof}
$\mathsf{W}$ and $\Nat$ are equivalent by \cref{kern:thm:witness}.

\emph{From $\mathsf{W}_{\mathrm r}$ to $\mathsf{W}$.}  The rules \rulename{Hyp},
\rulename{I0}, \rulename{IS}, \rulename{Tr}, \rulename{QE} and \rulename{DE}
are common, and an instance of $\rulename{DI}_{\mathrm{r}}$ is an instance of
$\rulename{DI}_{\mathrm{w}}$ with the same premise and conclusion
(\cref{wit:lem:recon-instance}).  So a derivation of
$\mathsf{W}_{\mathrm r}$ is a derivation of $\mathsf{W}$.

\emph{From $\Nat$ to $\mathsf{W}_{\mathrm r}$.}  By
\cref{wit:lem:recon-rules}(ii), (iv) and (v) and induction on $r$,
$(\mathrm T_r)$ and $(\mathrm J_r)$ hold for every $r$.  Translate a
derivation of $\Nat$ by induction on it.  \rulename{Hyp}, \rulename{I0},
\rulename{IS}, \rulename{Tr}, \rulename{QE} and \rulename{DE} are copied, the
premises being translated.  An instance of \rulename{QI} is replaced by the
derivation of \cref{wit:lem:recon-rules}(iii) from its translated premise, and
an instance of \rulename{DI} with the matrix $B$ by $(\mathrm J_{r(B)})$
applied to its translated premise.  Each replacement acts over the context of
the inference it replaces and introduces only fresh variables, so the
conditions on the eigenvariables of enclosing inferences are kept.
\end{proof}

\begin{corollary}[The bridge for $\mathsf{W}_{\mathrm r}$]
\label{wit:cor:recon-bridge}
$\mathsf{W}_{\mathrm r}$, $\Teq$ and $\iPRA$ are bi-interpretable with whole
numerical domains, actual equality, the identity graph $x=y$ in both numerical
composites and both formula round trips;
$\proves_{\mathsf{W}_{\mathrm r}}0=\Suc 0$ holds exactly when $\iPRA$ is
inconsistent; and weakening, substitution of terms, cut and the Leibniz rule
are admissible in $\mathsf{W}_{\mathrm r}$.
\end{corollary}

\begin{proof}
\Cref{wit:thm:recon,kern:prop:transfer}.  The calculus has the terms and
formulas of $\Nat$ and the same derivable sequents, so the translations of
\cref{base:thm:bridge} and the derivations of the round trips are used
unchanged.
\end{proof}

\begin{remark}[Substitution in $\mathsf{W}_{\mathrm r}$]
\label{wit:rem:recon-subst}
Substitution is admissible in $\mathsf{W}_{\mathrm r}$ through
\cref{wit:thm:recon}, not by a direct induction on its derivations.  Let $a$,
$b$ be distinct variables, $\Gamma$ a context containing $a=b$, and $z$ fresh.
$\rulename{DI}_{\mathrm{r}}$ with the witness $a$, the seed $a=b$ and the
premise $\Gamma,a=b\proves a=b$ by \rulename{Hyp} is in the mode of cases, has
the body $x=b$ and concludes $\D{x}{y}(x=b)$.  The substitution of $a$ for
$b$ turns the premise into $\Gamma',a=a\proves a=a$, whose body is $x=x$; the
rule then concludes $\D{x}{y}(x=x)$, whereas the substituted conclusion is
$\D{x}{y}(x=a)$.  A substitution can also turn a variable witness into the
literal $0$ and so change the mode, as in \cref{kern:rem:wsubst}.  The proof
of \cref{wit:thm:recon} therefore uses weakening, injective renaming and cut
only, which change no datum of an instance except the names of
eigenvariables.
\end{remark}

\begin{remark}[The package]
\label{wit:rem:recon-package}
\pkg{reconstructed\_introduction} calls the calculus $\mathcal S_7$ and the
rule synthesized introduction.  Its transport rule is transport in a literal
equation template, from $s=t$ and $E(s)$ to $E(t)$, in place of
\rulename{Tr}.  The two forms are interderivable over \rulename{I0}:
\cref{kern:lem:local}(iii) derives the transport in equation templates from
\rulename{Tr}; conversely the transport in equation templates derives
reflexivity, by transporting $\It{0}{t}{u}{u}=t$ along itself in the template
$v=t$, then symmetry, by transporting $s=s$ in the template $v=s$, and then
\rulename{Tr}, from $s=t$ and $h(s)=u$, by transporting $u=h(s)$ in the
template $u=h(v)$.  So \cref{wit:thm:recon} holds for both forms.  The package
proves soundness directly in $\Nat$, deriving an introduction in the mode of
progress by \rulename{QI}, \rulename{QE} and \rulename{DI} at the witness $0$;
the proof above goes through $\rulename{DI}_{\mathrm{w}}$ and
\cref{kern:thm:witness}.  Its binder rank counts nested binders instead of
occurrences of binders; the induction is the same, since the matrix of a
compound formula has smaller rank in both senses.
\Cref{wit:lem:recon-syntax}(iii) is the oversized marker of
\pkg{abstraction\_and\_refinements} (\cref{kern:lem:fmarker}) applied, before
the abstraction of $\Suc z$, to the template $B(v,\Suc z)$.  The package
contains an independent written review of the equivalence proof and of its
readable checker.  The theorem is not formalized in Lean or Rocq; the bridge
it uses is that of \cref{base:thm:bridge}, which is not proved again.  The
package's
separate exact-binding checkers of the native calculus do not implement
$\mathsf{W}_{\mathrm r}$.
\end{remark}

% ---------------------------------------------------------------------
\section{Checkers and sizes}
\label{wit:sec:sizes}

This section draws on the checker reports of \pkg{no\_witness},
\pkg{parameter\_witness}, \pkg{deterministic\_introduction},
\pkg{newest\_parameter}, \pkg{uniform\_raw}, \pkg{parameter\_intro},
\pkg{conditional\_copy}, \pkg{local\_introduction},
\pkg{equational\_transport}, \pkg{names\_and\_scopes},
\pkg{export\_specialization},
\pkg{vacuous\_specialization} and \pkg{reconstructed\_introduction}.  The checker
of \pkg{constructed\_computation} and the compact checkers of
\pkg{export\_specialization} and \pkg{vacuous\_specialization}, which check
the guarded scalar calculus and not the calculi of this chapter, belong to
another chapter, as do the exact-binding refinements of the native scalar
checker in \pkg{names\_and\_scopes}.  The exact-binding and padded sorted
checkers of the native calculus in \pkg{reconstructed\_introduction}
(313 / 357 / 590 and 298 / 336 / 573) belong to \cref{prov:ch:cert}.

All checkers are standalone Python modules that read one JSON certificate,
validate syntax, scopes and every inference, and accept or reject; none
calls a prover.  They share the encoding of syntax: the integer $0$ is
zero, negative integers are parameters counted from the outside, positive
integers are de Bruijn indices, and lists encode the constructors.  The
sizes are measured by the protocol of \cref{base:sec:measure} under
CPython 3.12.14.

\begin{longtable}{@{}p{0.36\linewidth}rrrp{0.36\linewidth}@{}}
\caption{Checkers of the chapter: significant tokens, AST nodes and
normalized bytes, with the matched control as tokens / nodes / bytes.}
\label{wit:tab:sizes}\\
\toprule
Program & Tokens & Nodes & Bytes & Matched control\\
\midrule
\endfirsthead
\toprule
Program & Tokens & Nodes & Bytes & Matched control\\
\midrule
\endhead
\bottomrule
\endfoot
\multicolumn{5}{@{}l}{\pkg{no\_witness}}\\
eight rules, $\Nat_{\nu}$ & 971 & 1078 & 2817 &
 supplied witness: 979 / 1086 / 2834\\
seven rules, $\Nat_{\mathrm{diag},\nu}$ & 924 & 1021 & 2753 &
 supplied witness: 932 / 1029 / 2770\\
\multicolumn{5}{@{}l}{\pkg{parameter\_witness}}\\
five families, $\mathsf{C}_{\beta}$ & 833 & 971 & 2947 &
 two-case raw clause checker: 837 / 977 / 2956\\
\multicolumn{5}{@{}l}{\pkg{deterministic\_introduction}}\\
$\mathsf{G}_{\mathrm d}$ & 1159 & 1339 & 3688 &
 $\mathsf{G}$: 1078 / 1252 / 3442\\
inherited native calculus, shared first traversal & 333 & 379 & 656 &
 incoming native control: 360 / 409 / 780\\
inherited guarded compact module & 267 & 303 & 506 &
 before the change: 267 / 302 / 508\\
\multicolumn{5}{@{}l}{\pkg{newest\_parameter}}\\
$\mathsf{G}_{\nu}$ & 1070 & 1243 & 3425 &
 $\mathsf{G}$: 1078 / 1252 / 3442; $\mathsf{G}_{\mathrm d}$:
 1159 / 1339 / 3688\\
\multicolumn{5}{@{}l}{\pkg{uniform\_raw}}\\
$\mathsf{G}_{\nu}$ & 1073 & 1248 & 3439 &
 $\mathsf{G}_{\mathrm d}$: 1159 / 1339 / 3688\\
\multicolumn{5}{@{}l}{\pkg{parameter\_intro}}\\
$\mathsf{P}_{\nu}$ & 630 & 701 & 1925 &
 supplied witness: 636 / 707 / 1967\\
$\mathsf{P}_{\nu}$ with a separate template check & 646 & 717 & 1985 & \\
$\mathsf{P}_{\nu}$ with fuel & 667 & 748 & 2056 &
 supplied witness, with fuel: 673 / 754 / 2098\\
\multicolumn{5}{@{}l}{\pkg{conditional\_copy}}\\
$\mathsf{K}$ & 646 & 748 & 1844 & $\mathsf{L}$: 651 / 755 / 1863\\
$\mathsf{K}_{\mathrm N}$, backward & 619 & 686 & 1907 &
 $\mathsf{P}_{\nu}$ of \pkg{parameter\_intro}: 630 / 701 / 1925\\
explicit equation at \rulename{Name} & 634 & 742 & 1835 & \\
$\mathsf{K}$ with fuel & 685 & 796 & 1995 & \\
\multicolumn{5}{@{}l}{\pkg{local\_introduction}}\\
$\mathsf{A}_{\nu}$ & 862 & 1011 & 3245 &
 $\mathsf{A}_{\nu}$ with a supplied \rulename{Close} witness:
 870 / 1022 / 3274\\
\multicolumn{5}{@{}l}{\pkg{equational\_transport}}\\
$\mathsf{A}^{=}_{\nu}$, \rulename{EqBack} & 869 & 1021 & 3263 &
 $\mathsf{A}_{\nu}$ of \pkg{local\_introduction}: 862 / 1011 / 3245\\
$\mathsf{A}^{\times}_{\nu}$, \rulename{XBack} & 867 & 1022 & 3267 &
 $\mathsf{A}_{\nu}$ of \pkg{local\_introduction}: 862 / 1011 / 3245\\
\multicolumn{5}{@{}l}{\pkg{names\_and\_scopes}}\\
$\mathsf{A}^{\mathrm n}_{\nu}$, readable module with documentation strings
 & 928 & 1087 & 3767 & none stated\\
\multicolumn{5}{@{}l}{\pkg{export\_specialization}}\\
$\mathcal X_\nu$, \rulename{Refl} primitive & 663 & 748 & 2093 &
 $\mathcal X$, \rulename{Refl} primitive: 669 / 754 / 2117\\
$\mathcal X^{\mathcal R}_\nu$, target of $\mathcal R$ read off the goal &
 654 & 731 & 2049 &
 $\mathcal X^{\mathcal R}$ of \pkg{clause\_discharge}: 653 / 729 / 2058\\
\multicolumn{5}{@{}l}{\pkg{vacuous\_specialization}}\\
$\mathsf V$ & 645 & 745 & 1830 &
 $\mathsf K$ of \pkg{conditional\_copy}: 646 / 748 / 1844\\
\multicolumn{5}{@{}l}{\pkg{reconstructed\_introduction}}\\
$\mathsf{W}_{\mathrm r}$, readable reference module with documentation
 string & 1896 & 2182 & 7639 & none stated; gzip-9 of the normalized
 source: 2249 bytes\\
\end{longtable}

The packages compare their programs with further programs of the corpus.
\pkg{parameter\_intro} states 636 / 708 / 2010 for the incoming checker with
an explicit name and 683 / 771 / 1990 for the incoming checker with an
inferred opening; its gzip sizes are 778 for its checker, 790 for the
matched control, 792 with the separate template check, and 826 and 838 for
the two variants with fuel.  \pkg{conditional\_copy} states
662 / 741 / 1950 for the checker of \pkg{uniform\_recursion}.  The compact
programs of the corpus use other input contracts and are smaller: 267
tokens, 302 nodes and 508 bytes; 256, 283 and 594; 263, 276 and 580; and,
with a validator of the public goal, 360, 409 and 780.  The three minima
256, 276 and 508 belong to different programs.  No package claims a global
minimum or a lower bound.

\paragraph{\pkg{no\_witness}.}
An introduction is the instruction \texttt{['di', matrix, premise]}; its
witness is the integer $-n$ in a scope of $n$ parameters, which is zero in
the empty scope.  The controls differ only by a supplied and validated
witness term.  Each matched pair differs by 8 tokens, 8 nodes and 17 bytes.
In the tagged syntax of the pinned programs the same change saves 3 bytes
and 1 node and adds 2 tokens.  Eighteen accepted and 33 rejected
certificates are run in two modes; the certificate for $X(2)$ checks the
part of the naming that uses numerical cases and not the predecessor or the
equivalence theorem.

\paragraph{\pkg{parameter\_witness}.}
The introduction field is a bit, decoded by the truth value of the JSON
value; it is never inserted as a term, and the witness is constructed from
the scope.  Certificates for the packet $T$ at $0$, for its closure and for
its instances at $1$ and $2$ are checked, and a certificate that uses the
eigenvariable of an introduction as its witness is rejected.  A separate
evaluator checks 447 numerical equalities of the predecessor construction,
among them the values of the predecessor at $0,\dots,7$.

\paragraph{\pkg{deterministic\_introduction}.}
The naming matrix is recognized as the literal list pattern of
$\mathrm{Nm}_r$ with $r$ valid in the parent scope.  The checker of
$\mathsf{G}_{\mathrm d}$ is larger than its control by 81 tokens, 87 nodes
and 246 bytes: the rule removes a datum from each introduction and costs
more code to recognize the exception.  The fixtures contain an introduction
that is valid before an extension of the scope, the rejected replay of the
same tree after the extension, and the repaired derivation of
\cref{wit:thm:weakening}.  The package also contains two changes to
inherited programs that do not implement $\mathsf{G}_{\mathrm d}$.  One
checks the grammar of the public goal during the first scope traversal,
using the parity of the traversal depth as the distinction between formulas
and terms; an induction on the goal shows that the shared pass succeeds
exactly when the two separate passes do and that the subsequent traces are
equal.  The other replaces $k-1$ by the bitwise complement of $-k$ in the
compact module, which saves two bytes and adds one node.

\paragraph{\pkg{newest\_parameter} and \pkg{uniform\_raw}.}
Both obtain the checker of $\mathsf{G}_{\nu}$ from the checker of
$\mathsf{G}_{\mathrm d}$ by deleting the function that selects the witness
and instantiating the matrix at $-n$.  \pkg{newest\_parameter} also contains
an unmeasured checker of the eight rules of $\Nat_{\nu}$, with certificates
for $X(\Suc t)$ at a parameter $t$ in scopes of zero, one, two and five
parameters, and a rejected replay of a closed derivation in a larger scope;
it also replays symbolic certificates of seven equation-case lemmas of the
predecessor construction.
\pkg{uniform\_raw} checks the whole closure $\D{a}{r}X(r)$ and its
instance at the successor of a compound term.

\paragraph{\pkg{parameter\_intro}.}
The recursor and the binder $\D{}{}$ both bind two variables in their last
child, so one traversal serves terms, formulas, scopes, substitution and
normalization.  The reduction of a successor counter substitutes the
previous counter first, into the outer of the two bound positions, and the
previous value second; the opposite order without a shift would capture
variables.  The witness of an introduction is the integer $1-n$ for a
context of $n$ clauses.  Omitting the witness field saves 42 bytes, 6
tokens and 6 nodes; omitting the separate template check, by
\cref{wit:lem:template}, saves 60 bytes, 16 tokens and 16 nodes.  The
variant with fuel decrements a budget at every call of the traversal and of
the inference checker; every accepted certificate has a sufficient budget,
so the two interfaces have the same theorems, and the bounded relation is
primitive recursive.  The certificates that name $1$ and $2$ without a
witness field need 1155 and 1345 units.

\paragraph{\pkg{conditional\_copy}.}
The forward checker of $\mathsf{K}$ infers every conclusion and compares
the goal with it.  Omitting the witness saves 19 bytes, 5 tokens and 7
nodes against $\mathsf{L}$.  The checker of $\mathsf{K}$ has fewer bytes
and more tokens and nodes than that of \pkg{parameter\_intro}; the backward
checker of $\mathsf{K}_{\mathrm N}$ is smaller in all three coordinates, by
18 bytes, 11 tokens and 15 nodes.  The package contains a translator from
$\mathsf{L}$ to $\mathsf{K}$ that implements the proof of
\cref{wit:thm:witness-removal} on arbitrary derivations.  Its outputs show
the growth of certificates, in inference nodes and bytes of JSON: an
introduction at the witness $1$ grows from 2 nodes and 74 bytes to 15 and
529; a conditional copy of depth one from 7 and 368 to 46 and 2527; a
derivation with a used progress assumption from 15 and 421 to 80 and 2868;
the congruence of the successor under two universal quantifiers from 25 and
1335 to 142 and 8909; a derivation in an open scope from 4 and 76 to 17 and
501.

\paragraph{\pkg{local\_introduction}.}
The input is a JSON object with exactly the keys \texttt{goal} and
\texttt{proof}; the goal is a closed formula with equations and $\D{}{}$,
validated in full, and the proof is a tagged tree with six tags: one for
each of \rulename{Use}, \rulename{Back}, \rulename{Alias}, \rulename{Open}
and \rulename{Close}, and one for the two computations, given by a counter,
a seed and a body.  Every consumed term is validated in
its scope, including an iterator body that a zero computation does not
evaluate; the source of \rulename{Back} must be the newest parameter and
not $0$; the unused premise field of a \rulename{Use} of a clause without
antecedent is ignored.  The checker of $\mathsf{A}_{\nu}$ and its control
share all code, including the diagonal test, and differ only in the
witness field of \rulename{Close}, supplied and validated in the control
and computed from the scope in $\mathsf{A}_{\nu}$, by the invariant that
the scope and the context have the same length.  The difference is 8
tokens, 11 nodes and 29 bytes.  The diagonal test is absent from the
witness-selected checkers of the Alias/Back basis; the short raw Alias/Back
checker of \pkg{alias\_back}, 773 / 916 / 2903, is smaller than both
programs of the row.  The package remeasures the compact incumbents of the
corpus under their own contract, a closed guarded public goal without a
grammar validator and a different certificate protocol: 248 / 277 / 609,
254 / 270 / 608 and 272 / 304 / 479, and 344 / 392 / 637 with the full
native goal and its alphabet validated; the envelope 248 / 270 / 479 is
attained by three different programs.  An untrusted compiler implements
the proof of \cref{wit:thm:ab-newest} for sources in the control calculus,
in which \rulename{Use} and \rulename{Back} are already at the newest
parameter, and replays every output through the checker of
$\mathsf{A}_{\nu}$; it does not parse the larger calculus
$\mathsf{A}_{\ast}$.  Eighteen source certificates are compiled and
checked, eighteen adversarial certificates are rejected, and 108 process
runs in the normal and the optimized interpreter mode pass; one open
example uses an embedding of parameters that does not preserve their order
and a map of clauses that is not contiguous.

\paragraph{\pkg{equational\_transport}.}
Both checkers have the input contract and the code of the checker of
$\mathsf{A}_{\nu}$ of \pkg{local\_introduction}, which is their matched
control, except for the backward node.  For \rulename{EqBack} it is
\texttt{['v', equation, major, minor]}; for \rulename{XBack} it is
\texttt{['v', h, major, minor]}: the checker infers $b=t$ from the major
premise and $\ell=r$ from the minor premise, requires $b$ to be the newest
parameter and $\ell\equiv h\subst{t}{b}$, and returns $r=h$ without
substituting in $r$.  The template $h$ is validated before the comparison,
because in this encoding a malformed list $[b,0]$ becomes the valid
successor $[0,0]$ after the substitution of $0$ for $b$.  \rulename{EqBack}
adds 7 tokens, 10 nodes and 18 bytes to the control and \rulename{XBack} 5
tokens, 11 nodes and 22 bytes; the gzip sizes at level 9 are 1044 for the
control, 1053 for \rulename{EqBack} and 1058 for \rulename{XBack}.  The checker of the source calculus
$\mathsf{A}_{\ast}$ measures 880 / 1035 / 3324, gzip 1048.  An untrusted
compiler implements \cref{wit:thm:eqback} from $\mathsf{A}_{\ast}$ to
$\mathsf{A}^{=}_{\nu}$, and a second one applies \cref{wit:rem:eqback-sim};
every output is replayed by the target checker.  Nineteen closed and
thirteen open translations, among them older and self-containing
\rulename{Back} sources, non-order-preserving embeddings, nested binders and
both modes, together with four open \rulename{XBack} instances, five
simulations of \rulename{EqBack} and adversarial certificates, pass in 174
checker runs under normal and optimized interpretation.  In one example 47
source inference nodes become 329 nodes of $\mathsf{A}^{=}_{\nu}$ and
148404 nodes of $\mathsf{A}^{\times}_{\nu}$, counted as occurrences in the
expanded tree.  An implementation review found that the first
$\mathsf{A}^{\times}_{\nu}$ converter traversed the ignored premise field of
a \rulename{Use} of a clause without antecedent; the final converter drops
it.

\paragraph{\pkg{names\_and\_scopes}.}
The checker of $\mathsf{A}^{\mathrm n}_{\nu}$ reads a JSON object with
exactly the keys \texttt{goal} and \texttt{proof}; the goal is a closed
well-formed formula, and the checker prints \texttt{ACCEPT} only after the
inferred conclusion equals the goal.  The encoding is that of the chapter,
with $[0,t]$ for $\Suc t$, $[n,a,h]$ for an iterator, $[s,t]$ for an equation
and $[B]$ for $\D{x}{y}B$, so that $\mathrm{Name}(t)$ is
$[[[1,t,3],t]]$; the term $t$ is validated at bound depth zero.  The proof
tags are one each for the computations, \rulename{Use}, \rulename{EqBack},
\rulename{EqAlias}, \rulename{Open}, \rulename{Close} and \rulename{Names};
both equation restrictions are checked, and the source of \rulename{EqBack}
must be the newest parameter.  The unused premise of a \rulename{Use} of a
clause without antecedent is an opaque annotation.  The module is a
readable reference implementation and is not compared with a matched
control; it is larger than the checkers of $\mathsf{A}^{=}_{\nu}$ and
$\mathsf{A}^{\times}_{\nu}$ of \pkg{equational\_transport}.  An untrusted
compiler implements the proof of \cref{wit:thm:ns-complete} from the liberal
calculus $\mathsf{A}_{\ast}$, with persistent families, the transport of
\cref{wit:thm:ns-rank} and the translations of \rulename{Close} and of full
\rulename{Alias}; it is not imported by the checker, which replays every
output.  Nineteen inherited closed fixtures, two further closed translations,
a direct certificate of $\mathrm{Name}(1)$ and thirteen open translations
(formula ranks $1$ to $3$, both modes, self-containing \rulename{Back}
replacements and an embedding that does not preserve the order of
parameters) pass in 60 runs under normal and optimized interpretation,
which also reject a compound \rulename{Alias}, false goals, malformed names
and scope escapes; a separate audit of binders checks 240 terms and 1440
conditions.  The largest generated open certificate has 24113 bytes.

\paragraph{\pkg{export\_specialization}.}
All four programs are complete closed-theorem checkers of the normalized
calculus on iterator terms: they read one JSON pair of a closed goal and a
proof, validate syntax and scopes, normalize, check every inference and
print \texttt{ACCEPT}; an exception rejects.  The matched pair with
primitive \rulename{Refl} isolates the restriction of \rulename{Use}: a use
packet loses its term field and its validation, and the checker
instantiates at the newest parameter.  The second pair keeps the
reflexivity clause at context position $0$ as in the control of
\pkg{clause\_discharge}; a use of position $0$ takes its term from the left
side of the validated equational goal, and every other use specializes at
the newest parameter.  This pair saves 9 bytes and adds one token and two
nodes; the first saves 24 bytes, 6 tokens and 6 nodes.  The validation of
the witness of \rulename{Intro} is kept in all four.  An untrusted
translator compiles the nine certificates of \pkg{clause\_discharge}
(congruence, symmetry, transitivity, induction, nested binders, the repair
at the witness $0$) to all four protocols and back; five open targets, an
injective renaming with clauses on both sides of the embedded clause, a
regression for capture under formula and iterator binders, and thirteen
negative controls are replayed.

\paragraph{\pkg{vacuous\_specialization}.}
The checker of $\mathsf V$ and its control, the checker of $\mathsf K$ of
\pkg{conditional\_copy}, read one JSON pair of a goal and a proof.  Terms
are $0$, rigid parameters $-1,-2,\dots$, positive de Bruijn indices and
$[n,a,h]$ for $R(n,a;\,i,u.\,h)$, with $u$ of index $1$ and $i$ of index $2$
in $h$; formulas are $[s,t]$ and $[B]$ for $\D{x}{y}B$.  The initial
context is $\mathsf{Rf}$, and a context of length $n$ has the parameters
$1-n,\dots,-1$, so that $1-n$ is the newest parameter, and $0$ at the root.
A use is $[j,q]$, a \rulename{Close} is $[B,q]$, an \rulename{Export} is
$[[\,],p,q]$ and opens at $-n$; the vacuous case compares the whole body of
the opened matrix with its instance at $0$.  The only differences from the
control are the deletion of the term of a use with its validation and the
enlarged selector.  The program \texttt{control\_liberal.py} of the package
is the checker of $\mathsf L$, the source of an untrusted compiler that implements the
translation of \cref{wit:thm:v-equiv} and whose outputs are rechecked; 39
focused checks and 8 command-line checks under normal and optimized
interpretation pass, among them the packet, reflexivity at a compound term,
conditional copy at three depths, a used induction hypothesis, nontrivial
embeddings of parameters and clauses, and the rejection of a surviving
fresh parameter and of rows at $0$ and $1$ that agree without a constant
matrix.

\paragraph{\pkg{reconstructed\_introduction}.}
The checker of $\mathsf{W}_{\mathrm r}$ reads one JSON pair of a closed goal
and a proof, in a locally nameless syntax with free parameters as strings,
bound indices $[\texttt{v},i]$, and the constructors $[\texttt{s},t]$,
$[\texttt{i},n,a,h]$, $[\texttt{=},s,t]$, $[\texttt{q},B]$ and $[\texttt{d},B]$,
with $y$ of index $0$ and $x$ of index $1$ below $\texttt{d}$.  It infers the
conclusion of every node and accepts when the conclusion of the root equals
the goal.  There is one record for each of the seven rules; an introduction
is $[\texttt{j},w,z,A,p]$.  The checker validates $w$ in the current scope,
extends the scope by $z$, validates the seed $A$ there, checks $p$ in the
context extended by $A$, and only then computes the body by root-first
replacement, tests that $z$ is not free in it and compares $A$ with the
required instance; the mode of progress requires the exact integer $0$ as
witness, so that Boolean and floating-point zeros are rejected, and an
equation as conclusion.  Every result is validated in the scope in which it
is returned, and \rulename{DE} rejects an escaping eigenvariable.  The
module is a readable reference implementation and is not compared with a
matched control.  It is checked by 21 positive and negative certificates in
42 runs under normal and optimized interpretation and by 1501 deterministic
checks of the syntax laws (opening and closing of binders, recovery at the
marker, separation of free labels from bound indices); among the negative
certificates are a wrong seed, the unmarked collision of
\cref{wit:sec:recon-rule} against the goal $\D{x}{y}(x=1)$, an escaping
eigenvariable and a progress step at a compound formula.

The restricted rules reduce the data of a certificate and, except for
$\mathsf{G}_{\mathrm d}$, the measured checker by a small amount.  They
increase the length of certificates, because a witness term is replaced by
a derivation of its name.

% ---------------------------------------------------------------------
\section{Further questions}
\label{wit:sec:questions}

The following items are stated by the packages without proof, are open, or
mark what the results of the chapter do not establish.

\begin{question}[Formalization of the scope-determined witness]
\label{wit:q:lean-newest}
The Lean development formalizes the restriction of the witness to the
literal $0$ or a variable, with the derivation of the names
(\cref{wit:rem:lean-complete}), and the predecessor in the equational
fragment (\cref{wit:sec:packet}).  It does not formalize the deterministic
selection of the witness, the newest parameter or $0$ in an empty scope,
with the regeneration of derivations under an extension of the scope.
Hence none of
\cref{wit:thm:newest,wit:thm:diagonal,wit:thm:R-newest,wit:thm:clause-bit,wit:thm:ab-newest,wit:thm:eqback,wit:thm:ns-complete,wit:thm:weakening,wit:thm:gd,wit:thm:guarded-uniform}
is formalized, and nothing in \cref{wit:sec:recursor,wit:sec:copy} on the
counter-aware recursor is formalized, in particular
\cref{wit:thm:recursor,wit:thm:copy,wit:thm:witness-removal}.  A
formalization of the first has to take derivations of the source with
realized contexts, or hereditary families, as its induction interface.  For
$\mathsf{G}_{\mathrm d}$ the selection with priority for a naming matrix,
the clause contexts and the guarded language are not formalized.  For the
recursor calculi a syntax with a normalizer and its bi-interpretation with
the iterator terms (\cref{ctr:lem:nf,ctr:thm:transfer}) would be needed
first;
\pkg{conditional\_copy} proposes a small interface for
\cref{wit:thm:copy} consisting of a predecessor that computes on
successors, gates, clause specialization, \rulename{Close} at the newest
parameter and the ordinary elimination, so that the theorem can be reused.
\Cref{wit:lem:template} and the translation of \cref{wit:thm:recursor} are
named by \pkg{parameter\_intro} as the next formalization task.
\pkg{local\_introduction} proposes a Lean definition of the calculus
$\mathsf{A}_{\nu}$ with ordered scopes and the induction statement of
\cref{wit:thm:ab-newest} over clause realizers, kept apart from the
invariant of closed certificates that the scope and the context have the
same length.
\end{question}

\begin{question}[Three families on raw iterator terms]
\label{wit:q:raw-three}
Can the introduction at the newest parameter be combined, on raw iterator
terms without normalization, with a three-family calculus whose equality
transport is recovered from packets?  The proofs of \cref{wit:sec:newest}
keep \rulename{Tr} as a primitive rule, and transferring them unchanged
would be circular.  \Cref{wit:sec:recursor,wit:sec:copy} answer the
question on the recursor with normalization, where $P(\Suc z)\equiv z$
holds by computation; on iterator terms no term has this property.  A
technical note of \pkg{uniform\_raw} sketches a raw calculus with a
universal reflexivity clause, a deterministic introduction, \rulename{Use}
and a discharge that evaluates, in which the opening and the evaluation
rules are recovered through the dummy major premise
$\D{x}{y}(\It{y}{0}{u}{0}=\It{y}{0}{u}{0})$ and formula transport is proved
before the names.  The note is not part of the proved results of that
package.
\end{question}

\begin{question}[Transport in terms only]
In $\mathsf{G}_{\mathrm d}$ and $\mathsf{G}_{\nu}$ the rule \rulename{Tr}
acts on arbitrary canonical formulas, and the restricted cut of
\cref{wit:lem:algebraic} uses this.  Can it be weakened to transport in
terms before the bootstrap, with formula transport proved without
unrestricted witnesses and without unrepaired weakening?
\end{question}

\begin{question}[The size of certificates]
The translations are effective.  No linear or polynomial bound on the size
of the translated certificates is claimed; the translations regenerate a
premise inside additional scopes and may duplicate it.  Determine the
growth, and whether sharing of regenerated subderivations in a checked
directed acyclic graph reduces it without enlarging the trusted rules.
\pkg{local\_introduction} asks the same for the repeated production of
persistent families in the proof of \cref{wit:thm:ab-newest}, and whether
the sharing can be checked without erasing the saving of the interface.
The translations of \cref{wit:thm:eqback} grow faster still: in an example
of \pkg{equational\_transport} 47 source nodes become 329 nodes of
$\mathsf{A}^{=}_{\nu}$ and 148404 of $\mathsf{A}^{\times}_{\nu}$, because
\cref{wit:lem:eqb-transport} recurses on the binders of the template and the
simulation of \cref{wit:rem:eqback-sim} repeats the major premise.
\end{question}

\begin{question}[Modes without a guard]
Can the scope determine the witness while the mode of the introduction is
also determined without a test on the matrix, with a direct proof of all
native consequences?  In $\Nat_{\nu}$ the progress rule \rulename{QI} is
kept; the result does not remove it, does not reduce the number of rules
and does not settle the independence of progress.
\end{question}

\begin{question}[The specialization term]
In the calculi of the chapter up to \cref{wit:sec:copy}, \rulename{QE}, or
\rulename{Use}, still takes an arbitrary term; the proofs specialize a
closure at $\Suc t$.  By
\cref{wit:thm:finite-range,wit:thm:finite-range-general} the target cannot
be restricted to a parameter or finitely many numerals together with the
witness in the identity language of $\Nat$, whose elimination substitutes
nothing.  \Cref{wit:thm:v-equiv} delivers the term by the vacuous case of an
inferred discharge instead.  Whether a calculus with the native elimination
\rulename{DE} and both restrictions, differently interpreted, is
bi-interpretable with $\iPRA$ is not decided.
\end{question}

\begin{question}[Newest specialization]
\label{wit:q:spec}
\begin{enumerate}[label=\textup{(\roman*)}]
\item \Cref{wit:thm:x-newest} is formalized against \Lean{XCPrf}, as
 \pkg{export\_specialization} proposes (\cref{wit:rem:x-lean}), without the
 node count; \cref{wit:thm:v-equiv} is not formalized, and \pkg{vacuous\_specialization} one of its compiler with
 the invariant on parameter and clause embeddings, beyond the soundness of
 the vacuous case.  Their translators are untrusted programs whose outputs
 are rechecked; their correctness is supported by finite tests only.
\item The size of translated certificates.  \Cref{wit:thm:x-newest} bounds
 the inference nodes by $n+3u$, but not the size of the normalized terms,
 the bytes of a certificate or the time of checking.  In $\mathsf V$ a use
 of a clause costs four inferences before its antecedent, and an
 introduction at an arbitrary witness invokes conditional copy.  A benchmark
 that reports proof size and checking time beside checker size is proposed.
\item Is there a larger class of cheaply recognized tautological closure
 matrices than those with $\widehat B\equiv B_0$ that keeps the discharge
 deterministic and makes the complete checker smaller?
\item Does the raw selector of \cref{wit:prop:x-raw}, with its condition on
 free variables, cost less in a complete checker than the selector on normal
 forms?  No raw checker is measured.
\item Can the newest specialization, or the vacuous selector, be built into
 the arity-coded compact checkers, so that one program checks the new
 calculus and is smaller than the compact records?  Can one module improve
 all three coordinates of the compact envelope at once under a fixed
 contract?  The packages ask, with \pkg{local\_introduction}, for a
 representation-stable measure that charges normalization, parsing and all
 decoding and admits nontrivial lower bounds; neither package proves
 minimality or a lower bound.
\item Membership-only majors.  \pkg{export\_specialization} considers the
 calculus in which the major premise of \rulename{QE} and of \rulename{DE}
 must be a literal member of the context instead of a derived formula, any
 computation equation may be added to the context, and all other rules are
 unchanged.  It argues that this calculus does not derive
 $\proves\It{v}{0}{u}{u}=0$, which $\Nat$ derives by equation induction:
 from equation-only contexts an introduction cannot end a derivation of an
 equation and neither elimination finds a compound major, so only
 hypotheses, computation equations and transport remain, and the free term
 algebra modulo the computation equations separates the two sides.  The
 article does not state this as a theorem; a complete proof would have to
 fix the induction on derivations and the equational countermodel exactly.
 It bounds which checker shortcuts preserve the open consequences of $\Nat$,
 not the size of every certificate design.
\end{enumerate}
\end{question}

\begin{question}[Names and scopes]
\label{wit:q:ns}
\pkg{names\_and\_scopes} asks the following.
\begin{enumerate}[label=\textup{(\roman*)}]
\item Can \rulename{Names} be encoded through the same mechanism as the
 computations, so that replacing full \rulename{Alias} by
 \rulename{EqAlias} and \rulename{Names} makes a complete checker smaller?
 A restriction of the rules alone does not; the readable checker of
 $\mathsf{A}^{\mathrm n}_{\nu}$ is not a compact program.
\item Can \rulename{EqAlias} be absorbed into the names and \rulename{Open}
 without an equally general operation elsewhere?  The extraction of
 $b=t$ from $\mathcal N_{b,t}$ (\cref{wit:lem:ns-base}(iv)) specializes a
 clause at $1$ through \rulename{EqAlias}, so a removal needs another
 bootstrap.
\item Can the evaluations of \cref{wit:thm:ns-obstruction,wit:thm:ns-templates},
 the reduction of \cref{wit:thm:ns-complete} and its compiler be formalized
 against the Lean development?  The persistent families give the induction
 structure; no part is machine-checked.
\item What is the smallest exact public grammar whose encoder has a total,
 explicit inverse on accepted goals?  A broader sorted grammar admits
 padding aliases whose derivability need not agree with that of the
 canonical spelling.
\end{enumerate}
The negative theorems are separations of the stated consequence relations.
Whether a calculus with equation-only \rulename{Back} and \rulename{Alias},
the introduction at the newest parameter and some other finite addition, or
a different interpretation, presents $\iPRA$ is not decided by them.
\end{question}

\begin{question}[Reconstructed introduction]
\label{wit:q:recon}
\pkg{reconstructed\_introduction} leaves the following open.
\begin{enumerate}[label=\textup{(\roman*)}]
\item Its target, a presentation whose complete checker lies below the
 compact envelope 238 / 255 / 455 of the corpus, is not attained: the
 readable checker of $\mathsf{W}_{\mathrm r}$ (\cref{wit:tab:sizes}) pays for
 the reverse replacement, the freshness test and the comparison of the seed.
 Can abstraction and the opening of binders share one traversal in a compact
 complete checker of $\mathsf{W}_{\mathrm r}$?
\item (Raised by the article.)  Can the matrix of the introduction and that
 of the elimination of a closure both be reconstructed, that is, is
 $\mathsf{W}_{\mathrm r}$ with \rulename{QE} replaced by \rulename{AQE} of
 \cref{kern:def:A7} equivalent to $\Nat$?  The proofs of
 \cref{wit:lem:recon-rules} and \cref{kern:lem:aqerank} suggest the order
 transport, elimination, introduction at each rank, but neither
 \pkg{reconstructed\_introduction} nor \pkg{abstraction\_and\_refinements}
 states or proves the combination.
\item \Cref{wit:thm:recon} is a written proof with an independent written
 review; it is not formalized in Lean or Rocq, and the correctness of the
 reference checker is supported by finite tests only.  No minimality of
 $\mathsf{W}_{\mathrm r}$ is claimed.
\end{enumerate}
\end{question}

\begin{question}[Less normalization]
The proofs of \cref{wit:sec:recursor,wit:sec:copy} use the literal
identities $P(\Suc t)\equiv t$ and the two instances of the selector.  Is
there a canonical syntax that guarantees these reductions without full
normalization, with its own coherence law in place of \cref{ctr:lem:nf}?
\end{question}

Further remarks of the packages on what is not shown:
\begin{itemize}
\item No package proves that its presentation is minimal, and no checker of
 a calculus of the chapter improves the compact programs of the corpus; the
 compact checkers of \pkg{export\_specialization} and
 \pkg{vacuous\_specialization} check another calculus.  Whether the restriction of the witness
 improves the compact records under a matched contract is open, and a
 comparison that fixes one validation contract for all competing calculi is
 proposed.  The question of \pkg{deterministic\_introduction} whether the
 selector by pattern can be recognized more cheaply is superseded by
 \cref{wit:thm:guarded-uniform}, which removes the selector.
\item \pkg{local\_introduction} records that its target, a presentation
 whose complete checker lies below the current compact programs, is not
 attained: the checker of $\mathsf{A}_{\nu}$ is smaller than its matched
 control only, and the diagonal test it needs is absent from the
 witness-selected Alias/Back checkers.  Whether the removal of the witness
 can reduce the smallest complete checker, with the diagonal or the mirror
 test and the alias machinery made cheaper while exact freshness and a
 stated input contract are kept, is open; no measurement of a mirror
 variant is made.  The package also asks for a formulation of a genuine
 lower bound on checker size, which has to fix the implementation
 language, the input contract, the accounting of dependencies, the
 convention for success and the resource model; the measured envelope and
 the obstruction theorems of the article do not provide one.
\item \pkg{equational\_transport} records that its target, a complete
 checker below the current envelope, is not attained: both restricted
 checkers are larger than their control (\cref{wit:tab:sizes}), since
 \rulename{Alias} still needs substitution in formulas and \rulename{XBack}
 must destructure its minor premise.  Its remeasured compact
 representatives, under their own contracts, are 239 / 261 / 549,
 245 / 256 / 580, 244 / 269 / 547, 252 / 261 / 593, 274 / 303 / 465,
 267 / 298 / 461 with a shifted depth encoding, and 330 / 372 / 599 with
 the native goal validated; the minima 239, 256 and 461 belong to different
 programs.  A compact variant with 255 nodes, obtained by folding a
 conditional into a pattern guard, is unsound: it accepts the false guarded
 goal $W(0)=W(\Suc 0)$, $W(t)=\It{0}{t}{u}{0}$, with an empty packet of
 the wrong arity, and the package includes it only as a rejected
 experiment.  The package argues that simple changes of the depth offset of
 the compact traversal lose the no-capture property of the opening, and that
 totalizing decoders which admit malformed syntax lose the formula round
 trips unless a normalization is added to the checking contract; these are
 limits of the rewrites tried, not lower bounds.  Whether the equational
 interface of \rulename{XBack} can be built into a smaller complete checker
 is open.
\item \pkg{names\_and\_scopes} records that its target, a single complete
 checker below the whole current envelope, is not attained: the checker of
 $\mathsf{A}^{\mathrm n}_{\nu}$ is a readable reference module of
 928 / 1087 / 3767, and its separate refinements of the exact-binding
 native checker do not reach the compact incumbents under their own
 contracts: 238 tokens and 455 bytes for two different seeded-sequent
 programs and 255 nodes for a compact theorem-input stream.  Whether one program can improve all three
 coordinates under a fixed input, output and validation contract, with the
 size of certificates and the time of checking reported, is open.  Its
 compiler is untrusted and checked by the replay of its outputs only.
\item \pkg{deterministic\_introduction} supplies a complete elaborator from
 certificates of $\mathsf{G}_{\mathrm d}$ to certificates of $\mathsf{G}$
 and illustrates the converse by one translated pair; \pkg{uniform\_raw}
 supplies no translator; the translator of \pkg{conditional\_copy} is not
 verified, and the compiler of \pkg{local\_introduction} is checked only by
 the replay of its outputs.  The correctness of the Python programs is supported by finite
 tests only.
\item In \cref{wit:thm:recursor-bi,wit:prop:liberal} the equivalence of the
 diagonal calculus with $\Nat$ and the bridge are used over the vocabulary
 extended by the fixed descriptions $R$.  The packages argue that the
 proofs are parametric in additional terms; the Lean development covers the
 iterator terms.
\end{itemize}

